% Options for packages loaded elsewhere
\PassOptionsToPackage{unicode}{hyperref}
\PassOptionsToPackage{hyphens}{url}
\PassOptionsToPackage{dvipsnames,svgnames,x11names}{xcolor}
%
\documentclass[
  10pt,
]{article}
\usepackage{amsmath,amssymb}
\usepackage{iftex}
\ifPDFTeX
  \usepackage[T1]{fontenc}
  \usepackage[utf8]{inputenc}
  \usepackage{textcomp} % provide euro and other symbols
\else % if luatex or xetex
  \usepackage{unicode-math} % this also loads fontspec
  \defaultfontfeatures{Scale=MatchLowercase}
  \defaultfontfeatures[\rmfamily]{Ligatures=TeX,Scale=1}
\fi
\usepackage{lmodern}
\ifPDFTeX\else
  % xetex/luatex font selection
\fi
% Use upquote if available, for straight quotes in verbatim environments
\IfFileExists{upquote.sty}{\usepackage{upquote}}{}
\IfFileExists{microtype.sty}{% use microtype if available
  \usepackage[]{microtype}
  \UseMicrotypeSet[protrusion]{basicmath} % disable protrusion for tt fonts
}{}
\makeatletter
\@ifundefined{KOMAClassName}{% if non-KOMA class
  \IfFileExists{parskip.sty}{%
    \usepackage{parskip}
  }{% else
    \setlength{\parindent}{0pt}
    \setlength{\parskip}{6pt plus 2pt minus 1pt}}
}{% if KOMA class
  \KOMAoptions{parskip=half}}
\makeatother
\usepackage{xcolor}
\usepackage{longtable,booktabs,array}
\usepackage{calc} % for calculating minipage widths
% Correct order of tables after \paragraph or \subparagraph
\usepackage{etoolbox}
\makeatletter
\patchcmd\longtable{\par}{\if@noskipsec\mbox{}\fi\par}{}{}
\makeatother
% Allow footnotes in longtable head/foot
\IfFileExists{footnotehyper.sty}{\usepackage{footnotehyper}}{\usepackage{footnote}}
\makesavenoteenv{longtable}
\setlength{\emergencystretch}{3em} % prevent overfull lines
\providecommand{\tightlist}{%
  \setlength{\itemsep}{0pt}\setlength{\parskip}{0pt}}
\setcounter{secnumdepth}{5}

\usepackage{amsmath,amssymb}
\usepackage{graphicx}
\usepackage{longtable,booktabs,array}
\usepackage{xcolor}
\usepackage[margin=2.0cm]{geometry}
\usepackage{microtype}
\setlength{\emergencystretch}{3em}
\setcounter{MaxMatrixCols}{20}
\pdfinfoomitdate=1
\pdftrailerid{}
\pdfsuppressptexinfo=-1
\renewcommand{\arraystretch}{1.12}
\usepackage{fvextra}
\RecustomVerbatimEnvironment{verbatim}{Verbatim}{breaklines,breakanywhere,fontsize=\small}
\fvset{breaklines,breakanywhere}
\allowdisplaybreaks
\sloppy
\ifLuaTeX
  \usepackage{selnolig}  % disable illegal ligatures
\fi
\IfFileExists{bookmark.sty}{\usepackage{bookmark}}{\usepackage{hyperref}}
\IfFileExists{xurl.sty}{\usepackage{xurl}}{} % add URL line breaks if available
\urlstyle{same}
\hypersetup{
  pdftitle={UM-3 field equations},
  pdfauthor={Underlying theory: Patrick L. Nash; Derived and verified by Claude (Anthropic); Field equations of UM-3, 2026-10-24; Repository once-ere/Dirac\_claude, folder universe\_model\_3},
  colorlinks=true,
  linkcolor={blue},
  filecolor={Maroon},
  citecolor={Blue},
  urlcolor={blue},
  pdfcreator={LaTeX via pandoc}}

\title{UM-3 field equations}
\usepackage{etoolbox}
\makeatletter
\providecommand{\subtitle}[1]{% add subtitle to \maketitle
  \apptocmd{\@title}{\par {\large #1 \par}}{}{}
}
\makeatother
\subtitle{The complete set of coupled field equations of
Einstein-Lovelock gravity, dirac16complex, the Higgs doublet with the ad
hoc vacuum, the gauge fields of SU(3), SU(2) and U(1), and three
generations of quarks and leptons with three right-handed neutrinos}
\author{Underlying theory: Patrick L. Nash \and Derived and verified by
Claude (Anthropic) \and Field equations of UM-3,
2026-10-24 \and Repository once-ere/Dirac\_claude, folder
universe\_model\_3}
\date{}

\begin{document}
\maketitle

{
\hypersetup{linkcolor=}
\setcounter{tocdepth}{3}
\tableofcontents
}
\hypertarget{how-to-read-this-document}{%
\section{How to read this document}\label{how-to-read-this-document}}

\hypertarget{what-this-document-contains}{%
\subsection{What this document
contains}\label{what-this-document-contains}}

This document states the complete set of coupled field equations of the
universe model UM-3, derives every one of them from one action, and
explains every term. It is self-contained: every equation it uses is
written in it, and every symbol of every equation is defined right after
the equation. The sections are:

\begin{itemize}
\tightlist
\item
  \textbf{Where the fields live} (this section): the author's metric has
  a fold, the seven-dimensional surface \(z = \pi/2\), that carries
  energy; the Standard Model lives on the fold, dirac16complex and
  gravity in the bulk.
\item
  \textbf{The action} (IN-1 to IN-16): every term of the action, written
  out and explained.
\item
  \textbf{Gravity} (FE-1 to FE-11): the Einstein-Lovelock equations in
  the bulk, the junction conditions at the fold, the energy-momentum
  tensor of every field and the conservation laws.
\item
  \textbf{The matter fields} (FE-12 to FE-26): the field equations of
  dirac16complex, of the Higgs doublet, of the gauge fields and of the
  quarks and leptons, and the equations of the ad hoc Higgs vacuum.
\item
  \textbf{The background} (FE-27 to FE-35): the author's metric, its
  exact curvature, the exact background equations in the bulk and at the
  fold, the author's vacuum, the static late-time state, the frozen
  extra times and the power series of the metric.
\item
  \textbf{dirac16complex: modes and Kohn-Sham equations} (FE-36 to
  FE-43): the equation of the hidden direction, the bulk continuum, the
  states bound to the fold, the Kohn-Sham equations with the mean fields
  of every other field, the Krein norm.
\item
  \textbf{The Standard Model on the fold} (FE-44 to FE-47): the
  four-dimensional couplings, the masses with the ad hoc vacuum, the
  mixing matrices and the four-dimensional couplings of the bound states
  of dirac16complex.
\item
  \textbf{The three neutrino flavours} (FE-48 to FE-53): the flavour
  evolution in vacuum, in matter, in dense neutrino gases (supernovae)
  and in the early Universe.
\item
  \textbf{The cosmological perturbations} (FE-54 to FE-58): the metric
  as a power series and the Einstein-Boltzmann equations of the CMB, the
  CNB and the BAO.
\item
  \textbf{The thermal history} (FE-59 to FE-62): the expansion of the
  fold, the degrees of freedom and the Boltzmann equations of the bound
  states of dirac16complex and of the right-handed neutrinos.
\item
  \textbf{The coupled system}: which field appears in which equation,
  and how the system is solved.
\item
  \textbf{What was verified}: the checks of the computations behind
  every derived statement.
\end{itemize}

\hypertarget{the-status-labels}{%
\subsection{The status labels}\label{the-status-labels}}

Every equation and every result carries one of the labels

\begin{itemize}
\tightlist
\item
  \textbf{AUTHOR}: from the author (his notebooks \texttt{4+4.nb} and
  \texttt{Dirac\_Einstein-Lovelock\_4+4\_Nash.nb}, and his order of
  2026-10-24: the coordinates, the metric, the Lovelock weights, the
  Dirac matrices, the chiral projectors, the list of fields, the Higgs
  potential and the ad hoc vacuum);
\item
  \textbf{RECORD}: from the record of this repository (the Lagrangian of
  dirac16complex with its canonical spin connection and its contact
  self-interaction, the Kohn-Sham treatment);
\item
  \textbf{STANDARD}: textbook physics;
\item
  \textbf{ASSUMED}: a choice made in UM-3 without derivation;
\item
  \textbf{INVENTED}: a new interaction of UM-3 (the portals of
  dirac16complex);
\item
  \textbf{DERIVED}: computed and verified in UM-3 (the verification is
  stated);
\item
  \textbf{ESTIMATE}: an order-of-magnitude evaluation;
\item
  \textbf{OPEN}: not settled.
\end{itemize}

\hypertarget{every-equation-is-followed-by-its-symbols}{%
\subsection{Every equation is followed by its
symbols}\label{every-equation-is-followed-by-its-symbols}}

Every display equation is followed by the list ``Symbols in (FE-n)'' (or
``Symbols in this equation'' for a display without a tag) of EVERY
symbol that appears in it: every quantity, field, coupling, constant and
operator, every index letter with the values it takes, every mark
(prime, dagger, transpose, star), the imaginary unit, Euler's number and
the differential. Each item gives the meaning, the unit (natural units
\(\hbar = c = 1\)), the value where one exists, and the status. A field
equation carries the tag FE-1 to FE-62 and a term of the action the tag
IN-1 to IN-16; the numbering is that of the inventory of UM-3 and is the
same in every UM-3 document. Displays without a tag are steps of a
derivation.

\hypertarget{coordinates-signature-and-units}{%
\subsection{Coordinates, signature and
units}\label{coordinates-signature-and-units}}

The author's coordinates are \(X = \{x_0, x_1, \dots, x_7\}\), always
with the hidden direction first: \(x_0\) is the hidden space-like
coordinate, on which the metric depends periodically with the period
\(\pi/(6H)\) through the angle \(z = 6Hx_0\); \(x_1, x_2, x_3\) are
ordinary space, \(x_4\) is the time, and \(x_5, x_6, x_7\) are three
extra time-like coordinates. The signature is \((+,+,+,+,-,-,-,-)\) in
this order; the four observable coordinates \(x_1, \dots, x_4\) have the
signature \((+,+,+,-)\). \(H\) is the author's constant (an inverse
length); this document keeps it as a free parameter. Units are natural,
\(\hbar = c = 1\); in eight dimensions the field dirac16complex has the
dimension GeV\(^{7/2}\) and the gravitational coupling
\(\kappa = 8\pi G_8\) the dimension GeV\(^{-6}\). A prime on a function
of the time \(x_4\) alone is its derivative \(\mathrm{d}/\mathrm{d}x_4\)
(in the cosmological perturbation equations, and only there, a dot is
the derivative with respect to the conformal time). The proper
coordinate of the hidden direction is \(\sigma\) (defined in FE-27):
\(\sigma = 0\) is the fold \(z = \pi/2\).

\hypertarget{the-fields-of-um-3-and-where-they-live}{%
\subsection{The fields of UM-3 and where they
live}\label{the-fields-of-um-3-and-where-they-live}}

The author orders a model universe with the following coupled fields
(his order of 2026-10-24, items 1 to 10); UM-3 adds the right-handed
quark singlets, the Yukawa couplings and three right-handed neutrinos
(accepted by the author on 2026-10-24 at 17:32), the stabiliser of the
bulk, the fold fluid and the induced gravity of the fold (ASSUMED,
explained below).

\begin{longtable}[]{@{}
  >{\raggedright\arraybackslash}p{(\columnwidth - 10\tabcolsep) * \real{0.1619}}
  >{\raggedright\arraybackslash}p{(\columnwidth - 10\tabcolsep) * \real{0.0952}}
  >{\raggedright\arraybackslash}p{(\columnwidth - 10\tabcolsep) * \real{0.2095}}
  >{\raggedright\arraybackslash}p{(\columnwidth - 10\tabcolsep) * \real{0.0762}}
  >{\raggedright\arraybackslash}p{(\columnwidth - 10\tabcolsep) * \real{0.2286}}
  >{\raggedright\arraybackslash}p{(\columnwidth - 10\tabcolsep) * \real{0.2286}}@{}}
\toprule\noalign{}
\begin{minipage}[b]{\linewidth}\raggedright
field
\end{minipage} & \begin{minipage}[b]{\linewidth}\raggedright
symbol
\end{minipage} & \begin{minipage}[b]{\linewidth}\raggedright
what it is
\end{minipage} & \begin{minipage}[b]{\linewidth}\raggedright
lives in
\end{minipage} & \begin{minipage}[b]{\linewidth}\raggedright
it couples to
\end{minipage} & \begin{minipage}[b]{\linewidth}\raggedright
field equations
\end{minipage} \\
\midrule\noalign{}
\endhead
\bottomrule\noalign{}
\endlastfoot
Einstein-Lovelock gravity & \(g_{MN}\) & the metric of eight dimensions
& everywhere & every field & FE-1 to FE-11, FE-27 to FE-35, FE-54 to
FE-56 \\
dirac16complex & \(\Psi\) & the author's 16-component complex spinor
field; its left- and right-handed halves \(\Psi_{\mathrm{L}}\),
\(\Psi_{\mathrm{R}}\) & the bulk & gravity (spin connection), itself,
every Standard-Model field at the fold (IN-16) & FE-6, FE-12 to FE-14,
FE-36 to FE-43, FE-47 \\
Higgs doublet & \(\Phi\) & the scalar doublet with
\(V = \mu^2\Phi^\dagger\Phi + \lambda(\Phi^\dagger\Phi)^2\), expanded
about the ad hoc vacuum \(\Phi = (\mathcal{A}, v + h)/\sqrt2\) & the
fold & gravity, \(W\), \(B\), the fermions (Yukawa),
\(\Psi_{\mathrm{L}}\) only & FE-7, FE-15, FE-26, FE-45 \\
three left-handed lepton doublets & \(\ell_i\) &
\((\nu_{\mathrm{L}i}, e_{\mathrm{L}i})\), \(i = 1, 2, 3\) & the fold &
gravity, \(W\), \(B\), \(\Phi\), \(\Psi_{\mathrm{L}}\) and
\(\Psi_{\mathrm{R}}\) & FE-9, FE-22, FE-48 to FE-53 \\
right-handed electron, muon, tau & \(e_{\mathrm{R}i}\) & the
charged-lepton singlets & the fold & gravity, \(B\), \(\Phi\),
\(\Psi_{\mathrm{L}}\), \(\Psi_{\mathrm{R}}\) & FE-9, FE-23 \\
three right-handed neutrinos & \(\nu_{\mathrm{R}i}\) & gauge singlets;
Dirac neutrino masses & the fold & gravity, \(\Phi\),
\(\Psi_{\mathrm{L}}\), \(\Psi_{\mathrm{R}}\) & FE-9, FE-24, FE-62 \\
gauge fields of \(SU(2)_L\) and \(U(1)_Y\) & \(W^1, W^2, W^3\), \(B\) &
the electroweak gauge fields & the fold & gravity, \(\Phi\), the
fermions, \(\Psi_{\mathrm{L}}\), \(\Psi_{\mathrm{R}}\) & FE-8, FE-17,
FE-18 \\
three quark generations & \(Q_i\), \(u_{\mathrm{R}i}\),
\(d_{\mathrm{R}i}\) & the quark doublets and singlets & the fold &
gravity, gluons, \(W\), \(B\), \(\Phi\), \(\Psi_{\mathrm{L}}\),
\(\Psi_{\mathrm{R}}\) & FE-9, FE-19 to FE-21 \\
gluons & \(G^A\) & the gauge field of \(SU(3)_C\) & the fold & gravity,
quarks, \(\Psi_{\mathrm{L}}\), \(\Psi_{\mathrm{R}}\) & FE-8, FE-16 \\
stabiliser & \(\mathcal{T}_{\mathrm{stab}}\) & a medium of the bulk that
keeps the late-time bulk static & the bulk & gravity & FE-3, FE-32 \\
fold fluid & \(\mathcal{S}_{\mathrm{ff}}\) & the thin source that the
kink of the metric requires at the fold & the fold & gravity & FE-4,
FE-5, FE-31, FE-32 \\
induced gravity of the fold & \(\kappa_4\) & an Einstein-Hilbert term of
the fold & the fold & gravity & FE-5, FE-59 \\
\end{longtable}

\textbf{Why the Standard Model lives on the fold (DERIVED).} In the
proper coordinate \(\sigma\) of the hidden direction (FE-27), one period
of \(x_0\) is the whole line \(-\infty < \sigma < \infty\) and the
author's metric is
\(\mathrm{d}\sigma^2 + \mathrm{e}^{-2H\lvert\sigma\rvert}\big(\mathrm{e}^{2a_4}\,\mathrm{d}\vec x^2 - \mathrm{e}^{-2a_4}b_4^2\,\mathrm{d}\vec x_5^2\big) - \mathrm{d}x_4^2\):
ordinary space is warped by \(\mathrm{e}^{-H\lvert\sigma\rvert}\), the
time \(x_4\) is not. A field confined to the fold \(\sigma = 0\) has the
light cone of the fold; a field that spreads into the bulk moves along
ordinary space with a speed set by the \(\sigma\)-average of
\(\mathrm{e}^{2H\lvert\sigma\rvert}\) over its profile. With the
author's metric a bulk photon moves \(\sqrt2\) times faster than light
on the fold, a bulk scalar \(\sqrt{3/2}\) times, a bulk fermion with a
profile of width \(1/M\) \(2M/(2M - H)\) times (DERIVED,
\texttt{um3\_bulk.py}, checks B1 to B3). Species with different profiles
would have different limiting speeds (excluded at the level \(10^{-15}\)
to \(10^{-20}\) for photons, electrons and nucleons), and three neutrino
flavours with three profiles would oscillate with energy-proportional
phases that observation excludes. The author's request for correct
flavour conversion therefore fixes the place of the Standard Model: all
of it lives on the fold, where it has one light cone, and its flavour
physics is that of the Standard Model with three Dirac neutrinos.
dirac16complex is the author's field of eight dimensions and lives in
the bulk; it couples to the Standard Model at the fold.

\textbf{Why the fold carries a source (DERIVED).} The warp
\(\mathrm{e}^{-2H\lvert\sigma\rvert}\) has a kink at the fold. The
Einstein-Lovelock equations, which are linear in the second derivatives
of the metric along \(\sigma\), then contain a delta function at
\(\sigma = 0\): the metric is a solution only if the fold carries a thin
source \(\mathcal{S}^a{}_b\) (FE-4). At the author's vacuum
\(\kappa\mathcal{S} = \mathrm{diag}(-10H, -6H, -10H)\) for (ordinary
space, time, extra times) (DERIVED, \texttt{um3\_fold.py}, check J5).
This source is the fold fluid together with the Standard Model and the
portals.

\textbf{Why the fold has its own Einstein-Hilbert term (ASSUMED).} The
transverse traceless gravitons of the bulk, with the author's Lovelock
weights, propagate with \(c_{\mathrm{g}}^2 = 15/2\) times the light cone
of the fold at the static late-time state, and they have no time-kinetic
term at all at the author's vacuum (DERIVED, \texttt{um3\_bulk.py}, B5).
Gravitational waves move with the speed of light to \(10^{-15}\)
(GW170817). UM-3 therefore gives the fold an induced Einstein-Hilbert
term (IN-9) that dominates four-dimensional gravity; the bulk
contributes a fraction \(r \lesssim 10^{-15}\) (FE-59).

\hypertarget{conventions}{%
\subsection{Conventions}\label{conventions}}

\begin{itemize}
\tightlist
\item
  Indices. \(M, N, P, Q\) (with numbered copies \(M_1, M_2, \dots\)):
  coordinate indices of eight dimensions; \(a, b, c, d\): frame indices
  of eight dimensions (the frame direction \((0)\) is the \(\sigma\)
  direction); \(\mu, \nu, \rho\): coordinate indices of the four
  observable directions \(x_1, \dots, x_4\); \(A, B, C = 1, \dots, 8\):
  colour (adjoint) indices; \(I, J, K = 1, 2, 3\): weak-isospin
  (adjoint) indices; \(i, j = 1, 2, 3\): generations (and the neutrino
  mass states); \(\alpha, \beta = e, \mu, \tau\): lepton flavours;
  \(p = 1, 2, 3\): the Lovelock order; \(X = \mathrm{L}, \mathrm{R}\):
  the two halves of dirac16complex; \(f \in \{Q, u, d, \ell, e, \nu\}\):
  the six kinds of Standard-Model multiplets. A repeated index is summed
  over its range.
\item
  Two kinds of Dirac matrices. The author's real \(16 \times 16\)
  matrices \(T_a\) (IN-5) act on dirac16complex; the four-dimensional
  complex \(4 \times 4\) matrices \(\hat\gamma^\mu\) (IN-14) act on the
  Standard-Model fermions. They are never multiplied with each other;
  the couplings between dirac16complex and the Standard Model are
  products of Lorentz tensors.
\item
  Normalisation of the fold. No field depends on the extra times
  \(x_5, x_6, x_7\) (ASSUMED: a dependence on them gives tachyonic
  four-dimensional modes, FE-38). Integrated over them, the extra times
  give their physical volume on the fold, \(V_*\); the Standard-Model
  fields are normalised as four-dimensional fields, the couplings are
  their measured four-dimensional values, and the fold action carries
  the factor \(1/V_*\) (IN-8).
\item
  A word or a particle name in a subscript is upright
  (\(m_{\mathrm{W}}\), \(\nu_{\mathrm{R}}\), \(\Psi_{\mathrm{L}}\)); an
  italic letter in a subscript is an index. \(\mathrm{e}\) is Euler's
  number, \(\mathrm{i}\) the imaginary unit, \(\mathrm{d}\) the
  differential; \(\delta\) is the Kronecker delta and
  \(\delta_{\mathrm{D}}\) the Dirac delta function.
\end{itemize}

\hypertarget{the-action}{%
\section{The action}\label{the-action}}

All field equations of UM-3 follow from one action by varying it with
respect to every field. This section writes out every term of the action
(IN-1 to IN-16) and says in plain words what it couples. The
Standard-Model part has exactly the gauge group, the fermion content
(three generations, each with a right-handed neutrino), the Higgs
potential and the gauge couplings of the Standard Model; it lives on the
fold. dirac16complex is a singlet of the Standard-Model gauge group; its
left-handed and right-handed halves couple to every Standard-Model field
at the fold through gauge-invariant portals with their own coupling
constants, and only its left-handed half couples to the Higgs doublet.

\hypertarget{the-sum-of-all-terms}{%
\subsection{The sum of all terms}\label{the-sum-of-all-terms}}

\[
\begin{aligned}
S_{\mathrm{tot}} &= S_{\mathrm{EL}} + S_{\mathrm{stab}} + S_\Psi + S_{\mathrm{fold}},\\
S_{\mathrm{fold}} &= \int\mathrm{d}^8X\,\sqrt{\lvert g\rvert}\;\delta_{\mathrm{D}}(\sigma)\,\frac{1}{V_*}\,
\big(\mathcal{L}_{\mathrm{IG}} + \mathcal{L}_{\mathrm{ff}} + \mathcal{L}_{\mathrm{H}} + \mathcal{L}_{\mathrm{gauge}}
+ \mathcal{L}_{\mathrm{f}} + \mathcal{L}_{\mathrm{Y}} + \mathcal{L}_{\mathrm{P}}\big)
\end{aligned} \tag{IN-1}
\] \textbf{Symbols in (IN-1).}

\begin{itemize}
\tightlist
\item
  \(S_{\mathrm{tot}}\): The action of UM-3 (the integral of the
  Lagrangian density over spacetime); its variation with respect to
  every field gives the field equations. Unit: dimensionless
  (\(\hbar = 1\)). Status: ASSUMED.
\item
  \(S_{\mathrm{EL}}\): The Einstein-Lovelock action of gravity in eight
  dimensions. Unit: dimensionless. Status: AUTHOR.
\item
  \(S_{\mathrm{stab}}\): The action of the stabiliser of the bulk (a
  medium specified by its energy-momentum tensor
  \(\mathcal{T}_{\mathrm{stab}}\)). Unit: dimensionless. Status:
  ASSUMED.
\item
  \(S_\Psi\): The action of dirac16complex in the bulk (kinetic term
  with the canonical spin connection, mass term, contact
  self-interaction). Unit: dimensionless. Status: RECORD.
\item
  \(S_{\mathrm{fold}}\): The action of the fold: the Standard Model, the
  portals of dirac16complex, the fold fluid and the induced gravity of
  the fold, all at \(\sigma = 0\). Unit: dimensionless. Status: ASSUMED.
\item
  \(\mathrm{d}\): The differential: \(\mathrm{d}X\) is an infinitesimal
  change of \(X\), \(\mathrm{d}X/\mathrm{d}Y\) the derivative of \(X\)
  with respect to \(Y\), \(\mathrm{d}^nX\) the volume element of \(n\)
  variables. Status: STANDARD.
\item
  \(X = \{x_0, \dots, x_7\}\): The author's coordinates of the
  eight-dimensional spacetime, \(X = \{x_0, x_1, \dots, x_7\}\) with the
  hidden direction first (Dirac\_Einstein-Lovelock\_4+4\_Nash.nb,
  section DEFCHART); \(X^M\) is the coordinate with the index \(M\) and
  \(\mathrm{d}^8X = \mathrm{d}x_0\,\mathrm{d}x_1\cdots\mathrm{d}x_7\)
  the coordinate volume element. Unit: length (each coordinate). Status:
  AUTHOR.
\item
  \(g_{MN}\): The metric tensor of the eight-dimensional spacetime
  (signature \(+,+,+,+,-,-,-,-\) in the order \(x_0, \dots, x_7\)):
  \(g_{MN}\,\mathrm{d}X^M\mathrm{d}X^N\) is the squared interval,
  \(g^{MN}\) its inverse, \(\lvert g\rvert\) the absolute value of its
  determinant (\(\sqrt{\lvert g\rvert}\,\mathrm{d}^8X\) is the proper
  volume element). Unit: dimensionless. Status: AUTHOR.
\item
  \(\delta_{\mathrm{D}}\): The Dirac delta function:
  \(\int\mathrm{d}\sigma\,\delta_{\mathrm{D}}(\sigma)F(\sigma) = F(0)\).
  Unit: 1/length (for \(\delta_{\mathrm{D}}(\sigma)\)). Status:
  STANDARD.
\item
  \(\sigma\): The proper coordinate of the hidden direction:
  \(\mathrm{d}\sigma = \lvert\cot z\rvert\,\mathrm{d}x_0\),
  \(\sigma = \mathrm{sgn}(\pi/2 - z)\ln(\sin z)/(6H)\) on one period
  \(0 < z < \pi\); \(\sigma = 0\) is the fold \(z = \pi/2\),
  \(\sigma \to \mp\infty\) the tips \(z \to 0, \pi\); one period of
  \(x_0\) is the whole line \(-\infty < \sigma < \infty\). Unit: length.
  Status: DERIVED.
\item
  \(V_*\): The physical volume of the extra times \(x_5, x_6, x_7\) on
  the fold in the late universe,
  \(V_* = V_{\mathrm{t}}(\mathrm{e}^{-a_4}b_4)^3\) at a time when
  \(\mathrm{e}^{-a_4}b_4\) is constant; it converts the
  eight-dimensional normalisation into the four-dimensional one. Unit:
  length\(^3\) (GeV\(^{-3}\)). Status: ASSUMED.
\item
  \(\mathcal{L}_{\mathrm{IG}}\): The Lagrangian of the induced gravity
  of the fold. Unit: GeV\(^4\). Status: ASSUMED.
\item
  \(\mathcal{L}_{\mathrm{ff}}\): The Lagrangian of the fold fluid (only
  its energy-momentum tensor \(\mathcal{T}_{\mathrm{ff}}\) enters the
  equations). Unit: GeV\(^4\). Status: ASSUMED.
\item
  \(\mathcal{L}_{\mathrm{H}}\): The Lagrangian of the Higgs doublet on
  the fold. Unit: GeV\(^4\). Status: AUTHOR.
\item
  \(\mathcal{L}_{\mathrm{gauge}}\): The Lagrangian of the gauge fields
  of \(SU(3)_C\times SU(2)_L\times U(1)_Y\) on the fold. Unit:
  GeV\(^4\). Status: STANDARD.
\item
  \(\mathcal{L}_{\mathrm{f}}\): The kinetic Lagrangian of the quarks and
  leptons on the fold. Unit: GeV\(^4\). Status: STANDARD.
\item
  \(\mathcal{L}_{\mathrm{Y}}\): The Yukawa Lagrangian of the quarks and
  leptons with the Higgs doublet. Unit: GeV\(^4\). Status: STANDARD.
\item
  \(\mathcal{L}_{\mathrm{P}}\): The Lagrangian of the portals of
  dirac16complex at the fold. Unit: GeV\(^4\). Status: INVENTED.
\end{itemize}

The action has four parts: Einstein-Lovelock gravity in eight dimensions
(\(S_{\mathrm{EL}}\), IN-2, IN-3), the stabiliser of the bulk
(\(S_{\mathrm{stab}}\), IN-7), the field dirac16complex in the bulk
(\(S_\Psi\), IN-4 to IN-6) and the fold (\(S_{\mathrm{fold}}\), IN-8 to
IN-16). The Dirac delta \(\delta_{\mathrm{D}}(\sigma)\) of the proper
hidden coordinate \(\sigma\) confines the terms of the fold to
\(\sigma = 0\); the factor \(1/V_*\) turns the integral over the extra
times into the normalisation of four-dimensional fields (IN-8). On the
fold live: the induced gravity of the fold
(\(\mathcal{L}_{\mathrm{IG}}\)), the fold fluid
(\(\mathcal{L}_{\mathrm{ff}}\)), the Higgs doublet
(\(\mathcal{L}_{\mathrm{H}}\)), the gauge fields
(\(\mathcal{L}_{\mathrm{gauge}}\)), the quarks and leptons
(\(\mathcal{L}_{\mathrm{f}}\)), their Yukawa couplings
(\(\mathcal{L}_{\mathrm{Y}}\)) and the portals of dirac16complex
(\(\mathcal{L}_{\mathrm{P}}\)). Varying \(S_{\mathrm{tot}}\) with
respect to the metric gives the gravitational equations (FE-1, FE-4);
with respect to \(\bar\Psi\) the equation of dirac16complex (FE-12,
FE-13); with respect to \(\Phi^\dagger\) the Higgs equation (FE-15);
with respect to the gauge fields the Yang-Mills equations (FE-16 to
FE-18); with respect to every fermion its Dirac equation (FE-19 to
FE-24). Status: ASSUMED (the choice and the place of the terms; the
content of each term is stated with its own label).

\hypertarget{gravity}{%
\subsection{Gravity}\label{gravity}}

\[
\begin{aligned}
S_{\mathrm{EL}} &= \frac{1}{2\kappa}\int\mathrm{d}^8X\,\sqrt{\lvert g\rvert}\,\Big(\sum_{p=1}^{3}\alpha_{(p)}\,\mathcal{L}_{(p)}
 - 2\Lambda\Big),\qquad \kappa = 8\pi G_8,\\
\alpha_{(p)} &= w_{(p)}\,H^{2-2p},\qquad w_{(1)} = 1,\quad w_{(2)} = \frac{1}{32},\quad w_{(3)} = \frac{1}{1152},
 \qquad \Lambda = -10H^2
\end{aligned} \tag{IN-2}
\] \textbf{Symbols in (IN-2).}

\begin{itemize}
\tightlist
\item
  \(S_{\mathrm{EL}}\): The Einstein-Lovelock action of gravity in eight
  dimensions. Unit: dimensionless. Status: AUTHOR.
\item
  \(\kappa\): The eight-dimensional gravitational coupling,
  \(\kappa = 8\pi G_8\). Unit: GeV\(^{-6}\). Status: AUTHOR.
\item
  \(\mathrm{d}\): The differential: \(\mathrm{d}X\) is an infinitesimal
  change of \(X\), \(\mathrm{d}X/\mathrm{d}Y\) the derivative of \(X\)
  with respect to \(Y\), \(\mathrm{d}^nX\) the volume element of \(n\)
  variables. Status: STANDARD.
\item
  \(X = \{x_0, \dots, x_7\}\): The author's coordinates of the
  eight-dimensional spacetime, \(X = \{x_0, x_1, \dots, x_7\}\) with the
  hidden direction first (Dirac\_Einstein-Lovelock\_4+4\_Nash.nb,
  section DEFCHART); \(X^M\) is the coordinate with the index \(M\) and
  \(\mathrm{d}^8X = \mathrm{d}x_0\,\mathrm{d}x_1\cdots\mathrm{d}x_7\)
  the coordinate volume element. Unit: length (each coordinate). Status:
  AUTHOR.
\item
  \(g_{MN}\): The metric tensor of the eight-dimensional spacetime
  (signature \(+,+,+,+,-,-,-,-\) in the order \(x_0, \dots, x_7\)):
  \(g_{MN}\,\mathrm{d}X^M\mathrm{d}X^N\) is the squared interval,
  \(g^{MN}\) its inverse, \(\lvert g\rvert\) the absolute value of its
  determinant (\(\sqrt{\lvert g\rvert}\,\mathrm{d}^8X\) is the proper
  volume element). Unit: dimensionless. Status: AUTHOR.
\item
  \(p\): The Lovelock order \(p = 1, 2, 3\) (Einstein, Gauss-Bonnet,
  cubic), as an index and as a number. Status: STANDARD.
\item
  \(\alpha_{(p)}\): The Lovelock coupling of order \(p\),
  \(\alpha_{(p)} = w_{(p)}H^{2-2p}\) (\(\alpha_{(1)} = 1\)). Unit:
  GeV\(^{2-2p}\). Status: AUTHOR.
\item
  \(\mathcal{L}_{(p)}\): The Lovelock density of order \(p\) (a scalar
  built from \(p\) Riemann tensors): \(\mathcal{L}_{(1)} = R\)
  (Einstein-Hilbert), \(\mathcal{L}_{(2)}\) the Gauss-Bonnet density,
  \(\mathcal{L}_{(3)}\) the cubic Lovelock density. Unit:
  1/length\(^{2p}\). Status: STANDARD.
\item
  \(\Lambda\): The eight-dimensional cosmological constant,
  \(\Lambda = -10H^2\) (the author writes the pure number
  \(\Lambda_{\mathrm{N}} = 10\) in
  \(0 = -\Lambda_{\mathrm{N}} + \sum_p H^{-2p}w_{(p)}\mathcal{E}_{(p)}\)).
  Unit: GeV\(^2\). Value: \(-10H^2\). Status: AUTHOR.
\item
  \(\pi\): The number pi, the ratio of the circumference of a circle to
  its diameter. Unit: dimensionless. Value: 3.141592654. Status:
  STANDARD.
\item
  \(G_8\): The eight-dimensional Newton constant
  (\(\kappa = 8\pi G_8\)). Unit: GeV\(^{-6}\). Status: AUTHOR.
\item
  \(w_{(p)}\): The author's Lovelock weights, \(w_{(1)} = 1\),
  \(w_{(2)} = 1/32\), \(w_{(3)} = 1/1152\)
  (Dirac\_Einstein-Lovelock\_4+4\_Nash.nb, In{[}134{]}). Unit:
  dimensionless. Value: \((1, 1/32, 1/1152)\). Status: AUTHOR.
\item
  \(H\): The author's constant (an inverse length) that fixes the period
  \(\pi/(6H)\) of \(x_0\), the Lovelock couplings and the cosmological
  constant \(\Lambda = -10H^2\); in UM-3 a free parameter. Unit: GeV
  (1/length). Status: AUTHOR.
\end{itemize}

Gravity is the Einstein-Lovelock theory in eight dimensions with the
author's weights. \(\mathcal{L}_{(1)}\) is the Einstein-Hilbert density
(the Ricci scalar), \(\mathcal{L}_{(2)}\) the Gauss-Bonnet density and
\(\mathcal{L}_{(3)}\) the cubic Lovelock density; in eight dimensions
the series ends with \(p = 3\) (the density of order 4 is a total
derivative in eight dimensions). The author writes the vacuum equations
as
\(0 = -\Lambda_{\mathrm{N}} + H^{-2}w_{(1)}\mathcal{E}_{(1)} + H^{-4}w_{(2)}\mathcal{E}_{(2)} + H^{-6}w_{(3)}\mathcal{E}_{(3)}\)
with the pure number \(\Lambda_{\mathrm{N}} = 10\)
(\texttt{Dirac\_Einstein-Lovelock\_4+4\_Nash.nb}, In{[}134{]});
multiplied by \(H^2\) this is FE-1 without sources and with
\(\Lambda = -\Lambda_{\mathrm{N}}H^2 = -10H^2\). Gravity couples to
every field through the metric that appears in every other term of the
action (in \(\sqrt{\lvert g\rvert}\), in the contractions of indices and
in the spin connection). Status: AUTHOR.

\[
\begin{aligned}
\mathcal{L}_{(p)} &= \frac{1}{2^{p}}\,\delta^{M_1N_1\cdots M_pN_p}_{P_1Q_1\cdots P_pQ_p}\,
 R^{P_1Q_1}{}_{M_1N_1}\cdots R^{P_pQ_p}{}_{M_pN_p},\\
\mathcal{L}_{(1)} &= R,\qquad \mathcal{L}_{(2)} = R^2 - 4R_{MN}R^{MN} + R_{MNPQ}R^{MNPQ}
\end{aligned} \tag{IN-3}
\] \textbf{Symbols in (IN-3).}

\begin{itemize}
\tightlist
\item
  \(\mathcal{L}_{(p)}\): The Lovelock density of order \(p\) (a scalar
  built from \(p\) Riemann tensors): \(\mathcal{L}_{(1)} = R\)
  (Einstein-Hilbert), \(\mathcal{L}_{(2)}\) the Gauss-Bonnet density,
  \(\mathcal{L}_{(3)}\) the cubic Lovelock density. Unit:
  1/length\(^{2p}\). Status: STANDARD.
\item
  \(p\): The Lovelock order \(p = 1, 2, 3\) (Einstein, Gauss-Bonnet,
  cubic), as an index and as a number. Status: STANDARD.
\item
  \(\delta\): The Kronecker delta: with two indices \(\delta^M_N = 1\)
  if \(M = N\) and 0 otherwise (likewise \(\delta_{X\mathrm{L}} = 1\)
  for \(X = \mathrm{L}\)); with upper and lower index strings the
  generalised Kronecker delta
  \(\delta^{M_1\cdots M_n}_{N_1\cdots N_n}\), the determinant of the
  \(n \times n\) matrix of the Kronecker deltas \(\delta^{M_a}_{N_b}\)
  (\(+1\) or \(-1\) if the upper indices are an even or odd permutation
  of the lower ones, 0 otherwise). Status: STANDARD.
\item
  \(M\): A coordinate index of eight dimensions (also with a numbered
  copy \(M_1, M_2, \dots\)): it runs over \(x_0, x_1, \dots, x_7\); a
  repeated index is summed. Status: STANDARD.
\item
  \(N\): A coordinate index of eight dimensions (also
  \(N_1, N_2, \dots\)), running over \(x_0, \dots, x_7\); a repeated
  index is summed. Status: STANDARD.
\item
  \(P\): A coordinate index of eight dimensions (also
  \(P_1, P_2, \dots\)), running over \(x_0, \dots, x_7\); a repeated
  index is summed. Status: STANDARD.
\item
  \(Q\): A coordinate index of eight dimensions (also
  \(Q_1, Q_2, \dots\)), running over \(x_0, \dots, x_7\); a repeated
  index is summed. Status: STANDARD.
\item
  \(R^{MN}{}_{PQ},\ R_{MN},\ R\): The curvature of the metric
  \(g_{MN}\): with four indices the Riemann tensor (convention
  \([\nabla_P, \nabla_Q]V^M = R^M{}_{NPQ}V^N\)), with two indices the
  Ricci tensor \(R_{MN} = R^P{}_{MPN}\), without indices the scalar
  curvature \(R = g^{MN}R_{MN}\); with frame indices
  \(R^{(a)(b)}{}_{(c)(d)}\) its frame components,
  \(R^{(a)(b)}{}_{(a)(b)}\) (no sum) the sectional curvature of the
  frame plane \((a, b)\). Unit: 1/length\(^2\). Status: STANDARD.
\end{itemize}

The Lovelock densities are built with the generalised Kronecker delta
from \(p\) Riemann tensors. Their key property (Lovelock's theorem,
STANDARD) is that their metric variation contains no derivative of the
metric higher than the second, so that the field equations FE-1 are
second-order differential equations, and that it is linear in the second
derivatives along any one direction (quasi-linearity): a thin source at
the fold produces a finite jump of first derivatives, the junction
conditions FE-4. Status: STANDARD.

\hypertarget{the-field-dirac16complex}{%
\subsection{The field dirac16complex}\label{the-field-dirac16complex}}

\[
S_\Psi = \int\mathrm{d}^8X\,\sqrt{\lvert g\rvert}\,\Big[\frac12\big(\bar\Psi\gamma^MD_M\Psi - (D_M\bar\Psi)\gamma^M\Psi\big)
 - m_\Psi\,\bar\Psi\Psi - \frac{\lambda_\Psi}{2}\big(\bar\Psi\Psi\big)^2\Big] \tag{IN-4}
\] \textbf{Symbols in (IN-4).}

\begin{itemize}
\tightlist
\item
  \(S_\Psi\): The action of dirac16complex in the bulk (kinetic term
  with the canonical spin connection, mass term, contact
  self-interaction). Unit: dimensionless. Status: RECORD.
\item
  \(\mathrm{d}\): The differential: \(\mathrm{d}X\) is an infinitesimal
  change of \(X\), \(\mathrm{d}X/\mathrm{d}Y\) the derivative of \(X\)
  with respect to \(Y\), \(\mathrm{d}^nX\) the volume element of \(n\)
  variables. Status: STANDARD.
\item
  \(X = \{x_0, \dots, x_7\}\): The author's coordinates of the
  eight-dimensional spacetime, \(X = \{x_0, x_1, \dots, x_7\}\) with the
  hidden direction first (Dirac\_Einstein-Lovelock\_4+4\_Nash.nb,
  section DEFCHART); \(X^M\) is the coordinate with the index \(M\) and
  \(\mathrm{d}^8X = \mathrm{d}x_0\,\mathrm{d}x_1\cdots\mathrm{d}x_7\)
  the coordinate volume element. Unit: length (each coordinate). Status:
  AUTHOR.
\item
  \(g_{MN}\): The metric tensor of the eight-dimensional spacetime
  (signature \(+,+,+,+,-,-,-,-\) in the order \(x_0, \dots, x_7\)):
  \(g_{MN}\,\mathrm{d}X^M\mathrm{d}X^N\) is the squared interval,
  \(g^{MN}\) its inverse, \(\lvert g\rvert\) the absolute value of its
  determinant (\(\sqrt{\lvert g\rvert}\,\mathrm{d}^8X\) is the proper
  volume element). Unit: dimensionless. Status: AUTHOR.
\item
  \(\Psi,\ \bar\Psi\): dirac16complex: the author's 16-component complex
  spinor field of eight dimensions;
  \(\bar\Psi = \Psi^\dagger\sigma_{16}\) is its Dirac adjoint; in a term
  of the fold, \(\Psi\) is its value at \(\sigma = 0\). Unit:
  GeV\(^{7/2}\). Status: AUTHOR.
\item
  \(\gamma^M,\ \gamma^{(a)}\): The Dirac matrices of dirac16complex
  along the coordinates, \(\gamma^M = \mathbf{e}_a{}^MT^a\),
  \(\gamma_M = g_{MN}\gamma^N\); \(\gamma^{(a)} = T^a\) along the frame
  direction \((a)\) (\(\gamma^{(0)}\) along \(\sigma\), \(\gamma^{(4)}\)
  along the time). Unit: dimensionless. Status: AUTHOR.
\item
  \(M\): A coordinate index of eight dimensions (also with a numbered
  copy \(M_1, M_2, \dots\)): it runs over \(x_0, x_1, \dots, x_7\); a
  repeated index is summed. Status: STANDARD.
\item
  \(D_M,\ D_\mu\): The covariant derivative: on dirac16complex
  \(D_M = \partial_M + \Omega_M\) (with
  \(D_M\bar\Psi = \partial_M\bar\Psi - \bar\Psi\Omega_M\)); on the
  fields of the fold the gauge and spin covariant derivative \(D_\mu\)
  of IN-13; \(D^\mu = \hat h^{\mu\nu}D_\nu\); on an adjoint field
  \((D_\mu X)^A = \partial_\mu X^A + g_3f^{ABC}G^B_\mu X^C\) (likewise
  for \(SU(2)\) with \(g_2\epsilon^{IJK}\)). Unit: 1/length. Status:
  STANDARD.
\item
  \(m_\Psi\): The mass parameter of dirac16complex in the bulk. Unit:
  GeV. Status: RECORD.
\item
  \(\lambda_\Psi\): The coupling of the contact self-interaction of
  dirac16complex. Unit: GeV\(^{-6}\). Status: RECORD.
\end{itemize}

dirac16complex \(\Psi\) is the author's complex spinor field of eight
dimensions with 16 components. Its Lagrangian is that of the record: the
kinetic term with the canonical spin connection (in the symmetric form,
which is real without a factor \(\mathrm{i}\) for the author's real
matrices, IN-5), the mass term and a contact self-interaction. It
couples to gravity through \(\sqrt{\lvert g\rvert}\), \(\gamma^M\) and
the spin connection \(\Omega_M\) inside \(D_M\) (IN-6), and to itself
through \(\lambda_\Psi\).
\(\bar\Psi\Psi = \bar\Psi_{\mathrm{L}}\Psi_{\mathrm{L}} + \bar\Psi_{\mathrm{R}}\Psi_{\mathrm{R}}\)
(the mass term does not mix the halves, IN-5), while the kinetic term
mixes them (it contains one Dirac matrix). Status: RECORD.

\[
\begin{aligned}
&\{T_a, T_b\} = 2\eta_{ab}\mathbf{1},\qquad \eta = \mathrm{diag}(1, 1, 1, 1, -1, -1, -1, -1),\qquad T^a = \eta^{ab}T_b,
 \qquad \gamma^M = \mathbf{e}_a{}^M\,T^a,\\
&\bar\Psi = \Psi^\dagger\sigma_{16},\qquad \sigma_{16} = T_0T_1T_2T_3,\qquad \Gamma_8 = T_0T_1T_2T_3T_4T_5T_6T_7,\\
&P_{\mathrm{L}} = \frac{\mathbf{1} - \Gamma_8}{2},\qquad P_{\mathrm{R}} = \frac{\mathbf{1} + \Gamma_8}{2},\qquad
 \Psi_{\mathrm{L}} = P_{\mathrm{L}}\Psi,\qquad \Psi_{\mathrm{R}} = P_{\mathrm{R}}\Psi
\end{aligned} \tag{IN-5}
\] \textbf{Symbols in (IN-5).}

\begin{itemize}
\tightlist
\item
  \(T_a,\ T^a\): The author's eight real \(16 \times 16\) Dirac matrices
  T16\^{}A{[}a{]} of 4+4.nb (\(a = 0, \dots, 7\); \(T_0, \dots, T_7\)),
  with \(\{T_a, T_b\} = 2\eta_{ab}\mathbf{1}\); \(T^a = \eta^{ab}T_b\).
  Unit: dimensionless. Status: AUTHOR.
\item
  \(a\): A frame (orthonormal) index running over \(0, 1, \dots, 7\)
  (frame direction \((0)\) along \(\sigma\), \((1), (2), (3)\) along
  ordinary space, \((4)\) along the time, \((5), (6), (7)\) along the
  extra times); raised and lowered with \(\eta_{ab}\); on the fold
  (\(a, b \in \{1, \dots, 7\}\)) the frame of the fold; a repeated index
  is summed. Status: STANDARD.
\item
  \(b\): A frame index (as \(a\)). Status: STANDARD.
\item
  \(\eta_{ab}\): The flat metric of the frame,
  \(\eta = \mathrm{diag}(1, 1, 1, 1, -1, -1, -1, -1)\) (eta4488 of
  4+4.nb); \(\eta^{ab}\) is its inverse (the same matrix). Unit:
  dimensionless. Status: AUTHOR.
\item
  \(\mathbf{1}\): The identity matrix of the size fixed by the context
  (\(16 \times 16\) on dirac16complex, \(4 \times 4\) on a
  four-dimensional spinor, \(3 \times 3\) in generation or flavour
  space). Unit: dimensionless. Status: STANDARD.
\item
  \(\mathrm{diag}\): A diagonal matrix with the listed entries on its
  diagonal (in the order of the coordinates \(x_0, x_1, \dots, x_7\) for
  a tensor of eight dimensions, \(x_1, \dots, x_7\) for a tensor of the
  fold, \(x_1, \dots, x_4\) for a four-dimensional tensor). Status:
  STANDARD.
\item
  \(\gamma^M,\ \gamma^{(a)}\): The Dirac matrices of dirac16complex
  along the coordinates, \(\gamma^M = \mathbf{e}_a{}^MT^a\),
  \(\gamma_M = g_{MN}\gamma^N\); \(\gamma^{(a)} = T^a\) along the frame
  direction \((a)\) (\(\gamma^{(0)}\) along \(\sigma\), \(\gamma^{(4)}\)
  along the time). Unit: dimensionless. Status: AUTHOR.
\item
  \(M\): A coordinate index of eight dimensions (also with a numbered
  copy \(M_1, M_2, \dots\)): it runs over \(x_0, x_1, \dots, x_7\); a
  repeated index is summed. Status: STANDARD.
\item
  \(\mathbf{e}_a{}^M,\ \mathbf{e}^a{}_M\): The vielbein (frame field) of
  the metric: \(g_{MN} = \eta_{ab}\,\mathbf{e}^a{}_M\mathbf{e}^b{}_N\);
  \(\mathbf{e}_a{}^M\) is its inverse;
  \(\mathbf{e}_{bN} = \eta_{bc}\mathbf{e}^c{}_N\); \(\mathbf{e}^{(a)}\)
  is the frame one-form \(\mathbf{e}^a{}_M\,\mathrm{d}X^M\). Unit:
  dimensionless. Status: STANDARD.
\item
  \(\Psi,\ \bar\Psi\): dirac16complex: the author's 16-component complex
  spinor field of eight dimensions;
  \(\bar\Psi = \Psi^\dagger\sigma_{16}\) is its Dirac adjoint; in a term
  of the fold, \(\Psi\) is its value at \(\sigma = 0\). Unit:
  GeV\(^{7/2}\). Status: AUTHOR.
\item
  \({}^\dagger\): Hermitian conjugation: the complex conjugate of the
  transpose of a matrix or spinor (for field operators the adjoint
  operator). Status: STANDARD.
\item
  \(\sigma_{16}\): The matrix of the Dirac adjoint,
  \(\sigma_{16} = T_0T_1T_2T_3\), real and symmetric; \(\sigma_{16}T_a\)
  is real and antisymmetric. Unit: dimensionless. Status: AUTHOR.
\item
  \(\Gamma_8\): The chirality matrix of eight dimensions,
  \(\Gamma_8 = T_0T_1\cdots T_7\) (the author's T16\^{}A{[}8{]})
  \(= \mathrm{diag}(-\mathbf{1}_8, \mathbf{1}_8)\);
  \(\Gamma_8^2 = \mathbf{1}\); it anticommutes with every \(T_a\). Unit:
  dimensionless. Status: AUTHOR.
\item
  \(P_{\mathrm{L}}\): The author's left-handed chiral projector,
  \(P_{\mathrm{L}} = (\mathbf{1} - \Gamma_8)/2\). Unit: dimensionless.
  Status: AUTHOR.
\item
  \(P_{\mathrm{R}}\): The author's right-handed chiral projector,
  \(P_{\mathrm{R}} = (\mathbf{1} + \Gamma_8)/2\). Unit: dimensionless.
  Status: AUTHOR.
\item
  \(\Psi_{\mathrm{L}}\): The left-handed half of dirac16complex,
  \(\Psi_{\mathrm{L}} = P_{\mathrm{L}}\Psi\) (the upper eight components
  in the basis of 4+4.nb). Unit: GeV\(^{7/2}\). Status: AUTHOR.
\item
  \(\Psi_{\mathrm{R}}\): The right-handed half of dirac16complex,
  \(\Psi_{\mathrm{R}} = P_{\mathrm{R}}\Psi\) (the lower eight
  components). Unit: GeV\(^{7/2}\). Status: AUTHOR.
\end{itemize}

\(T_a\) (\(a = 0, \dots, 7\)) are the author's eight real
\(16 \times 16\) Dirac matrices \texttt{T16\^{}A{[}a{]}} of
\texttt{4+4.nb}, and \(\Gamma_8\) is the author's chirality matrix
\texttt{T16\^{}A{[}8{]}}
\(= \mathrm{diag}(-\mathbf{1}_8, \mathbf{1}_8)\); \(P_{\mathrm{L}}\) and
\(P_{\mathrm{R}}\) are the author's chiral projectors (the provenance
document \texttt{dirac\ matrices.md} of UM-3 proves the Clifford
relations, the reality and the projectors; the facts used here are
restated in this paragraph). The Dirac adjoint uses
\(\sigma_{16} = T_0T_1T_2T_3\): \(\sigma_{16}\) is real and symmetric,
\(\sigma_{16}T_a\) is real and antisymmetric, and \(\sigma_{16}\)
commutes with \(\Gamma_8\) (it is a product of four Dirac matrices), so
\(\overline{P_X\Psi} = \bar\Psi P_X\) and every bilinear with an even
number of Dirac matrices joins a half with itself, every bilinear with
an odd number joins \(\Psi_{\mathrm{L}}\) with \(\Psi_{\mathrm{R}}\)
(DERIVED, \texttt{um3\_fields.py}, checks F1 and F2). Status: AUTHOR
(matrices, projectors), RECORD (adjoint).

\[
\begin{aligned}
&D_M\Psi = \partial_M\Psi + \Omega_M\Psi,\qquad D_M\bar\Psi = \partial_M\bar\Psi - \bar\Psi\,\Omega_M,\qquad
 \Omega_M = \frac12\,\omega_{Mab}\,\Sigma^{ab},\qquad \Sigma^{ab} = \frac14\big[T^a, T^b\big],\\
&\omega_{Mab} = \mathbf{e}_a{}^N\big(\partial_M\mathbf{e}_{bN} - \Gamma^P{}_{MN}\,\mathbf{e}_{bP}\big)
\end{aligned} \tag{IN-6}
\] \textbf{Symbols in (IN-6).}

\begin{itemize}
\tightlist
\item
  \(D_M,\ D_\mu\): The covariant derivative: on dirac16complex
  \(D_M = \partial_M + \Omega_M\) (with
  \(D_M\bar\Psi = \partial_M\bar\Psi - \bar\Psi\Omega_M\)); on the
  fields of the fold the gauge and spin covariant derivative \(D_\mu\)
  of IN-13; \(D^\mu = \hat h^{\mu\nu}D_\nu\); on an adjoint field
  \((D_\mu X)^A = \partial_\mu X^A + g_3f^{ABC}G^B_\mu X^C\) (likewise
  for \(SU(2)\) with \(g_2\epsilon^{IJK}\)). Unit: 1/length. Status:
  STANDARD.
\item
  \(M\): A coordinate index of eight dimensions (also with a numbered
  copy \(M_1, M_2, \dots\)): it runs over \(x_0, x_1, \dots, x_7\); a
  repeated index is summed. Status: STANDARD.
\item
  \(\partial\): A partial derivative: \(\partial_M\) with respect to the
  coordinate \(X^M\), \(\partial_\mu\) with respect to \(x^\mu\),
  \(\partial_\sigma\) with respect to \(\sigma\),
  \(\partial X/\partial Y\) the derivative of \(X\) with respect to
  \(Y\) at fixed other variables. Status: STANDARD.
\item
  \(\Psi,\ \bar\Psi\): dirac16complex: the author's 16-component complex
  spinor field of eight dimensions;
  \(\bar\Psi = \Psi^\dagger\sigma_{16}\) is its Dirac adjoint; in a term
  of the fold, \(\Psi\) is its value at \(\sigma = 0\). Unit:
  GeV\(^{7/2}\). Status: AUTHOR.
\item
  \(\Omega_M\): The spin connection acting on dirac16complex,
  \(\Omega_M = \frac12\omega_{Mab}\Sigma^{ab}\). Unit: 1/length. Status:
  RECORD.
\item
  \(\omega_{Mab}\): The canonical (Levi-Civita) spin connection of the
  vielbein, antisymmetric in \(a, b\). Unit: 1/length. Status: STANDARD.
\item
  \(a\): A frame (orthonormal) index running over \(0, 1, \dots, 7\)
  (frame direction \((0)\) along \(\sigma\), \((1), (2), (3)\) along
  ordinary space, \((4)\) along the time, \((5), (6), (7)\) along the
  extra times); raised and lowered with \(\eta_{ab}\); on the fold
  (\(a, b \in \{1, \dots, 7\}\)) the frame of the fold; a repeated index
  is summed. Status: STANDARD.
\item
  \(b\): A frame index (as \(a\)). Status: STANDARD.
\item
  \(\Sigma^{ab},\ \Sigma^{\mu\nu}\): The generators of Spin(4,4) on
  dirac16complex, \(\Sigma^{ab} = \frac14[T^a, T^b]\) (the scaled
  commutators of the author's matrices);
  \(\Sigma^{\mu\nu} = \mathbf{e}_a{}^\mu\mathbf{e}_b{}^\nu\Sigma^{ab}\)
  along the fold. Unit: dimensionless. Status: AUTHOR.
\item
  \(T_a,\ T^a\): The author's eight real \(16 \times 16\) Dirac matrices
  T16\^{}A{[}a{]} of 4+4.nb (\(a = 0, \dots, 7\); \(T_0, \dots, T_7\)),
  with \(\{T_a, T_b\} = 2\eta_{ab}\mathbf{1}\); \(T^a = \eta^{ab}T_b\).
  Unit: dimensionless. Status: AUTHOR.
\item
  \(\mathbf{e}_a{}^M,\ \mathbf{e}^a{}_M\): The vielbein (frame field) of
  the metric: \(g_{MN} = \eta_{ab}\,\mathbf{e}^a{}_M\mathbf{e}^b{}_N\);
  \(\mathbf{e}_a{}^M\) is its inverse;
  \(\mathbf{e}_{bN} = \eta_{bc}\mathbf{e}^c{}_N\); \(\mathbf{e}^{(a)}\)
  is the frame one-form \(\mathbf{e}^a{}_M\,\mathrm{d}X^M\). Unit:
  dimensionless. Status: STANDARD.
\item
  \(N\): A coordinate index of eight dimensions (also
  \(N_1, N_2, \dots\)), running over \(x_0, \dots, x_7\); a repeated
  index is summed. Status: STANDARD.
\item
  \(\Gamma^P{}_{MN}\): The Christoffel symbols of the metric \(g_{MN}\)
  (the Levi-Civita connection). Unit: 1/length. Status: STANDARD.
\item
  \(P\): A coordinate index of eight dimensions (also
  \(P_1, P_2, \dots\)), running over \(x_0, \dots, x_7\); a repeated
  index is summed. Status: STANDARD.
\end{itemize}

The canonical spin connection \(\omega_{Mab}\) of the vielbein
(antisymmetric in \(a, b\)) acts on the spinor through the generators
\(\Sigma^{ab}\) of Spin(4,4). With it the Dirac matrices are covariantly
constant (\(D_M\gamma^N = 0\)), the kinetic term of IN-4 is a scalar and
the field equation FE-12 is covariant. For the author's metric
\(\gamma^M\Omega_M\) is never zero (FE-14). Status: RECORD (the
canonical construction), DERIVED (the vielbein postulate for the
author's metric, \texttt{um3\_geometry.py}, G5).

\hypertarget{the-stabiliser-of-the-bulk}{%
\subsection{The stabiliser of the
bulk}\label{the-stabiliser-of-the-bulk}}

\[
\mathcal{T}_{\mathrm{stab}}{}^a{}_b = \mathrm{diag}\big(p_{\mathrm{s}}, p_{\mathrm{s}}, p_{\mathrm{s}}, p_{\mathrm{s}},
 -\rho_{\mathrm{s}}, p_{\mathrm{s}}, p_{\mathrm{s}}, p_{\mathrm{s}}\big),\qquad
\kappa\rho_{\mathrm{s}} = -\frac{H^2}{16},\qquad \kappa p_{\mathrm{s}} = -\frac{5H^2}{16} \tag{IN-7}
\] \textbf{Symbols in (IN-7).}

\begin{itemize}
\tightlist
\item
  \(\mathcal{T}_{\mathrm{stab}}{}^a{}_b\): The energy-momentum tensor of
  the stabiliser of the bulk (frame components). Unit: GeV\(^8\).
  Status: ASSUMED.
\item
  \(a\): A frame (orthonormal) index running over \(0, 1, \dots, 7\)
  (frame direction \((0)\) along \(\sigma\), \((1), (2), (3)\) along
  ordinary space, \((4)\) along the time, \((5), (6), (7)\) along the
  extra times); raised and lowered with \(\eta_{ab}\); on the fold
  (\(a, b \in \{1, \dots, 7\}\)) the frame of the fold; a repeated index
  is summed. Status: STANDARD.
\item
  \(b\): A frame index (as \(a\)). Status: STANDARD.
\item
  \(\mathrm{diag}\): A diagonal matrix with the listed entries on its
  diagonal (in the order of the coordinates \(x_0, x_1, \dots, x_7\) for
  a tensor of eight dimensions, \(x_1, \dots, x_7\) for a tensor of the
  fold, \(x_1, \dots, x_4\) for a four-dimensional tensor). Status:
  STANDARD.
\item
  \(p_{\mathrm{s}}\): The pressure of the stabiliser (equal along
  \(\sigma\), ordinary space and the extra times). Unit: GeV\(^8\).
  Value: \(\kappa p_{\mathrm{s}} = -5H^2/16\). Status: ASSUMED.
\item
  \(\rho_{\mathrm{s}}\): The energy density of the stabiliser. Unit:
  GeV\(^8\). Value: \(\kappa\rho_{\mathrm{s}} = -H^2/16\). Status:
  ASSUMED.
\item
  \(\kappa\): The eight-dimensional gravitational coupling,
  \(\kappa = 8\pi G_8\). Unit: GeV\(^{-6}\). Status: AUTHOR.
\item
  \(H\): The author's constant (an inverse length) that fixes the period
  \(\pi/(6H)\) of \(x_0\), the Lovelock couplings and the cosmological
  constant \(\Lambda = -10H^2\); in UM-3 a free parameter. Unit: GeV
  (1/length). Status: AUTHOR.
\end{itemize}

\(S_{\mathrm{stab}}\) is the action of a medium of the bulk; only its
energy-momentum tensor \(\mathcal{T}_{\mathrm{stab}}\) enters the field
equations, and IN-7 states it in the frame of the author's metric. With
the author's weights and \(\Lambda\), the bulk of a universe whose
ordinary space expands slowly (the late universe, \(a_4' \ll H\)) is not
a vacuum solution: its frame curvature is that of the static block
(FE-32), which requires exactly this source, a negative energy density
\(\rho_{\mathrm{s}} = -H^2/(16\kappa)\) with the isotropic pressure
\(p_{\mathrm{s}} = 5\rho_{\mathrm{s}}\) (DERIVED,
\texttt{um3\_geometry.py}, G4). The stabiliser couples only to gravity.
It is absent in the primordial era, whose bulk is the author's vacuum
(FE-31). Status: ASSUMED (its existence; its value is DERIVED from the
requirement of a static bulk); its microscopic origin is OPEN.

\hypertarget{the-fold}{%
\subsection{The fold}\label{the-fold}}

\[
\begin{aligned}
&\int\mathrm{d}^8X\,\sqrt{\lvert g\rvert}\;\delta_{\mathrm{D}}(\sigma)\,\frac{1}{V_*}\,F
 = \int\mathrm{d}^4x\,\sqrt{-\hat h}\;\mathcal{Z}\,F\big\vert_{\sigma=0},\qquad \mathrm{d}^4x = \mathrm{d}x_1\,\mathrm{d}x_2\,\mathrm{d}x_3\,\mathrm{d}x_4,\\
&\hat h_{\mu\nu} = \mathrm{diag}\big(\mathrm{e}^{2a_4}, \mathrm{e}^{2a_4}, \mathrm{e}^{2a_4}, -1\big),\qquad
 \sqrt{-\hat h} = \mathrm{e}^{3a_4},\qquad
 \mathcal{Z} = \frac{V_{\mathrm{t}}\big(\mathrm{e}^{-a_4}b_4\big)^3}{V_*},\qquad V_{\mathrm{t}} = \int\mathrm{d}x_5\,\mathrm{d}x_6\,\mathrm{d}x_7
\end{aligned} \tag{IN-8}
\] \textbf{Symbols in (IN-8).}

\begin{itemize}
\tightlist
\item
  \(\mathrm{d}\): The differential: \(\mathrm{d}X\) is an infinitesimal
  change of \(X\), \(\mathrm{d}X/\mathrm{d}Y\) the derivative of \(X\)
  with respect to \(Y\), \(\mathrm{d}^nX\) the volume element of \(n\)
  variables. Status: STANDARD.
\item
  \(X = \{x_0, \dots, x_7\}\): The author's coordinates of the
  eight-dimensional spacetime, \(X = \{x_0, x_1, \dots, x_7\}\) with the
  hidden direction first (Dirac\_Einstein-Lovelock\_4+4\_Nash.nb,
  section DEFCHART); \(X^M\) is the coordinate with the index \(M\) and
  \(\mathrm{d}^8X = \mathrm{d}x_0\,\mathrm{d}x_1\cdots\mathrm{d}x_7\)
  the coordinate volume element. Unit: length (each coordinate). Status:
  AUTHOR.
\item
  \(g_{MN}\): The metric tensor of the eight-dimensional spacetime
  (signature \(+,+,+,+,-,-,-,-\) in the order \(x_0, \dots, x_7\)):
  \(g_{MN}\,\mathrm{d}X^M\mathrm{d}X^N\) is the squared interval,
  \(g^{MN}\) its inverse, \(\lvert g\rvert\) the absolute value of its
  determinant (\(\sqrt{\lvert g\rvert}\,\mathrm{d}^8X\) is the proper
  volume element). Unit: dimensionless. Status: AUTHOR.
\item
  \(\delta_{\mathrm{D}}\): The Dirac delta function:
  \(\int\mathrm{d}\sigma\,\delta_{\mathrm{D}}(\sigma)F(\sigma) = F(0)\).
  Unit: 1/length (for \(\delta_{\mathrm{D}}(\sigma)\)). Status:
  STANDARD.
\item
  \(\sigma\): The proper coordinate of the hidden direction:
  \(\mathrm{d}\sigma = \lvert\cot z\rvert\,\mathrm{d}x_0\),
  \(\sigma = \mathrm{sgn}(\pi/2 - z)\ln(\sin z)/(6H)\) on one period
  \(0 < z < \pi\); \(\sigma = 0\) is the fold \(z = \pi/2\),
  \(\sigma \to \mp\infty\) the tips \(z \to 0, \pi\); one period of
  \(x_0\) is the whole line \(-\infty < \sigma < \infty\). Unit: length.
  Status: DERIVED.
\item
  \(V_*\): The physical volume of the extra times \(x_5, x_6, x_7\) on
  the fold in the late universe,
  \(V_* = V_{\mathrm{t}}(\mathrm{e}^{-a_4}b_4)^3\) at a time when
  \(\mathrm{e}^{-a_4}b_4\) is constant; it converts the
  eight-dimensional normalisation into the four-dimensional one. Unit:
  length\(^3\) (GeV\(^{-3}\)). Status: ASSUMED.
\item
  \(F\): Any function of the fields at the fold (used to state a rule of
  integration). Status: STANDARD.
\item
  \(x\): The coordinates of the fold: \(x_i\) (\(i = 1, 2, 3\)) of
  ordinary space, \(x_4\) the time, \(x_r\) (\(r = 5, 6, 7\)) the extra
  times; a field of the fold is a function of the four observable
  coordinates \(x = (x_1, x_2, x_3, x_4)\), and
  \(\mathrm{d}^4x = \mathrm{d}x_1\,\mathrm{d}x_2\,\mathrm{d}x_3\,\mathrm{d}x_4\).
  Unit: length (each coordinate). Status: AUTHOR.
\item
  \(x_1\): The first coordinate of ordinary 3-space. Unit: length.
  Status: AUTHOR.
\item
  \(x_2\): The second coordinate of ordinary 3-space. Unit: length.
  Status: AUTHOR.
\item
  \(x_3\): The third coordinate of ordinary 3-space. Unit: length.
  Status: AUTHOR.
\item
  \(x_4\): The time (\(x_4 \ge 0\) in the ad hoc vacuum); a prime on a
  function of \(x_4\) alone is \(\mathrm{d}/\mathrm{d}x_4\). Unit:
  length (= time, \(c = 1\)). Status: AUTHOR.
\item
  \(\hat h_{\mu\nu}\): The four-dimensional metric of the fold,
  \(\hat h_{\mu\nu} = \mathrm{diag}(\mathrm{e}^{2a_4}, \mathrm{e}^{2a_4}, \mathrm{e}^{2a_4}, -1)\)
  in the coordinates \(x_1, \dots, x_4\) (the induced metric at
  \(\sigma = 0\) without the extra times); \(\hat h^{\mu\nu}\) its
  inverse, \(\hat h\) its determinant,
  \(\sqrt{-\hat h} = \mathrm{e}^{3a_4}\); \(\hat h^{(1)}_{\mu\nu}\) its
  first-order cosmological perturbation. Unit: dimensionless. Status:
  DERIVED.
\item
  \(\mu\): A four-dimensional coordinate index running over the
  observable coordinates \(x_1, x_2, x_3, x_4\) of the fold (signature
  \(+,+,+,-\)); a repeated index is summed. Status: STANDARD.
\item
  \(\nu\): A four-dimensional coordinate index (as \(\mu\)). Status:
  STANDARD.
\item
  \(\mathrm{diag}\): A diagonal matrix with the listed entries on its
  diagonal (in the order of the coordinates \(x_0, x_1, \dots, x_7\) for
  a tensor of eight dimensions, \(x_1, \dots, x_7\) for a tensor of the
  fold, \(x_1, \dots, x_4\) for a four-dimensional tensor). Status:
  STANDARD.
\item
  \(\mathrm{e}\): Euler's number, the base of the natural logarithm;
  \(\mathrm{e}^{X}\) is the exponential of \(X\). Unit: dimensionless.
  Value: 2.718281828. Status: STANDARD.
\item
  \(a_4\): The author's metric function of the time \(x_4\) that scales
  ordinary 3-space: on the fold \(\mathrm{e}^{a_4}\) is the scale factor
  of cosmology and \(a_4'\) the Hubble rate. Unit: dimensionless.
  Status: AUTHOR.
\item
  \(\mathcal{Z}\): The fold factor,
  \(\mathcal{Z}(x_4) = V_{\mathrm{t}}(\mathrm{e}^{-a_4}b_4)^3/V_*\): the
  factor by which the size of the extra times multiplies every term of
  the fold; \(\mathcal{Z} = 1\) when the extra times are frozen at their
  late value. Unit: dimensionless. Value: 1 in the late universe.
  Status: DERIVED.
\item
  \(V_{\mathrm{t}}\): The coordinate volume of the (compact) extra
  times,
  \(V_{\mathrm{t}} = \int\mathrm{d}x_5\,\mathrm{d}x_6\,\mathrm{d}x_7\).
  Unit: length\(^3\). Status: ASSUMED.
\item
  \(b_4\): The author's metric function of the time \(x_4\) in the
  extra-time components, \(b_4 > 0\), never identically 1 (the author's
  order of 2026-10-24); the scale factor of the extra times is
  \(\mathrm{e}^{-a_4}b_4\) times the warp. Unit: dimensionless. Status:
  AUTHOR.
\item
  \(x_5\): The first of the three extra time-like coordinates. Unit:
  length. Status: AUTHOR.
\item
  \(x_6\): The second of the three extra time-like coordinates. Unit:
  length. Status: AUTHOR.
\item
  \(x_7\): The third of the three extra time-like coordinates. Unit:
  length. Status: AUTHOR.
\end{itemize}

The fold is the seven-dimensional surface \(\sigma = 0\) (\(z = \pi/2\))
with the coordinates \(x_1, \dots, x_7\). No field depends on the extra
times (ASSUMED, FE-38), so the integral over them gives their coordinate
volume \(V_{\mathrm{t}}\) (the extra times are compact) times their
scale factors \((\mathrm{e}^{-a_4}b_4)^3\) at the fold. \(V_*\) is the
physical volume of the extra times on the fold in the late universe,
where \(\mathrm{e}^{-a_4}b_4\) is constant (FE-33); with this
normalisation \(\mathcal{Z} = 1\) in the whole late universe, the fields
of the fold are four-dimensional fields in the four-dimensional metric
\(\hat h_{\mu\nu}\) of the fold, and their couplings are the measured
four-dimensional ones. \(\mathcal{Z}(x_4)\) is the factor by which the
extra times rescale every coupling of the fold when
\(\mathrm{e}^{-a_4}b_4\) changes (FE-44). \(F\) stands for any function
of the fields on the fold. Status: ASSUMED (the place of the fields),
DERIVED (the reduction).

\[
\mathcal{L}_{\mathrm{IG}} = \frac{\hat R}{2\kappa_4} \tag{IN-9}
\] \textbf{Symbols in (IN-9).}

\begin{itemize}
\tightlist
\item
  \(\mathcal{L}_{\mathrm{IG}}\): The Lagrangian of the induced gravity
  of the fold. Unit: GeV\(^4\). Status: ASSUMED.
\item
  \(\hat R\): The scalar curvature of the metric \(\hat h_{\mu\nu}\) of
  the fold. Unit: 1/length\(^2\). Status: STANDARD.
\item
  \(\kappa_4\): The coupling of the induced gravity of the fold,
  \(\kappa_4 = 8\pi G_{\mathrm{ind}}\); four-dimensional gravity on the
  fold has \(8\pi G_{\mathrm{N}} = \kappa_4/(1 + 5r/6)\) (FE-59). Unit:
  GeV\(^{-2}\). Value:
  \(\kappa_4 \simeq 8\pi G_{\mathrm{N}} = 8\pi/M_{\mathrm{Pl}}^2\).
  Status: ASSUMED.
\end{itemize}

The induced gravity of the fold: an Einstein-Hilbert term of the
four-dimensional metric \(\hat h_{\mu\nu}\) of the fold (\(\hat R\) its
scalar curvature), with the coupling
\(\kappa_4 = 8\pi G_{\mathrm{ind}}\). Such a term is generated by the
quantum fluctuations of every field that lives on the fold (induced
gravity); UM-3 assumes that it dominates four-dimensional gravity,
because the gravitons of the bulk move with the wrong speed (DERIVED,
\texttt{um3\_bulk.py}, B5: \(c_{\mathrm{g}}^2 = 15/2\) at the static
block). With it, four-dimensional gravity on the fold is general
relativity with Newton's constant \(G_{\mathrm{N}}\) (FE-59), up to a
fraction \(r\) of bulk gravity. Status: ASSUMED.

\[
\mathcal{T}_{\mathrm{ff}}{}^a{}_b = \mathrm{diag}\big(p_{\mathrm{ff}}, p_{\mathrm{ff}}, p_{\mathrm{ff}}, -\rho_{\mathrm{ff}},
 q_{\mathrm{ff}}, q_{\mathrm{ff}}, q_{\mathrm{ff}}\big) \tag{IN-10}
\] \textbf{Symbols in (IN-10).}

\begin{itemize}
\tightlist
\item
  \(\mathcal{T}_{\mathrm{ff}}{}^a{}_b\): The energy-momentum tensor of
  the fold fluid (frame components of the fold, four-dimensional
  normalisation). Unit: GeV\(^4\). Status: ASSUMED.
\item
  \(a\): A frame (orthonormal) index running over \(0, 1, \dots, 7\)
  (frame direction \((0)\) along \(\sigma\), \((1), (2), (3)\) along
  ordinary space, \((4)\) along the time, \((5), (6), (7)\) along the
  extra times); raised and lowered with \(\eta_{ab}\); on the fold
  (\(a, b \in \{1, \dots, 7\}\)) the frame of the fold; a repeated index
  is summed. Status: STANDARD.
\item
  \(b\): A frame index (as \(a\)). Status: STANDARD.
\item
  \(\mathrm{diag}\): A diagonal matrix with the listed entries on its
  diagonal (in the order of the coordinates \(x_0, x_1, \dots, x_7\) for
  a tensor of eight dimensions, \(x_1, \dots, x_7\) for a tensor of the
  fold, \(x_1, \dots, x_4\) for a four-dimensional tensor). Status:
  STANDARD.
\item
  \(p_{\mathrm{ff}}\): The pressure of the fold fluid along ordinary
  space. Unit: GeV\(^4\). Value: \(-61HV_*/(8\kappa)\) at the static
  block. Status: ASSUMED.
\item
  \(\rho_{\mathrm{ff}}\): The energy density of the fold fluid. Unit:
  GeV\(^4\). Value: \(31HV_*/(4\kappa)\) at the static block. Status:
  ASSUMED.
\item
  \(q_{\mathrm{ff}}\): The pressure of the fold fluid along the extra
  times. Unit: GeV\(^4\). Value: \(-61HV_*/(8\kappa)\) at the static
  block. Status: ASSUMED.
\end{itemize}

The fold fluid: a medium at rest on the fold (the frame indices \(a, b\)
run over the seven directions \(1, \dots, 7\) of the fold), with the
energy density \(\rho_{\mathrm{ff}}\), the pressure \(p_{\mathrm{ff}}\)
along ordinary space and the pressure \(q_{\mathrm{ff}}\) along the
extra times (four-dimensional normalisation, IN-8). The kink of the
author's metric at the fold requires a thin source (FE-4); the fold
fluid is the part of it that the Standard Model and the portals do not
supply: at the static block \(\rho_{\mathrm{ff}} = 31HV_*/(4\kappa)\),
\(p_{\mathrm{ff}} = q_{\mathrm{ff}} = -61HV_*/(8\kappa)\) (FE-32), at
the author's vacuum \(\rho_{\mathrm{ff}} = 6HV_*/\kappa\),
\(p_{\mathrm{ff}} = q_{\mathrm{ff}} = -10HV_*/\kappa\) less the
contribution of the induced gravity (FE-31). It plays the role of the
brane tension of brane-world models. \(\mathcal{L}_{\mathrm{ff}}\)
stands for its Lagrangian; only \(\mathcal{T}_{\mathrm{ff}}\) enters the
field equations. Status: ASSUMED (its existence; its values DERIVED from
the junction conditions); its microscopic origin is OPEN.

\hypertarget{the-higgs-doublet-with-the-ad-hoc-vacuum}{%
\subsection{The Higgs doublet with the ad hoc
vacuum}\label{the-higgs-doublet-with-the-ad-hoc-vacuum}}

\[
\mathcal{L}_{\mathrm{H}} = -\hat h^{\mu\nu}\big(D_\mu\Phi\big)^\dagger D_\nu\Phi - V(\Phi),\qquad
V(\Phi) = \mu^2\,\Phi^\dagger\Phi + \lambda\big(\Phi^\dagger\Phi\big)^2,\qquad \lambda > 0,\quad \mu^2 < 0 \tag{IN-11}
\] \textbf{Symbols in (IN-11).}

\begin{itemize}
\tightlist
\item
  \(\mathcal{L}_{\mathrm{H}}\): The Lagrangian of the Higgs doublet on
  the fold. Unit: GeV\(^4\). Status: AUTHOR.
\item
  \(\hat h_{\mu\nu}\): The four-dimensional metric of the fold,
  \(\hat h_{\mu\nu} = \mathrm{diag}(\mathrm{e}^{2a_4}, \mathrm{e}^{2a_4}, \mathrm{e}^{2a_4}, -1)\)
  in the coordinates \(x_1, \dots, x_4\) (the induced metric at
  \(\sigma = 0\) without the extra times); \(\hat h^{\mu\nu}\) its
  inverse, \(\hat h\) its determinant,
  \(\sqrt{-\hat h} = \mathrm{e}^{3a_4}\); \(\hat h^{(1)}_{\mu\nu}\) its
  first-order cosmological perturbation. Unit: dimensionless. Status:
  DERIVED.
\item
  \(\mu\): A four-dimensional coordinate index running over the
  observable coordinates \(x_1, x_2, x_3, x_4\) of the fold (signature
  \(+,+,+,-\)); a repeated index is summed. Status: STANDARD.
\item
  \(\nu\): A four-dimensional coordinate index (as \(\mu\)). Status:
  STANDARD.
\item
  \(D_M,\ D_\mu\): The covariant derivative: on dirac16complex
  \(D_M = \partial_M + \Omega_M\) (with
  \(D_M\bar\Psi = \partial_M\bar\Psi - \bar\Psi\Omega_M\)); on the
  fields of the fold the gauge and spin covariant derivative \(D_\mu\)
  of IN-13; \(D^\mu = \hat h^{\mu\nu}D_\nu\); on an adjoint field
  \((D_\mu X)^A = \partial_\mu X^A + g_3f^{ABC}G^B_\mu X^C\) (likewise
  for \(SU(2)\) with \(g_2\epsilon^{IJK}\)). Unit: 1/length. Status:
  STANDARD.
\item
  \(\Phi\): The Higgs doublet (two complex components, hypercharge
  \(1/2\)), a field of the fold. Unit: GeV. Status: AUTHOR.
\item
  \({}^\dagger\): Hermitian conjugation: the complex conjugate of the
  transpose of a matrix or spinor (for field operators the adjoint
  operator). Status: STANDARD.
\item
  \(V(\Phi)\): The Higgs potential,
  \(V = \mu^2\Phi^\dagger\Phi + \lambda(\Phi^\dagger\Phi)^2\). Unit:
  GeV\(^4\). Status: AUTHOR.
\item
  \(\mu^2\): The mass parameter of the Higgs potential, \(\mu^2 < 0\)
  (it enters only as \(\mu^2\)). Unit: GeV\(^2\). Value:
  \(\mu^2 = -\lambda v^2 = -m_h^2/2 = -(88.6\ \mathrm{GeV})^2\). Status:
  AUTHOR.
\item
  \(\lambda\): The quartic coupling of the Higgs potential,
  \(\lambda > 0\). Unit: dimensionless. Value:
  \(\lambda = m_h^2/(2v^2) = 0.129\). Status: AUTHOR.
\end{itemize}

The Higgs doublet \(\Phi\) (two complex components, hypercharge \(1/2\))
with the author's potential. In the signature \((+,+,+,-)\) of the fold
the kinetic term carries a minus sign, so that the energy density is
positive. Status: AUTHOR (the potential), STANDARD (the rest).

\[
\Phi = \frac{1}{\sqrt2}\begin{pmatrix}\mathcal{A}(x_4)\\ v + h\end{pmatrix},\qquad
\mathcal{A}(x_4) = v_{\mathcal{A}}\,\mathrm{e}^{-2a_{40}x_4},\qquad a_{40} > 0,\quad x_4 \ge 0,\qquad v^2 = -\frac{\mu^2}{\lambda} \tag{IN-12}
\] \textbf{Symbols in (IN-12).}

\begin{itemize}
\tightlist
\item
  \(\Phi\): The Higgs doublet (two complex components, hypercharge
  \(1/2\)), a field of the fold. Unit: GeV. Status: AUTHOR.
\item
  \(\mathcal{A}\): The author's ad hoc upper component of the Higgs
  vacuum, almightyAdHoc,
  \(\mathcal{A}(x_4) = v_{\mathcal{A}}\mathrm{e}^{-2a_{40}x_4}\). Unit:
  GeV. Status: AUTHOR.
\item
  \(x_4\): The time (\(x_4 \ge 0\) in the ad hoc vacuum); a prime on a
  function of \(x_4\) alone is \(\mathrm{d}/\mathrm{d}x_4\). Unit:
  length (= time, \(c = 1\)). Status: AUTHOR.
\item
  \(v\): The vacuum expectation value of the lower component of the
  Higgs doublet, \(v^2 = -\mu^2/\lambda\). Unit: GeV. Value: 246.22 GeV.
  Status: STANDARD.
\item
  \(h\): The Higgs field: the fluctuation of the lower component of the
  Higgs doublet about \(v\). Unit: GeV. Value: \(m_h = 125.25\) GeV.
  Status: STANDARD.
\item
  \(v_{\mathcal{A}}\): The amplitude of almightyAdHoc at \(x_4 = 0\);
  the author's almightyAdHoc \(= \mathrm{Exp}[-2\,a_{40}\,x_4]\) has
  \(v_{\mathcal{A}} = 1\) in the unit of the field. Unit: GeV. Value: 1
  (author). Status: AUTHOR.
\item
  \(\mathrm{e}\): Euler's number, the base of the natural logarithm;
  \(\mathrm{e}^{X}\) is the exponential of \(X\). Unit: dimensionless.
  Value: 2.718281828. Status: STANDARD.
\item
  \(a_{40}\): The decay constant of almightyAdHoc, \(a_{40} > 0\):
  \(\mathcal{A}\) falls by \(\mathrm{e}\) in the time \(1/(2a_{40})\).
  Unit: GeV (1/length). Status: AUTHOR.
\item
  \(\mu^2\): The mass parameter of the Higgs potential, \(\mu^2 < 0\)
  (it enters only as \(\mu^2\)). Unit: GeV\(^2\). Value:
  \(\mu^2 = -\lambda v^2 = -m_h^2/2 = -(88.6\ \mathrm{GeV})^2\). Status:
  AUTHOR.
\item
  \(\lambda\): The quartic coupling of the Higgs potential,
  \(\lambda > 0\). Unit: dimensionless. Value:
  \(\lambda = m_h^2/(2v^2) = 0.129\). Status: AUTHOR.
\end{itemize}

The author's ad hoc vacuum: the upper component of the Higgs doublet is
the prescribed function almightyAdHoc
\(\mathcal{A}(x_4) = \exp(-2a_{40}x_4)\) (the author writes it with the
amplitude \(v_{\mathcal{A}} = 1\) in the unit of the field; UM-3 keeps
\(v_{\mathcal{A}}\) to show where it enters), the lower component is the
vacuum expectation value \(v\) plus the Higgs field \(h\). The Goldstone
fields are absorbed by the gauge choice (unitary gauge). \(\mathcal{A}\)
decays with the rate \(2a_{40}\) from the beginning of the time
\(x_4 = 0\): it acts only in the primordial era. The ad hoc vacuum is
not an exact solution of the Higgs equation; the residuals are FE-26.
Status: AUTHOR.

\hypertarget{the-gauge-fields}{%
\subsection{The gauge fields}\label{the-gauge-fields}}

\[
\begin{aligned}
D_\mu &= \partial_\mu + \hat\omega_\mu - \mathrm{i}g_3\,G^A_\mu\,t^A - \mathrm{i}g_2\,W^I_\mu\,\frac{\tau^I}{2}
 - \mathrm{i}g_1\,Y_f\,B_\mu,\\
\mathcal{L}_{\mathrm{gauge}} &= -\frac14\,G^A_{\mu\nu}G^{A\mu\nu} - \frac14\,W^I_{\mu\nu}W^{I\mu\nu} - \frac14\,B_{\mu\nu}B^{\mu\nu},\\
G^A_{\mu\nu} &= \partial_\mu G^A_\nu - \partial_\nu G^A_\mu + g_3\,f^{ABC}G^B_\mu G^C_\nu,\qquad
W^I_{\mu\nu} = \partial_\mu W^I_\nu - \partial_\nu W^I_\mu + g_2\,\epsilon^{IJK}W^J_\mu W^K_\nu,\qquad
B_{\mu\nu} = \partial_\mu B_\nu - \partial_\nu B_\mu
\end{aligned} \tag{IN-13}
\] \textbf{Symbols in (IN-13).}

\begin{itemize}
\tightlist
\item
  \(D_M,\ D_\mu\): The covariant derivative: on dirac16complex
  \(D_M = \partial_M + \Omega_M\) (with
  \(D_M\bar\Psi = \partial_M\bar\Psi - \bar\Psi\Omega_M\)); on the
  fields of the fold the gauge and spin covariant derivative \(D_\mu\)
  of IN-13; \(D^\mu = \hat h^{\mu\nu}D_\nu\); on an adjoint field
  \((D_\mu X)^A = \partial_\mu X^A + g_3f^{ABC}G^B_\mu X^C\) (likewise
  for \(SU(2)\) with \(g_2\epsilon^{IJK}\)). Unit: 1/length. Status:
  STANDARD.
\item
  \(\mu\): A four-dimensional coordinate index running over the
  observable coordinates \(x_1, x_2, x_3, x_4\) of the fold (signature
  \(+,+,+,-\)); a repeated index is summed. Status: STANDARD.
\item
  \(\nu\): A four-dimensional coordinate index (as \(\mu\)). Status:
  STANDARD.
\item
  \(\partial\): A partial derivative: \(\partial_M\) with respect to the
  coordinate \(X^M\), \(\partial_\mu\) with respect to \(x^\mu\),
  \(\partial_\sigma\) with respect to \(\sigma\),
  \(\partial X/\partial Y\) the derivative of \(X\) with respect to
  \(Y\) at fixed other variables. Status: STANDARD.
\item
  \(\hat\omega_\mu\): The spin connection of the four-dimensional metric
  \(\hat h_{\mu\nu}\) of the fold acting on four-dimensional spinors.
  Unit: 1/length. Status: STANDARD.
\item
  \(\mathrm{i}\): The imaginary unit, \(\mathrm{i}^2 = -1\). Unit:
  dimensionless. Status: STANDARD.
\item
  \(g_3\): The \(SU(3)_C\) gauge coupling (the strong coupling \(g_s\)).
  Unit: dimensionless. Value: \(g_3(m_Z) = 1.2177\)
  (\(\alpha_s = 0.1180\)). Status: STANDARD.
\item
  \(G^A_\mu,\ G^A_{\mu\nu}\): The gluon field of \(SU(3)_C\) on the fold
  (\(A = 1, \dots, 8\)) and its field strength. Unit: GeV (field),
  GeV\(^2\) (field strength). Status: AUTHOR.
\item
  \(A\): A colour (adjoint SU(3)) index running over \(1, \dots, 8\); a
  repeated index is summed. Status: STANDARD.
\item
  \(B\): A colour (adjoint SU(3)) index (as \(A\)). Status: STANDARD.
\item
  \(C\): A colour (adjoint SU(3)) index (as \(A\)). Status: STANDARD.
\item
  \(t^A\): The generators of \(SU(3)\) on colour triplets,
  \(t^A = \lambda_{\mathrm{GM}}^A/2\) (the Gell-Mann matrices halved),
  \(\mathrm{tr}(t^At^B) = \delta^{AB}/2\). Unit: dimensionless. Status:
  STANDARD.
\item
  \(g_2\): The \(SU(2)_L\) gauge coupling (usually \(g\)). Unit:
  dimensionless. Value: \(g_2(m_Z) = 0.6517\). Status: STANDARD.
\item
  \(W^I_\mu,\ W^I_{\mu\nu}\): The gauge field of \(SU(2)_L\) on the
  fold, \(W^1_\mu, W^2_\mu, W^3_\mu\), and its field strength. Unit: GeV
  (field), GeV\(^2\) (field strength). Status: AUTHOR.
\item
  \(I\): A weak-isospin (adjoint SU(2)) index running over \(1, 2, 3\);
  a repeated index is summed. Status: STANDARD.
\item
  \(J\): A weak-isospin (adjoint SU(2)) index (as \(I\)). Status:
  STANDARD.
\item
  \(K\): A weak-isospin (adjoint SU(2)) index (as \(I\)). Status:
  STANDARD.
\item
  \(\tau^I\): The Pauli matrices acting on weak doublets
  (\(\tau^1, \tau^2, \tau^3\); \(\tau^2\) is the second one, not a
  square). Unit: dimensionless. Status: STANDARD.
\item
  \(g_1\): The \(U(1)_Y\) gauge coupling (usually \(g'\); not normalised
  for grand unification). Unit: dimensionless. Value:
  \(g_1(m_Z) = 0.3574\). Status: STANDARD.
\item
  \(Y_f\): The hypercharge of a field (IN-13: \(1/6\), \(2/3\),
  \(-1/3\), \(-1/2\), \(-1\), \(0\) for
  \(Q, u_{\mathrm{R}}, d_{\mathrm{R}}, \ell, e_{\mathrm{R}}, \nu_{\mathrm{R}}\),
  \(1/2\) for \(\Phi\)). Unit: dimensionless. Status: STANDARD.
\item
  \(f\): The label of a kind of Standard-Model multiplet:
  \(f \in \{Q, u, d, \ell, e, \nu\}\) (quark doublets, up-quark
  singlets, down-quark singlets, lepton doublets, charged-lepton
  singlets, right-handed neutrinos); a sum over \(f\) runs over the six.
  Status: STANDARD.
\item
  \(B_\mu,\ B_{\mu\nu}\): The hypercharge gauge field of \(U(1)_Y\) on
  the fold and its field strength. Unit: GeV (field), GeV\(^2\) (field
  strength). Status: AUTHOR.
\item
  \(\mathcal{L}_{\mathrm{gauge}}\): The Lagrangian of the gauge fields
  of \(SU(3)_C\times SU(2)_L\times U(1)_Y\) on the fold. Unit:
  GeV\(^4\). Status: STANDARD.
\item
  \(f^{ABC}\): The structure constants of \(SU(3)\),
  \([t^A, t^B] = \mathrm{i}f^{ABC}t^C\). Unit: dimensionless. Status:
  STANDARD.
\item
  \(\epsilon^{IJK}\): The Levi-Civita symbol of three indices (the
  structure constants of \(SU(2)\)), \(\epsilon^{123} = 1\). Unit:
  dimensionless. Status: STANDARD.
\end{itemize}

The gauge fields of \(SU(3)_C \times SU(2)_L \times U(1)_Y\) on the
fold, with the canonical normalisation (the couplings in the covariant
derivative): the gluons \(G^A_\mu\), the weak-isospin fields
\(W^1_\mu, W^2_\mu, W^3_\mu\) and the hypercharge field \(B_\mu\). The
covariant derivative of a field acts with the part that fits the field:
\(t^A\) on colour triplets, \(\tau^I/2\) on weak doublets, the
hypercharge \(Y_f\) of the field, and the four-dimensional spin
connection \(\hat\omega_\mu\) on spinors. The hypercharges (STANDARD)
are \(Y = 1/6\) for \(Q\), \(2/3\) for \(u_{\mathrm{R}}\), \(-1/3\) for
\(d_{\mathrm{R}}\), \(-1/2\) for \(\ell\), \(-1\) for
\(e_{\mathrm{R}}\), \(0\) for \(\nu_{\mathrm{R}}\) and \(1/2\) for
\(\Phi\). The gauge fields have no components along the extra times
(ASSUMED: such components would be scalars of wrong sign). Status:
STANDARD.

\hypertarget{the-quarks-and-leptons}{%
\subsection{The quarks and leptons}\label{the-quarks-and-leptons}}

\[
\begin{aligned}
&\mathcal{L}_{\mathrm{f}} = \sum_{\psi}\mathrm{i}\,\bar\psi\,\hat\gamma^\mu D_\mu\psi,\qquad
 \psi \in \{Q, u_{\mathrm{R}}, d_{\mathrm{R}}, \ell, e_{\mathrm{R}}, \nu_{\mathrm{R}}\},\\
&\{\hat\gamma^\mu, \hat\gamma^\nu\} = -2\hat h^{\mu\nu}\mathbf{1},\qquad \bar\psi = \psi^\dagger\beta,\qquad
 \beta = \hat\gamma^{(4)},\qquad \hat\gamma_5 = \mathrm{i}\,\hat\gamma^{(4)}\hat\gamma^{(1)}\hat\gamma^{(2)}\hat\gamma^{(3)},\\
&\hat P_{\mathrm{L}} = \frac{\mathbf{1} - \hat\gamma_5}{2},\quad \hat P_{\mathrm{R}} = \frac{\mathbf{1} + \hat\gamma_5}{2},\qquad
 \hat P_{\mathrm{L}}Q = Q,\quad \hat P_{\mathrm{L}}\ell = \ell,\quad \hat P_{\mathrm{R}}u_{\mathrm{R}} = u_{\mathrm{R}},\quad
 \hat P_{\mathrm{R}}d_{\mathrm{R}} = d_{\mathrm{R}},\quad \hat P_{\mathrm{R}}e_{\mathrm{R}} = e_{\mathrm{R}},\quad
 \hat P_{\mathrm{R}}\nu_{\mathrm{R}} = \nu_{\mathrm{R}}
\end{aligned} \tag{IN-14}
\] \textbf{Symbols in (IN-14).}

\begin{itemize}
\tightlist
\item
  \(\mathcal{L}_{\mathrm{f}}\): The kinetic Lagrangian of the quarks and
  leptons on the fold. Unit: GeV\(^4\). Status: STANDARD.
\item
  \(\psi\): The label of a Standard-Model fermion field in a sum:
  \(\psi\) runs over
  \(Q, u_{\mathrm{R}}, d_{\mathrm{R}}, \ell, e_{\mathrm{R}}, \nu_{\mathrm{R}}\).
  Status: STANDARD.
\item
  \(\mathrm{i}\): The imaginary unit, \(\mathrm{i}^2 = -1\). Unit:
  dimensionless. Status: STANDARD.
\item
  \(\psi,\ \bar\psi\): A Standard-Model fermion field (a
  four-dimensional Weyl field of the fold with three generations,
  written as a four-component spinor of fixed chirality);
  \(\bar\psi = \psi^\dagger\beta\); \(\psi_f\) the fields of the kind
  \(f\). Unit: GeV\(^{3/2}\). Status: STANDARD.
\item
  \(\hat\gamma^\mu\): The four-dimensional complex \(4 \times 4\) Dirac
  matrices of the fold,
  \(\hat\gamma^\mu = \hat{\mathbf{e}}_a{}^\mu\hat\gamma^{(a)}\),
  \(\{\hat\gamma^\mu, \hat\gamma^\nu\} = -2\hat h^{\mu\nu}\mathbf{1}\);
  \(\hat\gamma_\mu = \hat h_{\mu\nu}\hat\gamma^\nu\). Unit:
  dimensionless. Status: STANDARD.
\item
  \(\mu\): A four-dimensional coordinate index running over the
  observable coordinates \(x_1, x_2, x_3, x_4\) of the fold (signature
  \(+,+,+,-\)); a repeated index is summed. Status: STANDARD.
\item
  \(D_M,\ D_\mu\): The covariant derivative: on dirac16complex
  \(D_M = \partial_M + \Omega_M\) (with
  \(D_M\bar\Psi = \partial_M\bar\Psi - \bar\Psi\Omega_M\)); on the
  fields of the fold the gauge and spin covariant derivative \(D_\mu\)
  of IN-13; \(D^\mu = \hat h^{\mu\nu}D_\nu\); on an adjoint field
  \((D_\mu X)^A = \partial_\mu X^A + g_3f^{ABC}G^B_\mu X^C\) (likewise
  for \(SU(2)\) with \(g_2\epsilon^{IJK}\)). Unit: 1/length. Status:
  STANDARD.
\item
  \(Q,\ \bar Q\): The left-handed quark doublets of the three
  generations, \(Q = (u_{\mathrm{L}}, d_{\mathrm{L}})\) (colour
  triplets, hypercharge \(1/6\)). Unit: GeV\(^{3/2}\). Status: AUTHOR.
\item
  \(u_{\mathrm{R}}\): The right-handed up-type quarks \(u, c, t\)
  (singlets of \(SU(2)_L\), hypercharge \(2/3\)). Unit: GeV\(^{3/2}\).
  Status: STANDARD.
\item
  \(d_{\mathrm{R}}\): The right-handed down-type quarks \(d, s, b\)
  (hypercharge \(-1/3\)). Unit: GeV\(^{3/2}\). Status: STANDARD.
\item
  \(\ell,\ \bar\ell\): The left-handed lepton doublets of the three
  generations, \(\ell = (\nu_{\mathrm{L}}, e_{\mathrm{L}})\)
  (hypercharge \(-1/2\)). Unit: GeV\(^{3/2}\). Status: AUTHOR.
\item
  \(e_{\mathrm{R}}\): The right-handed electron, muon and tau (singlets,
  hypercharge \(-1\)). Unit: GeV\(^{3/2}\). Status: AUTHOR.
\item
  \(\nu_{\mathrm{R}}\): The three right-handed neutrinos (gauge
  singlets). Unit: GeV\(^{3/2}\). Status: ASSUMED.
\item
  \(\nu\): A four-dimensional coordinate index (as \(\mu\)). Status:
  STANDARD.
\item
  \(\hat h_{\mu\nu}\): The four-dimensional metric of the fold,
  \(\hat h_{\mu\nu} = \mathrm{diag}(\mathrm{e}^{2a_4}, \mathrm{e}^{2a_4}, \mathrm{e}^{2a_4}, -1)\)
  in the coordinates \(x_1, \dots, x_4\) (the induced metric at
  \(\sigma = 0\) without the extra times); \(\hat h^{\mu\nu}\) its
  inverse, \(\hat h\) its determinant,
  \(\sqrt{-\hat h} = \mathrm{e}^{3a_4}\); \(\hat h^{(1)}_{\mu\nu}\) its
  first-order cosmological perturbation. Unit: dimensionless. Status:
  DERIVED.
\item
  \(\mathbf{1}\): The identity matrix of the size fixed by the context
  (\(16 \times 16\) on dirac16complex, \(4 \times 4\) on a
  four-dimensional spinor, \(3 \times 3\) in generation or flavour
  space). Unit: dimensionless. Status: STANDARD.
\item
  \({}^\dagger\): Hermitian conjugation: the complex conjugate of the
  transpose of a matrix or spinor (for field operators the adjoint
  operator). Status: STANDARD.
\item
  \(\beta\): The matrix of the four-dimensional Dirac adjoint,
  \(\beta = \hat\gamma^{(4)}\) (Hermitian, \(\beta^2 = \mathbf{1}\)).
  Unit: dimensionless. Status: STANDARD.
\item
  \(\hat\gamma_5\): The four-dimensional chirality matrix,
  \(\hat\gamma_5 = \mathrm{i}\hat\gamma^{(4)}\hat\gamma^{(1)}\hat\gamma^{(2)}\hat\gamma^{(3)}\),
  \(\hat\gamma_5^2 = \mathbf{1}\). Unit: dimensionless. Status:
  STANDARD.
\item
  \(\hat P_{\mathrm{L}}\): The four-dimensional left-handed projector,
  \(\hat P_{\mathrm{L}} = (\mathbf{1} - \hat\gamma_5)/2\). Unit:
  dimensionless. Status: STANDARD.
\item
  \(\hat P_{\mathrm{R}}\): The four-dimensional right-handed projector,
  \(\hat P_{\mathrm{R}} = (\mathbf{1} + \hat\gamma_5)/2\). Unit:
  dimensionless. Status: STANDARD.
\end{itemize}

Every Standard-Model fermion is a four-dimensional Weyl field on the
fold, written as a four-component spinor with a fixed chirality: the
quark doublets \(Q\) and the lepton doublets \(\ell\) are left-handed,
the singlets \(u_{\mathrm{R}}\), \(d_{\mathrm{R}}\), \(e_{\mathrm{R}}\),
\(\nu_{\mathrm{R}}\) right-handed. Each symbol stands for the three
generations (a vector in generation space) and, for quarks, the three
colours. \(\hat\gamma^\mu = \hat{\mathbf{e}}_a{}^\mu\hat\gamma^{(a)}\)
are complex \(4 \times 4\) Dirac matrices of the fold (with the frame
matrices \(\hat\gamma^{(a)}\), \(a = 1, \dots, 4\),
\((\hat\gamma^{(4)})^2 = \mathbf{1}\),
\((\hat\gamma^{(1)})^2 = (\hat\gamma^{(2)})^2 = (\hat\gamma^{(3)})^2 = -\mathbf{1}\));
\(\beta = \hat\gamma^{(4)}\) is Hermitian and
\((\hat\gamma^\mu)^\dagger = \beta\hat\gamma^\mu\beta\). With these
conventions the free Dirac equation is
\((-\mathrm{i}\hat\gamma^\mu\partial_\mu + m)\psi = 0\) and
\(\bar\psi\beta\psi\)-type energy densities are positive. Status:
STANDARD.

\[
\mathcal{L}_{\mathrm{Y}} = -\big(\bar Q\,\mathcal{Y}^{u}\,\tilde\Phi\,u_{\mathrm{R}} + \bar Q\,\mathcal{Y}^{d}\,\Phi\,d_{\mathrm{R}}
 + \bar\ell\,\mathcal{Y}^{e}\,\Phi\,e_{\mathrm{R}} + \bar\ell\,\mathcal{Y}^{\nu}\,\tilde\Phi\,\nu_{\mathrm{R}}\big)
 + \mathrm{h.c.},\qquad \tilde\Phi = \mathrm{i}\tau^2\Phi^* \tag{IN-15}
\] \textbf{Symbols in (IN-15).}

\begin{itemize}
\tightlist
\item
  \(\mathcal{L}_{\mathrm{Y}}\): The Yukawa Lagrangian of the quarks and
  leptons with the Higgs doublet. Unit: GeV\(^4\). Status: STANDARD.
\item
  \(Q,\ \bar Q\): The left-handed quark doublets of the three
  generations, \(Q = (u_{\mathrm{L}}, d_{\mathrm{L}})\) (colour
  triplets, hypercharge \(1/6\)). Unit: GeV\(^{3/2}\). Status: AUTHOR.
\item
  \(\mathcal{Y}^{u},\ \mathcal{Y}^{d},\ \mathcal{Y}^{e},\ \mathcal{Y}^{\nu}\):
  The Yukawa matrices (\(3 \times 3\) complex matrices in generation
  space) of the up quarks, the down quarks, the charged leptons and the
  neutrinos. Unit: dimensionless. Status: STANDARD.
\item
  \(\tilde\Phi\): The conjugate Higgs doublet,
  \(\tilde\Phi = \mathrm{i}\tau^2\Phi^*\) (hypercharge \(-1/2\)). Unit:
  GeV. Status: STANDARD.
\item
  \(u_{\mathrm{R}}\): The right-handed up-type quarks \(u, c, t\)
  (singlets of \(SU(2)_L\), hypercharge \(2/3\)). Unit: GeV\(^{3/2}\).
  Status: STANDARD.
\item
  \(\Phi\): The Higgs doublet (two complex components, hypercharge
  \(1/2\)), a field of the fold. Unit: GeV. Status: AUTHOR.
\item
  \(d_{\mathrm{R}}\): The right-handed down-type quarks \(d, s, b\)
  (hypercharge \(-1/3\)). Unit: GeV\(^{3/2}\). Status: STANDARD.
\item
  \(\ell,\ \bar\ell\): The left-handed lepton doublets of the three
  generations, \(\ell = (\nu_{\mathrm{L}}, e_{\mathrm{L}})\)
  (hypercharge \(-1/2\)). Unit: GeV\(^{3/2}\). Status: AUTHOR.
\item
  \(e_{\mathrm{R}}\): The right-handed electron, muon and tau (singlets,
  hypercharge \(-1\)). Unit: GeV\(^{3/2}\). Status: AUTHOR.
\item
  \(\nu_{\mathrm{R}}\): The three right-handed neutrinos (gauge
  singlets). Unit: GeV\(^{3/2}\). Status: ASSUMED.
\item
  \(\mathrm{h.c.}\): The Hermitian conjugate of the preceding terms.
  Status: STANDARD.
\item
  \(\mathrm{i}\): The imaginary unit, \(\mathrm{i}^2 = -1\). Unit:
  dimensionless. Status: STANDARD.
\item
  \(\tau^I\): The Pauli matrices acting on weak doublets
  (\(\tau^1, \tau^2, \tau^3\); \(\tau^2\) is the second one, not a
  square). Unit: dimensionless. Status: STANDARD.
\item
  \({}^*\): Complex conjugation of a number, or of each entry of a
  matrix or column. Status: STANDARD.
\end{itemize}

The Yukawa couplings, with the complex \(3 \times 3\) matrices
\(\mathcal{Y}^{u}\), \(\mathcal{Y}^{d}\), \(\mathcal{Y}^{e}\),
\(\mathcal{Y}^{\nu}\) in generation space; \(\tilde\Phi\) is the
conjugate doublet (hypercharge \(-1/2\)). In the vacuum they give the
quarks, the charged leptons and the neutrinos their Dirac masses
(FE-46); with three right-handed neutrinos the neutrino masses are Dirac
masses and the lepton mixing is that of the PMNS matrix. No Majorana
mass of \(\nu_{\mathrm{R}}\) is included (lepton number is conserved).
Status: STANDARD (with \(\nu_{\mathrm{R}}\): accepted by the author,
2026-10-24 17:32).

\hypertarget{the-portals-of-dirac16complex}{%
\subsection{The portals of
dirac16complex}\label{the-portals-of-dirac16complex}}

\[
\begin{aligned}
\mathcal{L}_{\mathrm{P}} = {}&-c_\Phi\,\big(\Phi^\dagger\Phi\big)\,\bar\Psi P_{\mathrm{L}}\Psi
 - \sum_{X}\Big[\frac14\big(c^{G}_X\,G^A_{\mu\nu}G^{A\mu\nu} + c^{W}_X\,W^I_{\mu\nu}W^{I\mu\nu}
 + c^{B}_X\,B_{\mu\nu}B^{\mu\nu}\big)\,\bar\Psi P_X\Psi\\
&\qquad + \frac{d_X}{2}\,B_{\mu\nu}\,\mathcal{D}_X^{\mu\nu} + \sum_{f} c^{f}_X\,J_f^\mu\,j_{X\mu}\Big],\\
\mathcal{D}_X^{\mu\nu} = {}&\mathrm{i}\,\bar\Psi\,\Sigma^{\mu\nu}P_X\Psi,\qquad
j_X^\mu = \mathrm{i}\big(\bar\Psi P_X D^\mu\Psi - (D^\mu\bar\Psi)P_X\Psi\big),\qquad
J_f^\mu = \bar\psi_f\,\hat\gamma^\mu\psi_f
\end{aligned} \tag{IN-16}
\] \textbf{Symbols in (IN-16).}

\begin{itemize}
\tightlist
\item
  \(\mathcal{L}_{\mathrm{P}}\): The Lagrangian of the portals of
  dirac16complex at the fold. Unit: GeV\(^4\). Status: INVENTED.
\item
  \(c_\Phi\): The Higgs-portal coupling of the left-handed half of
  dirac16complex (the right-handed half has none). Unit: GeV\(^{-5}\).
  Status: INVENTED.
\item
  \(\Phi\): The Higgs doublet (two complex components, hypercharge
  \(1/2\)), a field of the fold. Unit: GeV. Status: AUTHOR.
\item
  \({}^\dagger\): Hermitian conjugation: the complex conjugate of the
  transpose of a matrix or spinor (for field operators the adjoint
  operator). Status: STANDARD.
\item
  \(\Psi,\ \bar\Psi\): dirac16complex: the author's 16-component complex
  spinor field of eight dimensions;
  \(\bar\Psi = \Psi^\dagger\sigma_{16}\) is its Dirac adjoint; in a term
  of the fold, \(\Psi\) is its value at \(\sigma = 0\). Unit:
  GeV\(^{7/2}\). Status: AUTHOR.
\item
  \(P_{\mathrm{L}}\): The author's left-handed chiral projector,
  \(P_{\mathrm{L}} = (\mathbf{1} - \Gamma_8)/2\). Unit: dimensionless.
  Status: AUTHOR.
\item
  \(X\): The label of a half of dirac16complex:
  \(X \in \{\mathrm{L}, \mathrm{R}\}\) (left-handed or right-handed
  half); a sum over \(X\) runs over both. Status: STANDARD.
\item
  \(c^{G}_X\): The gluon-kinetic portal coupling of the half \(X\) of
  dirac16complex. Unit: GeV\(^{-7}\). Status: INVENTED.
\item
  \(G^A_\mu,\ G^A_{\mu\nu}\): The gluon field of \(SU(3)_C\) on the fold
  (\(A = 1, \dots, 8\)) and its field strength. Unit: GeV (field),
  GeV\(^2\) (field strength). Status: AUTHOR.
\item
  \(A\): A colour (adjoint SU(3)) index running over \(1, \dots, 8\); a
  repeated index is summed. Status: STANDARD.
\item
  \(\mu\): A four-dimensional coordinate index running over the
  observable coordinates \(x_1, x_2, x_3, x_4\) of the fold (signature
  \(+,+,+,-\)); a repeated index is summed. Status: STANDARD.
\item
  \(\nu\): A four-dimensional coordinate index (as \(\mu\)). Status:
  STANDARD.
\item
  \(c^{W}_X\): The \(W\)-kinetic portal coupling of the half \(X\).
  Unit: GeV\(^{-7}\). Status: INVENTED.
\item
  \(W^I_\mu,\ W^I_{\mu\nu}\): The gauge field of \(SU(2)_L\) on the
  fold, \(W^1_\mu, W^2_\mu, W^3_\mu\), and its field strength. Unit: GeV
  (field), GeV\(^2\) (field strength). Status: AUTHOR.
\item
  \(I\): A weak-isospin (adjoint SU(2)) index running over \(1, 2, 3\);
  a repeated index is summed. Status: STANDARD.
\item
  \(c^{B}_X\): The \(B\)-kinetic portal coupling of the half \(X\).
  Unit: GeV\(^{-7}\). Status: INVENTED.
\item
  \(B_\mu,\ B_{\mu\nu}\): The hypercharge gauge field of \(U(1)_Y\) on
  the fold and its field strength. Unit: GeV (field), GeV\(^2\) (field
  strength). Status: AUTHOR.
\item
  \(P_X\): The chiral projector of the half \(X\):
  \(P_X = P_{\mathrm{L}}\) for \(X = \mathrm{L}\), \(P_{\mathrm{R}}\)
  for \(X = \mathrm{R}\). Unit: dimensionless. Status: AUTHOR.
\item
  \(d_X\): The dipole-portal coupling of the half \(X\) of
  dirac16complex with the hypercharge field strength. Unit:
  GeV\(^{-5}\). Status: INVENTED.
\item
  \(\mathcal{D}_X^{\mu\nu}\): The dipole (antisymmetric tensor) density
  of the half \(X\) of dirac16complex at the fold,
  \(\mathcal{D}_X^{\mu\nu} = \mathrm{i}\bar\Psi\Sigma^{\mu\nu}P_X\Psi\)
  (real, DERIVED). Unit: GeV\(^7\). Status: DERIVED.
\item
  \(f\): The label of a kind of Standard-Model multiplet:
  \(f \in \{Q, u, d, \ell, e, \nu\}\) (quark doublets, up-quark
  singlets, down-quark singlets, lepton doublets, charged-lepton
  singlets, right-handed neutrinos); a sum over \(f\) runs over the six.
  Status: STANDARD.
\item
  \(c^{f}_X\): The current-portal coupling of the half \(X\) of
  dirac16complex with the fermion-number current of the multiplets of
  kind \(f\) (flavour blind). Unit: GeV\(^{-7}\). Status: INVENTED.
\item
  \(J_f^\mu\): The fermion-number current of the Standard-Model
  multiplets of kind \(f\), \(J_f^\mu = \bar\psi_f\hat\gamma^\mu\psi_f\)
  summed over the three generations and the gauge components; with other
  labels the gauge currents (FE-16 to FE-18). Unit: GeV\(^3\). Status:
  STANDARD.
\item
  \(j_X^\mu\): The convective current of the half \(X\) of
  dirac16complex,
  \(j_X^\mu = \mathrm{i}(\bar\Psi P_XD^\mu\Psi - (D^\mu\bar\Psi)P_X\Psi)\)
  (real and joining the half with itself, DERIVED). Unit: GeV\(^8\).
  Status: DERIVED.
\item
  \(\mathrm{i}\): The imaginary unit, \(\mathrm{i}^2 = -1\). Unit:
  dimensionless. Status: STANDARD.
\item
  \(\Sigma^{ab},\ \Sigma^{\mu\nu}\): The generators of Spin(4,4) on
  dirac16complex, \(\Sigma^{ab} = \frac14[T^a, T^b]\) (the scaled
  commutators of the author's matrices);
  \(\Sigma^{\mu\nu} = \mathbf{e}_a{}^\mu\mathbf{e}_b{}^\nu\Sigma^{ab}\)
  along the fold. Unit: dimensionless. Status: AUTHOR.
\item
  \(D_M,\ D_\mu\): The covariant derivative: on dirac16complex
  \(D_M = \partial_M + \Omega_M\) (with
  \(D_M\bar\Psi = \partial_M\bar\Psi - \bar\Psi\Omega_M\)); on the
  fields of the fold the gauge and spin covariant derivative \(D_\mu\)
  of IN-13; \(D^\mu = \hat h^{\mu\nu}D_\nu\); on an adjoint field
  \((D_\mu X)^A = \partial_\mu X^A + g_3f^{ABC}G^B_\mu X^C\) (likewise
  for \(SU(2)\) with \(g_2\epsilon^{IJK}\)). Unit: 1/length. Status:
  STANDARD.
\item
  \(\psi,\ \bar\psi\): A Standard-Model fermion field (a
  four-dimensional Weyl field of the fold with three generations,
  written as a four-component spinor of fixed chirality);
  \(\bar\psi = \psi^\dagger\beta\); \(\psi_f\) the fields of the kind
  \(f\). Unit: GeV\(^{3/2}\). Status: STANDARD.
\item
  \(\hat\gamma^\mu\): The four-dimensional complex \(4 \times 4\) Dirac
  matrices of the fold,
  \(\hat\gamma^\mu = \hat{\mathbf{e}}_a{}^\mu\hat\gamma^{(a)}\),
  \(\{\hat\gamma^\mu, \hat\gamma^\nu\} = -2\hat h^{\mu\nu}\mathbf{1}\);
  \(\hat\gamma_\mu = \hat h_{\mu\nu}\hat\gamma^\nu\). Unit:
  dimensionless. Status: STANDARD.
\end{itemize}

The new interactions of UM-3 (INVENTED), the order of the author that
the left-handed and the right-handed halves of dirac16complex each
couple to every Standard-Model field with its own coupling, and that
only the left-handed half couples to the Higgs doublet. \(\Psi\) is the
value of dirac16complex at the fold. Every term is a product of a
gauge-invariant Standard-Model operator and a Lorentz-tensor bilinear of
one half \(X\) of dirac16complex, so every term is invariant under
\(SU(3)_C \times SU(2)_L \times U(1)_Y\) and under the Lorentz
transformations of the fold; none changes the flavour of a quark or a
lepton (consistent with QCD and with the electroweak theory, DERIVED by
construction):

\begin{itemize}
\tightlist
\item
  the Higgs portal: the scalar density \(\bar\Psi P_{\mathrm{L}}\Psi\)
  of the left-handed half times \(\Phi^\dagger\Phi\), with the coupling
  \(c_\Phi\); the right-handed half has no Higgs coupling;
\item
  the gauge-kinetic portals: the scalar density of each half times the
  gauge-invariant kinetic densities of the gluons, of the \(W\) fields
  and of the \(B\) field, with the six couplings \(c^{G}_X\),
  \(c^{W}_X\), \(c^{B}_X\): they make the gauge couplings of the fold
  depend on the density of dirac16complex at the fold (FE-16 to FE-18)
  and couple the bound states of dirac16complex to gluons, photons and
  weak bosons;
\item
  the dipole portal: the antisymmetric tensor \(\mathcal{D}_X^{\mu\nu}\)
  of each half times the hypercharge field strength (the only
  Standard-Model field strength that is gauge invariant by itself), with
  the couplings \(d_X\): an electric and magnetic dipole moment of
  dirac16complex with respect to the photon and the \(Z\);
\item
  the current portals: the convective current \(j_X^\mu\) of each half
  (the only real vector current that joins a half with itself:
  \(\bar\Psi\gamma^\mu P_X\Psi\) joins \(\Psi_{\mathrm{L}}\) with
  \(\Psi_{\mathrm{R}}\) and is not real, DERIVED,
  \texttt{um3\_fields.py}, F1) times the fermion-number current
  \(J_f^\mu\) of each kind \(f\) of Standard-Model multiplet, with the
  twelve couplings \(c^{f}_X\); \(J_f^\mu\) sums over the three
  generations and over the gauge components, so the portal is flavour
  blind.
\end{itemize}

The 21 couplings are \(c_\Phi\), \(c^{G}_{\mathrm{L}}\),
\(c^{G}_{\mathrm{R}}\), \(c^{W}_{\mathrm{L}}\), \(c^{W}_{\mathrm{R}}\),
\(c^{B}_{\mathrm{L}}\), \(c^{B}_{\mathrm{R}}\), \(d_{\mathrm{L}}\),
\(d_{\mathrm{R}}\) and \(c^{f}_{\mathrm{L}}\), \(c^{f}_{\mathrm{R}}\)
for \(f = Q, u, d, \ell, e, \nu\). Every term is bilinear in \(\Psi\):
the action is invariant under
\(\Psi \to \mathrm{e}^{\mathrm{i}\theta}\Psi\) (the charge of
dirac16complex is conserved, and its lightest state is stable), and
baryon and lepton numbers are conserved. Status: INVENTED.

\hypertarget{gravity-1}{%
\section{Gravity}\label{gravity-1}}

\hypertarget{the-einstein-lovelock-equations-in-the-bulk}{%
\subsection{The Einstein-Lovelock equations in the
bulk}\label{the-einstein-lovelock-equations-in-the-bulk}}

\[
\sum_{p=1}^{3}\alpha_{(p)}\,\mathcal{E}_{(p)}{}^M{}_N + \Lambda\,\delta^M_N
 = \kappa\big(\mathcal{T}_\Psi{}^M{}_N + \mathcal{T}_{\mathrm{stab}}{}^M{}_N\big)\qquad (\sigma \ne 0) \tag{FE-1}
\] \textbf{Symbols in (FE-1).}

\begin{itemize}
\tightlist
\item
  \(p\): The Lovelock order \(p = 1, 2, 3\) (Einstein, Gauss-Bonnet,
  cubic), as an index and as a number. Status: STANDARD.
\item
  \(\alpha_{(p)}\): The Lovelock coupling of order \(p\),
  \(\alpha_{(p)} = w_{(p)}H^{2-2p}\) (\(\alpha_{(1)} = 1\)). Unit:
  GeV\(^{2-2p}\). Status: AUTHOR.
\item
  \(\mathcal{E}_{(p)}{}^M{}_N\): The Lovelock tensor of order \(p\): the
  variation of \(\sqrt{\lvert g\rvert}\,\mathcal{L}_{(p)}\) with respect
  to the metric (\(\mathcal{E}_{(1)}\) the Einstein tensor,
  \(\mathcal{E}_{(2)}\) the Gauss-Bonnet tensor); with frame indices its
  frame components. Unit: 1/length\(^{2p}\). Status: STANDARD.
\item
  \(M\): A coordinate index of eight dimensions (also with a numbered
  copy \(M_1, M_2, \dots\)): it runs over \(x_0, x_1, \dots, x_7\); a
  repeated index is summed. Status: STANDARD.
\item
  \(N\): A coordinate index of eight dimensions (also
  \(N_1, N_2, \dots\)), running over \(x_0, \dots, x_7\); a repeated
  index is summed. Status: STANDARD.
\item
  \(\Lambda\): The eight-dimensional cosmological constant,
  \(\Lambda = -10H^2\) (the author writes the pure number
  \(\Lambda_{\mathrm{N}} = 10\) in
  \(0 = -\Lambda_{\mathrm{N}} + \sum_p H^{-2p}w_{(p)}\mathcal{E}_{(p)}\)).
  Unit: GeV\(^2\). Value: \(-10H^2\). Status: AUTHOR.
\item
  \(\delta\): The Kronecker delta: with two indices \(\delta^M_N = 1\)
  if \(M = N\) and 0 otherwise (likewise \(\delta_{X\mathrm{L}} = 1\)
  for \(X = \mathrm{L}\)); with upper and lower index strings the
  generalised Kronecker delta
  \(\delta^{M_1\cdots M_n}_{N_1\cdots N_n}\), the determinant of the
  \(n \times n\) matrix of the Kronecker deltas \(\delta^{M_a}_{N_b}\)
  (\(+1\) or \(-1\) if the upper indices are an even or odd permutation
  of the lower ones, 0 otherwise). Status: STANDARD.
\item
  \(\kappa\): The eight-dimensional gravitational coupling,
  \(\kappa = 8\pi G_8\). Unit: GeV\(^{-6}\). Status: AUTHOR.
\item
  \(\mathcal{T}_\Psi\): The energy-momentum tensor of dirac16complex,
  \(\mathcal{T}_{MN} = -(2/\sqrt{\lvert g\rvert})\,\delta S_\Psi/\delta g^{MN}\).
  Unit: GeV\(^8\). Status: STANDARD.
\item
  \(\mathcal{T}_{\mathrm{stab}}{}^a{}_b\): The energy-momentum tensor of
  the stabiliser of the bulk (frame components). Unit: GeV\(^8\).
  Status: ASSUMED.
\item
  \(\sigma\): The proper coordinate of the hidden direction:
  \(\mathrm{d}\sigma = \lvert\cot z\rvert\,\mathrm{d}x_0\),
  \(\sigma = \mathrm{sgn}(\pi/2 - z)\ln(\sin z)/(6H)\) on one period
  \(0 < z < \pi\); \(\sigma = 0\) is the fold \(z = \pi/2\),
  \(\sigma \to \mp\infty\) the tips \(z \to 0, \pi\); one period of
  \(x_0\) is the whole line \(-\infty < \sigma < \infty\). Unit: length.
  Status: DERIVED.
\end{itemize}

The variation of \(S_{\mathrm{tot}}\) with respect to the metric, away
from the fold. The left-hand side is the Lovelock tensor of the author's
theory; the right-hand side is \(\kappa\) times the energy-momentum
tensor of the fields of the bulk: dirac16complex (FE-6) and the
stabiliser (IN-7). With the energy-momentum tensor defined by
\(\mathcal{T}_{MN} = -(2/\sqrt{\lvert g\rvert})\,\delta S_{\mathrm{matter}}/\delta g^{MN}\),
the time-time component
\(\mathcal{T}_{x_4x_4} = -\mathcal{T}^{x_4}{}_{x_4}\) is the energy
density seen by an observer at rest. Status: AUTHOR (the theory),
STANDARD (the variation).

\[
\begin{aligned}
\mathcal{E}_{(p)}{}^M{}_N &= -\frac{1}{2^{p+1}}\,\delta^{MM_1N_1\cdots M_pN_p}_{NP_1Q_1\cdots P_pQ_p}\,
 R^{P_1Q_1}{}_{M_1N_1}\cdots R^{P_pQ_p}{}_{M_pN_p},\\
\mathcal{E}_{(1)}{}^M{}_N &= R^M{}_N - \frac12\,\delta^M_N\,R,\\
\mathcal{E}_{(2)}{}^M{}_N &= 2\big(R\,R^M{}_N - 2R^M{}_PR^P{}_N - 2R^{MP}{}_{NQ}R^Q{}_P + R^{MPQS}R_{NPQS}\big)
 - \frac12\,\delta^M_N\,\mathcal{L}_{(2)}
\end{aligned} \tag{FE-2}
\] \textbf{Symbols in (FE-2).}

\begin{itemize}
\tightlist
\item
  \(\mathcal{E}_{(p)}{}^M{}_N\): The Lovelock tensor of order \(p\): the
  variation of \(\sqrt{\lvert g\rvert}\,\mathcal{L}_{(p)}\) with respect
  to the metric (\(\mathcal{E}_{(1)}\) the Einstein tensor,
  \(\mathcal{E}_{(2)}\) the Gauss-Bonnet tensor); with frame indices its
  frame components. Unit: 1/length\(^{2p}\). Status: STANDARD.
\item
  \(p\): The Lovelock order \(p = 1, 2, 3\) (Einstein, Gauss-Bonnet,
  cubic), as an index and as a number. Status: STANDARD.
\item
  \(M\): A coordinate index of eight dimensions (also with a numbered
  copy \(M_1, M_2, \dots\)): it runs over \(x_0, x_1, \dots, x_7\); a
  repeated index is summed. Status: STANDARD.
\item
  \(N\): A coordinate index of eight dimensions (also
  \(N_1, N_2, \dots\)), running over \(x_0, \dots, x_7\); a repeated
  index is summed. Status: STANDARD.
\item
  \(\delta\): The Kronecker delta: with two indices \(\delta^M_N = 1\)
  if \(M = N\) and 0 otherwise (likewise \(\delta_{X\mathrm{L}} = 1\)
  for \(X = \mathrm{L}\)); with upper and lower index strings the
  generalised Kronecker delta
  \(\delta^{M_1\cdots M_n}_{N_1\cdots N_n}\), the determinant of the
  \(n \times n\) matrix of the Kronecker deltas \(\delta^{M_a}_{N_b}\)
  (\(+1\) or \(-1\) if the upper indices are an even or odd permutation
  of the lower ones, 0 otherwise). Status: STANDARD.
\item
  \(P\): A coordinate index of eight dimensions (also
  \(P_1, P_2, \dots\)), running over \(x_0, \dots, x_7\); a repeated
  index is summed. Status: STANDARD.
\item
  \(Q\): A coordinate index of eight dimensions (also
  \(Q_1, Q_2, \dots\)), running over \(x_0, \dots, x_7\); a repeated
  index is summed. Status: STANDARD.
\item
  \(R^{MN}{}_{PQ},\ R_{MN},\ R\): The curvature of the metric
  \(g_{MN}\): with four indices the Riemann tensor (convention
  \([\nabla_P, \nabla_Q]V^M = R^M{}_{NPQ}V^N\)), with two indices the
  Ricci tensor \(R_{MN} = R^P{}_{MPN}\), without indices the scalar
  curvature \(R = g^{MN}R_{MN}\); with frame indices
  \(R^{(a)(b)}{}_{(c)(d)}\) its frame components,
  \(R^{(a)(b)}{}_{(a)(b)}\) (no sum) the sectional curvature of the
  frame plane \((a, b)\). Unit: 1/length\(^2\). Status: STANDARD.
\item
  \(S\): A coordinate index of eight dimensions, running over
  \(x_0, \dots, x_7\); a repeated index is summed. Status: STANDARD.
\item
  \(\mathcal{L}_{(p)}\): The Lovelock density of order \(p\) (a scalar
  built from \(p\) Riemann tensors): \(\mathcal{L}_{(1)} = R\)
  (Einstein-Hilbert), \(\mathcal{L}_{(2)}\) the Gauss-Bonnet density,
  \(\mathcal{L}_{(3)}\) the cubic Lovelock density. Unit:
  1/length\(^{2p}\). Status: STANDARD.
\end{itemize}

The Lovelock tensor of order \(p\): the variation of
\(\sqrt{\lvert g\rvert}\,\mathcal{L}_{(p)}\) with respect to the metric.
For \(p = 1\) it is the Einstein tensor, for \(p = 2\) the Gauss-Bonnet
tensor (written out in the third line); the cubic tensor
\(\mathcal{E}_{(3)}\) is the first line with \(p = 3\). UM-3 evaluates
the first line exactly by enumerating the non-zero frame components of
the curvature (\texttt{um3\_lovelock.py}); the second and third lines
are independent checks of it (DERIVED, \texttt{um3\_geometry.py}, G3:
the Einstein tensor from the coordinate Ricci tensor and the explicit
Gauss-Bonnet tensor at random points agree with the first line). Status:
STANDARD.

\[
\mathcal{E}_{(p)}{}^M{}_M = -\frac{8 - 2p}{2}\,\mathcal{L}_{(p)},\qquad \nabla_M\,\mathcal{E}_{(p)}{}^M{}_N = 0 \tag{FE-3}
\] \textbf{Symbols in (FE-3).}

\begin{itemize}
\tightlist
\item
  \(\mathcal{E}_{(p)}{}^M{}_N\): The Lovelock tensor of order \(p\): the
  variation of \(\sqrt{\lvert g\rvert}\,\mathcal{L}_{(p)}\) with respect
  to the metric (\(\mathcal{E}_{(1)}\) the Einstein tensor,
  \(\mathcal{E}_{(2)}\) the Gauss-Bonnet tensor); with frame indices its
  frame components. Unit: 1/length\(^{2p}\). Status: STANDARD.
\item
  \(p\): The Lovelock order \(p = 1, 2, 3\) (Einstein, Gauss-Bonnet,
  cubic), as an index and as a number. Status: STANDARD.
\item
  \(M\): A coordinate index of eight dimensions (also with a numbered
  copy \(M_1, M_2, \dots\)): it runs over \(x_0, x_1, \dots, x_7\); a
  repeated index is summed. Status: STANDARD.
\item
  \(N\): A coordinate index of eight dimensions (also
  \(N_1, N_2, \dots\)), running over \(x_0, \dots, x_7\); a repeated
  index is summed. Status: STANDARD.
\item
  \(\mathcal{L}_{(p)}\): The Lovelock density of order \(p\) (a scalar
  built from \(p\) Riemann tensors): \(\mathcal{L}_{(1)} = R\)
  (Einstein-Hilbert), \(\mathcal{L}_{(2)}\) the Gauss-Bonnet density,
  \(\mathcal{L}_{(3)}\) the cubic Lovelock density. Unit:
  1/length\(^{2p}\). Status: STANDARD.
\item
  \(\nabla\): The covariant derivative of a tensor with the Levi-Civita
  connection of the metric \(g_{MN}\) (Christoffel symbols
  \(\Gamma^P{}_{MN}\)). Status: STANDARD.
\end{itemize}

Two identities of every Lovelock tensor in eight dimensions: its trace
is a multiple of the Lovelock density, and its covariant divergence
vanishes (the contracted Bianchi identity). The second makes the
energy-momentum tensor of the sources conserved (FE-11). Status:
STANDARD (DERIVED for the author's metric with general \(a_4(x_4)\),
\(b_4(x_4)\): \texttt{um3\_geometry.py}, G3).

\hypertarget{the-junction-conditions-at-the-fold}{%
\subsection{The junction conditions at the
fold}\label{the-junction-conditions-at-the-fold}}

\[
\begin{aligned}
&\kappa\,\mathcal{S}^a{}_b = \lim_{\varepsilon\to0}\int_{-\varepsilon}^{\varepsilon}\mathrm{d}\sigma\,
 \sum_{p=1}^{3}\alpha_{(p)}\,\mathcal{E}_{(p)}{}^a{}_b
 = \mathfrak{F}^a{}_b\big\vert_{\sigma=0^+} - \mathfrak{F}^a{}_b\big\vert_{\sigma=0^-},\\
&\frac{\partial\mathfrak{F}^a{}_b}{\partial(\partial_\sigma\mathfrak{f})}
 = \frac{\partial}{\partial(\partial_\sigma^2\mathfrak{f})}\sum_{p=1}^{3}\alpha_{(p)}\,\mathcal{E}_{(p)}{}^a{}_b,
 \qquad \mathfrak{f} \in \{\mathfrak{a}, \mathfrak{b}, \mathfrak{n}\}
\end{aligned} \tag{FE-4}
\] \textbf{Symbols in (FE-4).}

\begin{itemize}
\tightlist
\item
  \(\kappa\): The eight-dimensional gravitational coupling,
  \(\kappa = 8\pi G_8\). Unit: GeV\(^{-6}\). Status: AUTHOR.
\item
  \(\mathcal{S}^a{}_b\): The thin source of the fold: its
  energy-momentum tensor per unit proper length across the fold
  (eight-dimensional normalisation, frame components along the fold);
  \(\kappa\mathcal{S}\) is the jump of the potential \(\mathfrak{F}\).
  Unit: GeV\(^7\). Status: DERIVED.
\item
  \(a\): A frame (orthonormal) index running over \(0, 1, \dots, 7\)
  (frame direction \((0)\) along \(\sigma\), \((1), (2), (3)\) along
  ordinary space, \((4)\) along the time, \((5), (6), (7)\) along the
  extra times); raised and lowered with \(\eta_{ab}\); on the fold
  (\(a, b \in \{1, \dots, 7\}\)) the frame of the fold; a repeated index
  is summed. Status: STANDARD.
\item
  \(b\): A frame index (as \(a\)). Status: STANDARD.
\item
  \(\varepsilon\): The half-thickness of a layer around the fold, taken
  to zero. Unit: length. Status: STANDARD.
\item
  \(\mathrm{d}\): The differential: \(\mathrm{d}X\) is an infinitesimal
  change of \(X\), \(\mathrm{d}X/\mathrm{d}Y\) the derivative of \(X\)
  with respect to \(Y\), \(\mathrm{d}^nX\) the volume element of \(n\)
  variables. Status: STANDARD.
\item
  \(\sigma\): The proper coordinate of the hidden direction:
  \(\mathrm{d}\sigma = \lvert\cot z\rvert\,\mathrm{d}x_0\),
  \(\sigma = \mathrm{sgn}(\pi/2 - z)\ln(\sin z)/(6H)\) on one period
  \(0 < z < \pi\); \(\sigma = 0\) is the fold \(z = \pi/2\),
  \(\sigma \to \mp\infty\) the tips \(z \to 0, \pi\); one period of
  \(x_0\) is the whole line \(-\infty < \sigma < \infty\). Unit: length.
  Status: DERIVED.
\item
  \(p\): The Lovelock order \(p = 1, 2, 3\) (Einstein, Gauss-Bonnet,
  cubic), as an index and as a number. Status: STANDARD.
\item
  \(\alpha_{(p)}\): The Lovelock coupling of order \(p\),
  \(\alpha_{(p)} = w_{(p)}H^{2-2p}\) (\(\alpha_{(1)} = 1\)). Unit:
  GeV\(^{2-2p}\). Status: AUTHOR.
\item
  \(\mathcal{E}_{(p)}{}^M{}_N\): The Lovelock tensor of order \(p\): the
  variation of \(\sqrt{\lvert g\rvert}\,\mathcal{L}_{(p)}\) with respect
  to the metric (\(\mathcal{E}_{(1)}\) the Einstein tensor,
  \(\mathcal{E}_{(2)}\) the Gauss-Bonnet tensor); with frame indices its
  frame components. Unit: 1/length\(^{2p}\). Status: STANDARD.
\item
  \(\mathfrak{F}^a{}_b\): The junction potential: a function of the
  first \(\sigma\)-derivatives of the metric functions whose derivatives
  are the coefficients of the second \(\sigma\)-derivatives in
  \(\sum_p\alpha_{(p)}\mathcal{E}_{(p)}{}^a{}_b\) (FE-4). Unit: GeV.
  Status: DERIVED.
\item
  \(\partial\): A partial derivative: \(\partial_M\) with respect to the
  coordinate \(X^M\), \(\partial_\mu\) with respect to \(x^\mu\),
  \(\partial_\sigma\) with respect to \(\sigma\),
  \(\partial X/\partial Y\) the derivative of \(X\) with respect to
  \(Y\) at fixed other variables. Status: STANDARD.
\item
  \(\sigma\): The label \(\sigma\) in \(\partial_\sigma\): the
  derivative with respect to the proper hidden coordinate \(\sigma\).
  Status: STANDARD.
\item
  \(\mathfrak{f}\): Any one of the three metric functions
  \(\mathfrak{a}, \mathfrak{b}, \mathfrak{n}\) of the general warped
  metric (FE-34). Unit: dimensionless. Status: DERIVED.
\item
  \(\mathfrak{a}\): The logarithm of the scale factor of ordinary space
  in the general warped metric FE-34, a function of \(\sigma\) and
  \(x_4\); for the author's metric
  \(\mathfrak{a} = a_4 - H\lvert\sigma\rvert\); \(\mathfrak{a}_n\) its
  term of order \(n\) in the power series FE-35. Unit: dimensionless.
  Status: DERIVED.
\item
  \(\mathfrak{b}\): The logarithm of the scale factor of the extra times
  in the general warped metric FE-34; for the author's metric
  \(\mathfrak{b} = \ln b_4 - a_4 - H\lvert\sigma\rvert\);
  \(\mathfrak{b}_n\) its term of order \(n\). Unit: dimensionless.
  Status: DERIVED.
\item
  \(\mathfrak{n}\): The logarithm of the lapse in the general warped
  metric FE-34 (\(g_{x_4x_4} = -\mathrm{e}^{2\mathfrak{n}}\)); for the
  author's metric \(\mathfrak{n} = 0\); \(\mathfrak{n}_n\) its term of
  order \(n\). Unit: dimensionless. Status: DERIVED.
\end{itemize}

The field equations at the fold. The author's metric (FE-27) has a kink
at \(\sigma = 0\); in general (FE-34) the metric functions
\(\mathfrak{a}, \mathfrak{b}, \mathfrak{n}\) have first
\(\sigma\)-derivatives that jump at the fold, so their second
derivatives contain \(\delta_{\mathrm{D}}(\sigma)\). Because every
Lovelock tensor is linear in the second \(\sigma\)-derivatives
(quasi-linearity, DERIVED for the general metric: \texttt{um3\_fold.py},
J3), FE-1 extended across the fold reads
\(\sum_p\alpha_{(p)}\mathcal{E}_{(p)} + \Lambda = \kappa(\mathcal{T}_{\mathrm{bulk}} + \mathcal{S}\,\delta_{\mathrm{D}}(\sigma))\):
the fold carries the thin source \(\mathcal{S}^a{}_b\) (energy per unit
area of the fold; \(a, b\) along the fold), equal to the integral of the
Lovelock tensor across the fold. The coefficients of the second
derivatives form an exact differential of the first derivatives
(DERIVED, J4), with the potential \(\mathfrak{F}^a{}_b\); the integral
is therefore the difference of \(\mathfrak{F}\) on the two sides. The
\((\sigma, \sigma)\) and \((\sigma, x_4)\) components contain no second
\(\sigma\)-derivative: they carry no thin part and hold on each side
separately. Status: DERIVED.

\[
\mathcal{S}^a{}_b = \frac{1}{V_*}\Big(\mathcal{T}_{\mathrm{H}}{}^a{}_b + \mathcal{T}_{\mathrm{gauge}}{}^a{}_b + \mathcal{T}_{\mathrm{f}}{}^a{}_b
 + \mathcal{T}_{\mathrm{P}}{}^a{}_b + \mathcal{T}_{\mathrm{ff}}{}^a{}_b + \mathcal{T}_{\mathrm{IG}}{}^a{}_b\Big),\qquad
\mathcal{T}_{\mathrm{IG}}{}^\mu{}_\nu = -\frac{1}{\kappa_4}\,\hat G^\mu{}_\nu,\qquad
\mathcal{T}^r{}_r = \mathcal{L}_{\mathrm{fold}} \tag{FE-5}
\] \textbf{Symbols in (FE-5).}

\begin{itemize}
\tightlist
\item
  \(\mathcal{S}^a{}_b\): The thin source of the fold: its
  energy-momentum tensor per unit proper length across the fold
  (eight-dimensional normalisation, frame components along the fold);
  \(\kappa\mathcal{S}\) is the jump of the potential \(\mathfrak{F}\).
  Unit: GeV\(^7\). Status: DERIVED.
\item
  \(a\): A frame (orthonormal) index running over \(0, 1, \dots, 7\)
  (frame direction \((0)\) along \(\sigma\), \((1), (2), (3)\) along
  ordinary space, \((4)\) along the time, \((5), (6), (7)\) along the
  extra times); raised and lowered with \(\eta_{ab}\); on the fold
  (\(a, b \in \{1, \dots, 7\}\)) the frame of the fold; a repeated index
  is summed. Status: STANDARD.
\item
  \(b\): A frame index (as \(a\)). Status: STANDARD.
\item
  \(V_*\): The physical volume of the extra times \(x_5, x_6, x_7\) on
  the fold in the late universe,
  \(V_* = V_{\mathrm{t}}(\mathrm{e}^{-a_4}b_4)^3\) at a time when
  \(\mathrm{e}^{-a_4}b_4\) is constant; it converts the
  eight-dimensional normalisation into the four-dimensional one. Unit:
  length\(^3\) (GeV\(^{-3}\)). Status: ASSUMED.
\item
  \(\mathcal{T}_{\mathrm{H}}\): The energy-momentum tensor of the Higgs
  doublet (four-dimensional normalisation). Unit: GeV\(^4\). Status:
  STANDARD.
\item
  \(\mathcal{T}_{\mathrm{gauge}}\): The energy-momentum tensor of the
  gauge fields. Unit: GeV\(^4\). Status: STANDARD.
\item
  \(\mathcal{T}_{\mathrm{f}}\): The energy-momentum tensor of the quarks
  and leptons. Unit: GeV\(^4\). Status: STANDARD.
\item
  \(\mathcal{T}_{\mathrm{P}}\): The energy-momentum tensor of the
  portals of dirac16complex. Unit: GeV\(^4\). Status: DERIVED.
\item
  \(\mathcal{T}_{\mathrm{ff}}{}^a{}_b\): The energy-momentum tensor of
  the fold fluid (frame components of the fold, four-dimensional
  normalisation). Unit: GeV\(^4\). Status: ASSUMED.
\item
  \(\mathcal{T}_{\mathrm{IG}}\): The energy-momentum tensor of the
  induced gravity of the fold,
  \(\mathcal{T}_{\mathrm{IG}}{}^\mu{}_\nu = -\hat G^\mu{}_\nu/\kappa_4\).
  Unit: GeV\(^4\). Status: DERIVED.
\item
  \(\mu\): A four-dimensional coordinate index running over the
  observable coordinates \(x_1, x_2, x_3, x_4\) of the fold (signature
  \(+,+,+,-\)); a repeated index is summed. Status: STANDARD.
\item
  \(\nu\): A four-dimensional coordinate index (as \(\mu\)). Status:
  STANDARD.
\item
  \(\kappa_4\): The coupling of the induced gravity of the fold,
  \(\kappa_4 = 8\pi G_{\mathrm{ind}}\); four-dimensional gravity on the
  fold has \(8\pi G_{\mathrm{N}} = \kappa_4/(1 + 5r/6)\) (FE-59). Unit:
  GeV\(^{-2}\). Value:
  \(\kappa_4 \simeq 8\pi G_{\mathrm{N}} = 8\pi/M_{\mathrm{Pl}}^2\).
  Status: ASSUMED.
\item
  \(\hat G^\mu{}_\nu\): The Einstein tensor of the four-dimensional
  metric \(\hat h_{\mu\nu}\) of the fold. Unit: 1/length\(^2\). Status:
  STANDARD.
\item
  \(r\): An extra-time direction of the fold: \(r \in \{5, 6, 7\}\)
  (\(x_5, x_6, x_7\)); not summed (the three components are equal).
  Status: STANDARD.
\item
  \(\mathcal{T}\): The energy-momentum tensor of any one term of the
  fold (in \(\mathcal{T}^r{}_r = \mathcal{L}_{\mathrm{fold}}\)). Unit:
  GeV\(^4\). Status: STANDARD.
\item
  \(\mathcal{L}_{\mathrm{fold}}\): The Lagrangian of any one term of the
  fold (four-dimensional normalisation). Unit: GeV\(^4\). Status:
  STANDARD.
\end{itemize}

The thin source of the fold is the energy-momentum tensor of everything
that lives on the fold, divided by \(V_*\) (IN-8): the Higgs doublet
(FE-7), the gauge fields (FE-8), the quarks and leptons (FE-9), the
portals (FE-10), the fold fluid (IN-10) and the induced gravity, whose
energy-momentum tensor is \(-1/\kappa_4\) times the Einstein tensor
\(\hat G^\mu{}_\nu\) of the metric of the fold. Because the fields of
the fold do not depend on the extra times and have no components along
them, the extra times enter only through \(\sqrt{\lvert g\rvert}\): for
every term of the fold the component along an extra time is the
Lagrangian of the term,
\(\mathcal{T}^r{}_r = \mathcal{L}_{\mathrm{fold}}\) (DERIVED:
\(\delta\sqrt{\lvert g\rvert} = -\frac12\sqrt{\lvert g\rvert}g_{rr}\delta g^{rr}\)).
Combining FE-4 and FE-5 gives the Einstein equations of the fold
(FE-59). Status: DERIVED.

\hypertarget{the-energy-momentum-tensors}{%
\subsection{The energy-momentum
tensors}\label{the-energy-momentum-tensors}}

\[
{\mathcal{T}_\Psi}_{MN} = -\frac14\Big(\bar\Psi\gamma_MD_N\Psi + \bar\Psi\gamma_ND_M\Psi - (D_M\bar\Psi)\gamma_N\Psi
 - (D_N\bar\Psi)\gamma_M\Psi\Big) + g_{MN}\,\mathcal{L}_\Psi \tag{FE-6}
\] \textbf{Symbols in (FE-6).}

\begin{itemize}
\tightlist
\item
  \(\mathcal{T}_\Psi\): The energy-momentum tensor of dirac16complex,
  \(\mathcal{T}_{MN} = -(2/\sqrt{\lvert g\rvert})\,\delta S_\Psi/\delta g^{MN}\).
  Unit: GeV\(^8\). Status: STANDARD.
\item
  \(M\): A coordinate index of eight dimensions (also with a numbered
  copy \(M_1, M_2, \dots\)): it runs over \(x_0, x_1, \dots, x_7\); a
  repeated index is summed. Status: STANDARD.
\item
  \(N\): A coordinate index of eight dimensions (also
  \(N_1, N_2, \dots\)), running over \(x_0, \dots, x_7\); a repeated
  index is summed. Status: STANDARD.
\item
  \(\Psi,\ \bar\Psi\): dirac16complex: the author's 16-component complex
  spinor field of eight dimensions;
  \(\bar\Psi = \Psi^\dagger\sigma_{16}\) is its Dirac adjoint; in a term
  of the fold, \(\Psi\) is its value at \(\sigma = 0\). Unit:
  GeV\(^{7/2}\). Status: AUTHOR.
\item
  \(\gamma^M,\ \gamma^{(a)}\): The Dirac matrices of dirac16complex
  along the coordinates, \(\gamma^M = \mathbf{e}_a{}^MT^a\),
  \(\gamma_M = g_{MN}\gamma^N\); \(\gamma^{(a)} = T^a\) along the frame
  direction \((a)\) (\(\gamma^{(0)}\) along \(\sigma\), \(\gamma^{(4)}\)
  along the time). Unit: dimensionless. Status: AUTHOR.
\item
  \(D_M,\ D_\mu\): The covariant derivative: on dirac16complex
  \(D_M = \partial_M + \Omega_M\) (with
  \(D_M\bar\Psi = \partial_M\bar\Psi - \bar\Psi\Omega_M\)); on the
  fields of the fold the gauge and spin covariant derivative \(D_\mu\)
  of IN-13; \(D^\mu = \hat h^{\mu\nu}D_\nu\); on an adjoint field
  \((D_\mu X)^A = \partial_\mu X^A + g_3f^{ABC}G^B_\mu X^C\) (likewise
  for \(SU(2)\) with \(g_2\epsilon^{IJK}\)). Unit: 1/length. Status:
  STANDARD.
\item
  \(g_{MN}\): The metric tensor of the eight-dimensional spacetime
  (signature \(+,+,+,+,-,-,-,-\) in the order \(x_0, \dots, x_7\)):
  \(g_{MN}\,\mathrm{d}X^M\mathrm{d}X^N\) is the squared interval,
  \(g^{MN}\) its inverse, \(\lvert g\rvert\) the absolute value of its
  determinant (\(\sqrt{\lvert g\rvert}\,\mathrm{d}^8X\) is the proper
  volume element). Unit: dimensionless. Status: AUTHOR.
\item
  \(\mathcal{L}_\Psi\): The Lagrangian of dirac16complex, the square
  bracket of IN-4. Unit: GeV\(^8\). Status: RECORD.
\end{itemize}

The symmetric energy-momentum tensor of dirac16complex, from the
variation of the vielbein in IN-4 (\(\mathcal{L}_\Psi\) is the
Lagrangian in the square bracket of IN-4; the variation of the spin
connection adds only the divergence of the spin current, which drops out
of the symmetric part). On solutions of FE-12,
\(\mathcal{L}_\Psi = \frac{\lambda_\Psi}{2}(\bar\Psi\Psi)^2\). For a
stationary state \(\Psi \propto \mathrm{e}^{-\mathrm{i}\omega x_4}\) of
the free field the energy density is
\(\mathcal{T}_{\Psi\,x_4x_4} = -\omega\,\mathrm{i}\bar\Psi\gamma^{x_4}\Psi\):
its sign is the product of the sign of the frequency and of the Krein
density (FE-43). Status: STANDARD (DERIVED for the energy density).

\[
{\mathcal{T}_{\mathrm{H}}}_{\mu\nu} = (D_\mu\Phi)^\dagger D_\nu\Phi + (D_\nu\Phi)^\dagger D_\mu\Phi + \hat h_{\mu\nu}\,\mathcal{L}_{\mathrm{H}} \tag{FE-7}
\] \textbf{Symbols in (FE-7).}

\begin{itemize}
\tightlist
\item
  \(\mathcal{T}_{\mathrm{H}}\): The energy-momentum tensor of the Higgs
  doublet (four-dimensional normalisation). Unit: GeV\(^4\). Status:
  STANDARD.
\item
  \(\mu\): A four-dimensional coordinate index running over the
  observable coordinates \(x_1, x_2, x_3, x_4\) of the fold (signature
  \(+,+,+,-\)); a repeated index is summed. Status: STANDARD.
\item
  \(\nu\): A four-dimensional coordinate index (as \(\mu\)). Status:
  STANDARD.
\item
  \(D_M,\ D_\mu\): The covariant derivative: on dirac16complex
  \(D_M = \partial_M + \Omega_M\) (with
  \(D_M\bar\Psi = \partial_M\bar\Psi - \bar\Psi\Omega_M\)); on the
  fields of the fold the gauge and spin covariant derivative \(D_\mu\)
  of IN-13; \(D^\mu = \hat h^{\mu\nu}D_\nu\); on an adjoint field
  \((D_\mu X)^A = \partial_\mu X^A + g_3f^{ABC}G^B_\mu X^C\) (likewise
  for \(SU(2)\) with \(g_2\epsilon^{IJK}\)). Unit: 1/length. Status:
  STANDARD.
\item
  \(\Phi\): The Higgs doublet (two complex components, hypercharge
  \(1/2\)), a field of the fold. Unit: GeV. Status: AUTHOR.
\item
  \({}^\dagger\): Hermitian conjugation: the complex conjugate of the
  transpose of a matrix or spinor (for field operators the adjoint
  operator). Status: STANDARD.
\item
  \(\hat h_{\mu\nu}\): The four-dimensional metric of the fold,
  \(\hat h_{\mu\nu} = \mathrm{diag}(\mathrm{e}^{2a_4}, \mathrm{e}^{2a_4}, \mathrm{e}^{2a_4}, -1)\)
  in the coordinates \(x_1, \dots, x_4\) (the induced metric at
  \(\sigma = 0\) without the extra times); \(\hat h^{\mu\nu}\) its
  inverse, \(\hat h\) its determinant,
  \(\sqrt{-\hat h} = \mathrm{e}^{3a_4}\); \(\hat h^{(1)}_{\mu\nu}\) its
  first-order cosmological perturbation. Unit: dimensionless. Status:
  DERIVED.
\item
  \(\mathcal{L}_{\mathrm{H}}\): The Lagrangian of the Higgs doublet on
  the fold. Unit: GeV\(^4\). Status: AUTHOR.
\end{itemize}

The energy-momentum tensor of the Higgs doublet (four-dimensional
normalisation): its time-time component is
\(\lvert D_4\Phi\rvert^2 + \hat h^{ij}(D_i\Phi)^\dagger D_j\Phi + V \ge V\).
Status: STANDARD.

\[
{\mathcal{T}_{\mathrm{gauge}}}_{\mu\nu} = G^A_{\mu\rho}G^A{}_\nu{}^\rho + W^I_{\mu\rho}W^I{}_\nu{}^\rho + B_{\mu\rho}B_\nu{}^\rho
 + \hat h_{\mu\nu}\,\mathcal{L}_{\mathrm{gauge}} \tag{FE-8}
\] \textbf{Symbols in (FE-8).}

\begin{itemize}
\tightlist
\item
  \(\mathcal{T}_{\mathrm{gauge}}\): The energy-momentum tensor of the
  gauge fields. Unit: GeV\(^4\). Status: STANDARD.
\item
  \(\mu\): A four-dimensional coordinate index running over the
  observable coordinates \(x_1, x_2, x_3, x_4\) of the fold (signature
  \(+,+,+,-\)); a repeated index is summed. Status: STANDARD.
\item
  \(\nu\): A four-dimensional coordinate index (as \(\mu\)). Status:
  STANDARD.
\item
  \(G^A_\mu,\ G^A_{\mu\nu}\): The gluon field of \(SU(3)_C\) on the fold
  (\(A = 1, \dots, 8\)) and its field strength. Unit: GeV (field),
  GeV\(^2\) (field strength). Status: AUTHOR.
\item
  \(A\): A colour (adjoint SU(3)) index running over \(1, \dots, 8\); a
  repeated index is summed. Status: STANDARD.
\item
  \(\rho\): A four-dimensional coordinate index (as \(\mu\)). Status:
  STANDARD.
\item
  \(W^I_\mu,\ W^I_{\mu\nu}\): The gauge field of \(SU(2)_L\) on the
  fold, \(W^1_\mu, W^2_\mu, W^3_\mu\), and its field strength. Unit: GeV
  (field), GeV\(^2\) (field strength). Status: AUTHOR.
\item
  \(I\): A weak-isospin (adjoint SU(2)) index running over \(1, 2, 3\);
  a repeated index is summed. Status: STANDARD.
\item
  \(B_\mu,\ B_{\mu\nu}\): The hypercharge gauge field of \(U(1)_Y\) on
  the fold and its field strength. Unit: GeV (field), GeV\(^2\) (field
  strength). Status: AUTHOR.
\item
  \(\hat h_{\mu\nu}\): The four-dimensional metric of the fold,
  \(\hat h_{\mu\nu} = \mathrm{diag}(\mathrm{e}^{2a_4}, \mathrm{e}^{2a_4}, \mathrm{e}^{2a_4}, -1)\)
  in the coordinates \(x_1, \dots, x_4\) (the induced metric at
  \(\sigma = 0\) without the extra times); \(\hat h^{\mu\nu}\) its
  inverse, \(\hat h\) its determinant,
  \(\sqrt{-\hat h} = \mathrm{e}^{3a_4}\); \(\hat h^{(1)}_{\mu\nu}\) its
  first-order cosmological perturbation. Unit: dimensionless. Status:
  DERIVED.
\item
  \(\mathcal{L}_{\mathrm{gauge}}\): The Lagrangian of the gauge fields
  of \(SU(3)_C\times SU(2)_L\times U(1)_Y\) on the fold. Unit:
  GeV\(^4\). Status: STANDARD.
\end{itemize}

The energy-momentum tensor of the gauge fields; for each field its
energy density is \((\vec E^2 + \vec B^2)/2\) in terms of the electric
and magnetic components. Status: STANDARD.

\[
{\mathcal{T}_{\mathrm{f}}}_{\mu\nu} = -\frac{\mathrm{i}}{4}\sum_\psi\Big(\bar\psi\hat\gamma_\mu D_\nu\psi + \bar\psi\hat\gamma_\nu D_\mu\psi
 - (D_\mu\bar\psi)\hat\gamma_\nu\psi - (D_\nu\bar\psi)\hat\gamma_\mu\psi\Big) + \hat h_{\mu\nu}\big(\mathcal{L}_{\mathrm{f}} + \mathcal{L}_{\mathrm{Y}}\big) \tag{FE-9}
\] \textbf{Symbols in (FE-9).}

\begin{itemize}
\tightlist
\item
  \(\mathcal{T}_{\mathrm{f}}\): The energy-momentum tensor of the quarks
  and leptons. Unit: GeV\(^4\). Status: STANDARD.
\item
  \(\mu\): A four-dimensional coordinate index running over the
  observable coordinates \(x_1, x_2, x_3, x_4\) of the fold (signature
  \(+,+,+,-\)); a repeated index is summed. Status: STANDARD.
\item
  \(\nu\): A four-dimensional coordinate index (as \(\mu\)). Status:
  STANDARD.
\item
  \(\mathrm{i}\): The imaginary unit, \(\mathrm{i}^2 = -1\). Unit:
  dimensionless. Status: STANDARD.
\item
  \(\psi\): The label of a Standard-Model fermion field in a sum:
  \(\psi\) runs over
  \(Q, u_{\mathrm{R}}, d_{\mathrm{R}}, \ell, e_{\mathrm{R}}, \nu_{\mathrm{R}}\).
  Status: STANDARD.
\item
  \(\psi,\ \bar\psi\): A Standard-Model fermion field (a
  four-dimensional Weyl field of the fold with three generations,
  written as a four-component spinor of fixed chirality);
  \(\bar\psi = \psi^\dagger\beta\); \(\psi_f\) the fields of the kind
  \(f\). Unit: GeV\(^{3/2}\). Status: STANDARD.
\item
  \(\hat\gamma^\mu\): The four-dimensional complex \(4 \times 4\) Dirac
  matrices of the fold,
  \(\hat\gamma^\mu = \hat{\mathbf{e}}_a{}^\mu\hat\gamma^{(a)}\),
  \(\{\hat\gamma^\mu, \hat\gamma^\nu\} = -2\hat h^{\mu\nu}\mathbf{1}\);
  \(\hat\gamma_\mu = \hat h_{\mu\nu}\hat\gamma^\nu\). Unit:
  dimensionless. Status: STANDARD.
\item
  \(D_M,\ D_\mu\): The covariant derivative: on dirac16complex
  \(D_M = \partial_M + \Omega_M\) (with
  \(D_M\bar\Psi = \partial_M\bar\Psi - \bar\Psi\Omega_M\)); on the
  fields of the fold the gauge and spin covariant derivative \(D_\mu\)
  of IN-13; \(D^\mu = \hat h^{\mu\nu}D_\nu\); on an adjoint field
  \((D_\mu X)^A = \partial_\mu X^A + g_3f^{ABC}G^B_\mu X^C\) (likewise
  for \(SU(2)\) with \(g_2\epsilon^{IJK}\)). Unit: 1/length. Status:
  STANDARD.
\item
  \(\hat h_{\mu\nu}\): The four-dimensional metric of the fold,
  \(\hat h_{\mu\nu} = \mathrm{diag}(\mathrm{e}^{2a_4}, \mathrm{e}^{2a_4}, \mathrm{e}^{2a_4}, -1)\)
  in the coordinates \(x_1, \dots, x_4\) (the induced metric at
  \(\sigma = 0\) without the extra times); \(\hat h^{\mu\nu}\) its
  inverse, \(\hat h\) its determinant,
  \(\sqrt{-\hat h} = \mathrm{e}^{3a_4}\); \(\hat h^{(1)}_{\mu\nu}\) its
  first-order cosmological perturbation. Unit: dimensionless. Status:
  DERIVED.
\item
  \(\mathcal{L}_{\mathrm{f}}\): The kinetic Lagrangian of the quarks and
  leptons on the fold. Unit: GeV\(^4\). Status: STANDARD.
\item
  \(\mathcal{L}_{\mathrm{Y}}\): The Yukawa Lagrangian of the quarks and
  leptons with the Higgs doublet. Unit: GeV\(^4\). Status: STANDARD.
\end{itemize}

The energy-momentum tensor of the quarks and leptons (with the kinetic
term in its symmetric form). For a free particle of energy \(E\) at rest
in the frame of the fold its time-time component is
\(E\,\psi^\dagger\psi > 0\). Status: STANDARD.

\[
\begin{aligned}
{\mathcal{T}_{\mathrm{P}}}_{\mu\nu} = {}&\sum_X\Big[\big(c^{G}_X\,G^A_{\mu\rho}G^A{}_\nu{}^\rho + c^{W}_X\,W^I_{\mu\rho}W^I{}_\nu{}^\rho
 + c^{B}_X\,B_{\mu\rho}B_\nu{}^\rho\big)\,\bar\Psi P_X\Psi\\
&+ \frac{d_X}{2}\big(B_{\mu\rho}\,\mathcal{D}_{X\nu}{}^\rho + B_{\nu\rho}\,\mathcal{D}_{X\mu}{}^\rho\big)
 + \sum_f\frac{c^{f}_X}{2}\big(J_{f\mu}\,j_{X\nu} + J_{f\nu}\,j_{X\mu}\big)\Big] + \hat h_{\mu\nu}\,\mathcal{L}_{\mathrm{P}}
\end{aligned} \tag{FE-10}
\] \textbf{Symbols in (FE-10).}

\begin{itemize}
\tightlist
\item
  \(\mathcal{T}_{\mathrm{P}}\): The energy-momentum tensor of the
  portals of dirac16complex. Unit: GeV\(^4\). Status: DERIVED.
\item
  \(\mu\): A four-dimensional coordinate index running over the
  observable coordinates \(x_1, x_2, x_3, x_4\) of the fold (signature
  \(+,+,+,-\)); a repeated index is summed. Status: STANDARD.
\item
  \(\nu\): A four-dimensional coordinate index (as \(\mu\)). Status:
  STANDARD.
\item
  \(X\): The label of a half of dirac16complex:
  \(X \in \{\mathrm{L}, \mathrm{R}\}\) (left-handed or right-handed
  half); a sum over \(X\) runs over both. Status: STANDARD.
\item
  \(c^{G}_X\): The gluon-kinetic portal coupling of the half \(X\) of
  dirac16complex. Unit: GeV\(^{-7}\). Status: INVENTED.
\item
  \(G^A_\mu,\ G^A_{\mu\nu}\): The gluon field of \(SU(3)_C\) on the fold
  (\(A = 1, \dots, 8\)) and its field strength. Unit: GeV (field),
  GeV\(^2\) (field strength). Status: AUTHOR.
\item
  \(A\): A colour (adjoint SU(3)) index running over \(1, \dots, 8\); a
  repeated index is summed. Status: STANDARD.
\item
  \(\rho\): A four-dimensional coordinate index (as \(\mu\)). Status:
  STANDARD.
\item
  \(c^{W}_X\): The \(W\)-kinetic portal coupling of the half \(X\).
  Unit: GeV\(^{-7}\). Status: INVENTED.
\item
  \(W^I_\mu,\ W^I_{\mu\nu}\): The gauge field of \(SU(2)_L\) on the
  fold, \(W^1_\mu, W^2_\mu, W^3_\mu\), and its field strength. Unit: GeV
  (field), GeV\(^2\) (field strength). Status: AUTHOR.
\item
  \(I\): A weak-isospin (adjoint SU(2)) index running over \(1, 2, 3\);
  a repeated index is summed. Status: STANDARD.
\item
  \(c^{B}_X\): The \(B\)-kinetic portal coupling of the half \(X\).
  Unit: GeV\(^{-7}\). Status: INVENTED.
\item
  \(B_\mu,\ B_{\mu\nu}\): The hypercharge gauge field of \(U(1)_Y\) on
  the fold and its field strength. Unit: GeV (field), GeV\(^2\) (field
  strength). Status: AUTHOR.
\item
  \(\Psi,\ \bar\Psi\): dirac16complex: the author's 16-component complex
  spinor field of eight dimensions;
  \(\bar\Psi = \Psi^\dagger\sigma_{16}\) is its Dirac adjoint; in a term
  of the fold, \(\Psi\) is its value at \(\sigma = 0\). Unit:
  GeV\(^{7/2}\). Status: AUTHOR.
\item
  \(P_X\): The chiral projector of the half \(X\):
  \(P_X = P_{\mathrm{L}}\) for \(X = \mathrm{L}\), \(P_{\mathrm{R}}\)
  for \(X = \mathrm{R}\). Unit: dimensionless. Status: AUTHOR.
\item
  \(d_X\): The dipole-portal coupling of the half \(X\) of
  dirac16complex with the hypercharge field strength. Unit:
  GeV\(^{-5}\). Status: INVENTED.
\item
  \(\mathcal{D}_X^{\mu\nu}\): The dipole (antisymmetric tensor) density
  of the half \(X\) of dirac16complex at the fold,
  \(\mathcal{D}_X^{\mu\nu} = \mathrm{i}\bar\Psi\Sigma^{\mu\nu}P_X\Psi\)
  (real, DERIVED). Unit: GeV\(^7\). Status: DERIVED.
\item
  \(f\): The label of a kind of Standard-Model multiplet:
  \(f \in \{Q, u, d, \ell, e, \nu\}\) (quark doublets, up-quark
  singlets, down-quark singlets, lepton doublets, charged-lepton
  singlets, right-handed neutrinos); a sum over \(f\) runs over the six.
  Status: STANDARD.
\item
  \(c^{f}_X\): The current-portal coupling of the half \(X\) of
  dirac16complex with the fermion-number current of the multiplets of
  kind \(f\) (flavour blind). Unit: GeV\(^{-7}\). Status: INVENTED.
\item
  \(J_f^\mu\): The fermion-number current of the Standard-Model
  multiplets of kind \(f\), \(J_f^\mu = \bar\psi_f\hat\gamma^\mu\psi_f\)
  summed over the three generations and the gauge components; with other
  labels the gauge currents (FE-16 to FE-18). Unit: GeV\(^3\). Status:
  STANDARD.
\item
  \(j_X^\mu\): The convective current of the half \(X\) of
  dirac16complex,
  \(j_X^\mu = \mathrm{i}(\bar\Psi P_XD^\mu\Psi - (D^\mu\bar\Psi)P_X\Psi)\)
  (real and joining the half with itself, DERIVED). Unit: GeV\(^8\).
  Status: DERIVED.
\item
  \(\hat h_{\mu\nu}\): The four-dimensional metric of the fold,
  \(\hat h_{\mu\nu} = \mathrm{diag}(\mathrm{e}^{2a_4}, \mathrm{e}^{2a_4}, \mathrm{e}^{2a_4}, -1)\)
  in the coordinates \(x_1, \dots, x_4\) (the induced metric at
  \(\sigma = 0\) without the extra times); \(\hat h^{\mu\nu}\) its
  inverse, \(\hat h\) its determinant,
  \(\sqrt{-\hat h} = \mathrm{e}^{3a_4}\); \(\hat h^{(1)}_{\mu\nu}\) its
  first-order cosmological perturbation. Unit: dimensionless. Status:
  DERIVED.
\item
  \(\mathcal{L}_{\mathrm{P}}\): The Lagrangian of the portals of
  dirac16complex at the fold. Unit: GeV\(^4\). Status: INVENTED.
\end{itemize}

The energy-momentum tensor of the portals (from the variation of IN-16
with respect to the metric of the fold, which enters through the raised
indices, through \(\hat\gamma^\mu\) in \(J_f^\mu\) and through
\(\sqrt{-\hat h}\)); the Higgs portal enters only through
\(\hat h_{\mu\nu}\mathcal{L}_{\mathrm{P}}\). It is a source of the fold.
Status: DERIVED.

\hypertarget{the-conservation-laws}{%
\subsection{The conservation laws}\label{the-conservation-laws}}

\[
\begin{aligned}
&\nabla_M\big(\mathcal{T}_\Psi{}^M{}_N + \mathcal{T}_{\mathrm{stab}}{}^M{}_N\big) = 0\qquad (\sigma \ne 0),\\
&\hat\nabla_a\,\mathcal{S}^a{}_b = -\Big(\mathcal{T}_{\mathrm{bulk}}{}^{(0)}{}_b\big\vert_{\sigma=0^+}
 - \mathcal{T}_{\mathrm{bulk}}{}^{(0)}{}_b\big\vert_{\sigma=0^-}\Big)
\end{aligned} \tag{FE-11}
\] \textbf{Symbols in (FE-11).}

\begin{itemize}
\tightlist
\item
  \(\nabla\): The covariant derivative of a tensor with the Levi-Civita
  connection of the metric \(g_{MN}\) (Christoffel symbols
  \(\Gamma^P{}_{MN}\)). Status: STANDARD.
\item
  \(\mathcal{T}_\Psi\): The energy-momentum tensor of dirac16complex,
  \(\mathcal{T}_{MN} = -(2/\sqrt{\lvert g\rvert})\,\delta S_\Psi/\delta g^{MN}\).
  Unit: GeV\(^8\). Status: STANDARD.
\item
  \(\mathcal{T}_{\mathrm{stab}}{}^a{}_b\): The energy-momentum tensor of
  the stabiliser of the bulk (frame components). Unit: GeV\(^8\).
  Status: ASSUMED.
\item
  \(M\): A coordinate index of eight dimensions (also with a numbered
  copy \(M_1, M_2, \dots\)): it runs over \(x_0, x_1, \dots, x_7\); a
  repeated index is summed. Status: STANDARD.
\item
  \(N\): A coordinate index of eight dimensions (also
  \(N_1, N_2, \dots\)), running over \(x_0, \dots, x_7\); a repeated
  index is summed. Status: STANDARD.
\item
  \(\sigma\): The proper coordinate of the hidden direction:
  \(\mathrm{d}\sigma = \lvert\cot z\rvert\,\mathrm{d}x_0\),
  \(\sigma = \mathrm{sgn}(\pi/2 - z)\ln(\sin z)/(6H)\) on one period
  \(0 < z < \pi\); \(\sigma = 0\) is the fold \(z = \pi/2\),
  \(\sigma \to \mp\infty\) the tips \(z \to 0, \pi\); one period of
  \(x_0\) is the whole line \(-\infty < \sigma < \infty\). Unit: length.
  Status: DERIVED.
\item
  \(\hat\nabla\): The covariant derivative of the fold: with the
  Levi-Civita connection of the metric of the fold (\(\hat h_{\mu\nu}\)
  for four-dimensional tensors, the induced metric of the
  seven-dimensional fold for \(\hat\nabla_a\mathcal{S}^a{}_b\)). Status:
  STANDARD.
\item
  \(\mathcal{S}^a{}_b\): The thin source of the fold: its
  energy-momentum tensor per unit proper length across the fold
  (eight-dimensional normalisation, frame components along the fold);
  \(\kappa\mathcal{S}\) is the jump of the potential \(\mathfrak{F}\).
  Unit: GeV\(^7\). Status: DERIVED.
\item
  \(a\): A frame (orthonormal) index running over \(0, 1, \dots, 7\)
  (frame direction \((0)\) along \(\sigma\), \((1), (2), (3)\) along
  ordinary space, \((4)\) along the time, \((5), (6), (7)\) along the
  extra times); raised and lowered with \(\eta_{ab}\); on the fold
  (\(a, b \in \{1, \dots, 7\}\)) the frame of the fold; a repeated index
  is summed. Status: STANDARD.
\item
  \(b\): A frame index (as \(a\)). Status: STANDARD.
\item
  \(\mathcal{T}_{\mathrm{bulk}}\): The energy-momentum tensor of all
  sources of the bulk,
  \(\mathcal{T}_\Psi + \mathcal{T}_{\mathrm{stab}}\);
  \(\mathcal{T}_{\mathrm{bulk}}{}^{(0)}{}_b\) is its flux along
  \(\sigma\) (frame components). Unit: GeV\(^8\). Status: STANDARD.
\end{itemize}

The energy and momentum of the bulk are conserved away from the fold
(FE-3 and FE-1). At the fold the thin source exchanges energy and
momentum with the bulk: its covariant divergence along the fold equals
minus the jump of the flux \(\mathcal{T}_{\mathrm{bulk}}{}^{(0)}{}_b\)
of the bulk sources along \(\sigma\) (the frame component \((0), b\)).
For the author's metric with general \(a_4(x_4)\), \(b_4(x_4)\) this is
an identity of the junction formula FE-30 and the bulk equation FE-29
(DERIVED, \texttt{um3\_fold.py}, J6). The fluxes of dirac16complex into
the bulk carry its share of this exchange. Status: DERIVED.

\hypertarget{the-matter-fields}{%
\section{The matter fields}\label{the-matter-fields}}

\hypertarget{dirac16complex}{%
\subsection{dirac16complex}\label{dirac16complex}}

\[
\gamma^MD_M\Psi - m_\Psi\Psi - \lambda_\Psi\big(\bar\Psi\Psi\big)\Psi = 0\qquad (\sigma \ne 0) \tag{FE-12}
\] \textbf{Symbols in (FE-12).}

\begin{itemize}
\tightlist
\item
  \(\gamma^M,\ \gamma^{(a)}\): The Dirac matrices of dirac16complex
  along the coordinates, \(\gamma^M = \mathbf{e}_a{}^MT^a\),
  \(\gamma_M = g_{MN}\gamma^N\); \(\gamma^{(a)} = T^a\) along the frame
  direction \((a)\) (\(\gamma^{(0)}\) along \(\sigma\), \(\gamma^{(4)}\)
  along the time). Unit: dimensionless. Status: AUTHOR.
\item
  \(M\): A coordinate index of eight dimensions (also with a numbered
  copy \(M_1, M_2, \dots\)): it runs over \(x_0, x_1, \dots, x_7\); a
  repeated index is summed. Status: STANDARD.
\item
  \(D_M,\ D_\mu\): The covariant derivative: on dirac16complex
  \(D_M = \partial_M + \Omega_M\) (with
  \(D_M\bar\Psi = \partial_M\bar\Psi - \bar\Psi\Omega_M\)); on the
  fields of the fold the gauge and spin covariant derivative \(D_\mu\)
  of IN-13; \(D^\mu = \hat h^{\mu\nu}D_\nu\); on an adjoint field
  \((D_\mu X)^A = \partial_\mu X^A + g_3f^{ABC}G^B_\mu X^C\) (likewise
  for \(SU(2)\) with \(g_2\epsilon^{IJK}\)). Unit: 1/length. Status:
  STANDARD.
\item
  \(\Psi,\ \bar\Psi\): dirac16complex: the author's 16-component complex
  spinor field of eight dimensions;
  \(\bar\Psi = \Psi^\dagger\sigma_{16}\) is its Dirac adjoint; in a term
  of the fold, \(\Psi\) is its value at \(\sigma = 0\). Unit:
  GeV\(^{7/2}\). Status: AUTHOR.
\item
  \(m_\Psi\): The mass parameter of dirac16complex in the bulk. Unit:
  GeV. Status: RECORD.
\item
  \(\lambda_\Psi\): The coupling of the contact self-interaction of
  dirac16complex. Unit: GeV\(^{-6}\). Status: RECORD.
\item
  \(\sigma\): The proper coordinate of the hidden direction:
  \(\mathrm{d}\sigma = \lvert\cot z\rvert\,\mathrm{d}x_0\),
  \(\sigma = \mathrm{sgn}(\pi/2 - z)\ln(\sin z)/(6H)\) on one period
  \(0 < z < \pi\); \(\sigma = 0\) is the fold \(z = \pi/2\),
  \(\sigma \to \mp\infty\) the tips \(z \to 0, \pi\); one period of
  \(x_0\) is the whole line \(-\infty < \sigma < \infty\). Unit: length.
  Status: DERIVED.
\end{itemize}

The field equation of dirac16complex in the bulk: the variation of IN-4
with respect to \(\bar\Psi\) (the second half of the symmetric kinetic
term, integrated by parts with \(D_M\gamma^M = 0\), doubles the first).
With the author's real matrices the operator \(\gamma^MD_M - m_\Psi\) is
real: the field equation does not contain \(\mathrm{i}\), and a real
solution exists for every complex one. Status: RECORD.

\[
\begin{aligned}
&\Psi\big\vert_{\sigma=0^+} = \exp\Big(\frac{1}{V_*}\,\gamma^{(0)}\,\mathbb{M}\Big)\,\Psi\big\vert_{\sigma=0^-},\\
&\mathbb{M} = c_\Phi\big(\Phi^\dagger\Phi\big)P_{\mathrm{L}} + \sum_X\Big[\frac14\big(c^{G}_X\,G^A_{\mu\nu}G^{A\mu\nu}
 + c^{W}_X\,W^I_{\mu\nu}W^{I\mu\nu} + c^{B}_X\,B_{\mu\nu}B^{\mu\nu}\big)P_X + \frac{d_X}{2}\,B_{\mu\nu}\,\mathrm{i}\Sigma^{\mu\nu}P_X\\
&\qquad\qquad + \sum_f c^{f}_X\big(2\mathrm{i}\,J_f^\mu P_XD_\mu + \mathrm{i}(\hat\nabla_\mu J_f^\mu)P_X\big)\Big]
\end{aligned} \tag{FE-13}
\] \textbf{Symbols in (FE-13).}

\begin{itemize}
\tightlist
\item
  \(\Psi,\ \bar\Psi\): dirac16complex: the author's 16-component complex
  spinor field of eight dimensions;
  \(\bar\Psi = \Psi^\dagger\sigma_{16}\) is its Dirac adjoint; in a term
  of the fold, \(\Psi\) is its value at \(\sigma = 0\). Unit:
  GeV\(^{7/2}\). Status: AUTHOR.
\item
  \(\sigma\): The proper coordinate of the hidden direction:
  \(\mathrm{d}\sigma = \lvert\cot z\rvert\,\mathrm{d}x_0\),
  \(\sigma = \mathrm{sgn}(\pi/2 - z)\ln(\sin z)/(6H)\) on one period
  \(0 < z < \pi\); \(\sigma = 0\) is the fold \(z = \pi/2\),
  \(\sigma \to \mp\infty\) the tips \(z \to 0, \pi\); one period of
  \(x_0\) is the whole line \(-\infty < \sigma < \infty\). Unit: length.
  Status: DERIVED.
\item
  \(V_*\): The physical volume of the extra times \(x_5, x_6, x_7\) on
  the fold in the late universe,
  \(V_* = V_{\mathrm{t}}(\mathrm{e}^{-a_4}b_4)^3\) at a time when
  \(\mathrm{e}^{-a_4}b_4\) is constant; it converts the
  eight-dimensional normalisation into the four-dimensional one. Unit:
  length\(^3\) (GeV\(^{-3}\)). Status: ASSUMED.
\item
  \(\gamma^M,\ \gamma^{(a)}\): The Dirac matrices of dirac16complex
  along the coordinates, \(\gamma^M = \mathbf{e}_a{}^MT^a\),
  \(\gamma_M = g_{MN}\gamma^N\); \(\gamma^{(a)} = T^a\) along the frame
  direction \((a)\) (\(\gamma^{(0)}\) along \(\sigma\), \(\gamma^{(4)}\)
  along the time). Unit: dimensionless. Status: AUTHOR.
\item
  \(\mathbb{M}\): The fold operator of dirac16complex: the
  \(16 \times 16\) matrix (with a derivative along the fold) built from
  the Standard-Model fields at the fold by varying the portals with
  respect to \(\bar\Psi\) (FE-13). Unit: GeV\(^{-3}\) (so that
  \(\mathbb{M}/V_*\) is dimensionless). Status: DERIVED.
\item
  \(c_\Phi\): The Higgs-portal coupling of the left-handed half of
  dirac16complex (the right-handed half has none). Unit: GeV\(^{-5}\).
  Status: INVENTED.
\item
  \(\Phi\): The Higgs doublet (two complex components, hypercharge
  \(1/2\)), a field of the fold. Unit: GeV. Status: AUTHOR.
\item
  \({}^\dagger\): Hermitian conjugation: the complex conjugate of the
  transpose of a matrix or spinor (for field operators the adjoint
  operator). Status: STANDARD.
\item
  \(P_{\mathrm{L}}\): The author's left-handed chiral projector,
  \(P_{\mathrm{L}} = (\mathbf{1} - \Gamma_8)/2\). Unit: dimensionless.
  Status: AUTHOR.
\item
  \(X\): The label of a half of dirac16complex:
  \(X \in \{\mathrm{L}, \mathrm{R}\}\) (left-handed or right-handed
  half); a sum over \(X\) runs over both. Status: STANDARD.
\item
  \(c^{G}_X\): The gluon-kinetic portal coupling of the half \(X\) of
  dirac16complex. Unit: GeV\(^{-7}\). Status: INVENTED.
\item
  \(G^A_\mu,\ G^A_{\mu\nu}\): The gluon field of \(SU(3)_C\) on the fold
  (\(A = 1, \dots, 8\)) and its field strength. Unit: GeV (field),
  GeV\(^2\) (field strength). Status: AUTHOR.
\item
  \(A\): A colour (adjoint SU(3)) index running over \(1, \dots, 8\); a
  repeated index is summed. Status: STANDARD.
\item
  \(\mu\): A four-dimensional coordinate index running over the
  observable coordinates \(x_1, x_2, x_3, x_4\) of the fold (signature
  \(+,+,+,-\)); a repeated index is summed. Status: STANDARD.
\item
  \(\nu\): A four-dimensional coordinate index (as \(\mu\)). Status:
  STANDARD.
\item
  \(c^{W}_X\): The \(W\)-kinetic portal coupling of the half \(X\).
  Unit: GeV\(^{-7}\). Status: INVENTED.
\item
  \(W^I_\mu,\ W^I_{\mu\nu}\): The gauge field of \(SU(2)_L\) on the
  fold, \(W^1_\mu, W^2_\mu, W^3_\mu\), and its field strength. Unit: GeV
  (field), GeV\(^2\) (field strength). Status: AUTHOR.
\item
  \(I\): A weak-isospin (adjoint SU(2)) index running over \(1, 2, 3\);
  a repeated index is summed. Status: STANDARD.
\item
  \(c^{B}_X\): The \(B\)-kinetic portal coupling of the half \(X\).
  Unit: GeV\(^{-7}\). Status: INVENTED.
\item
  \(B_\mu,\ B_{\mu\nu}\): The hypercharge gauge field of \(U(1)_Y\) on
  the fold and its field strength. Unit: GeV (field), GeV\(^2\) (field
  strength). Status: AUTHOR.
\item
  \(P_X\): The chiral projector of the half \(X\):
  \(P_X = P_{\mathrm{L}}\) for \(X = \mathrm{L}\), \(P_{\mathrm{R}}\)
  for \(X = \mathrm{R}\). Unit: dimensionless. Status: AUTHOR.
\item
  \(d_X\): The dipole-portal coupling of the half \(X\) of
  dirac16complex with the hypercharge field strength. Unit:
  GeV\(^{-5}\). Status: INVENTED.
\item
  \(\mathrm{i}\): The imaginary unit, \(\mathrm{i}^2 = -1\). Unit:
  dimensionless. Status: STANDARD.
\item
  \(\Sigma^{ab},\ \Sigma^{\mu\nu}\): The generators of Spin(4,4) on
  dirac16complex, \(\Sigma^{ab} = \frac14[T^a, T^b]\) (the scaled
  commutators of the author's matrices);
  \(\Sigma^{\mu\nu} = \mathbf{e}_a{}^\mu\mathbf{e}_b{}^\nu\Sigma^{ab}\)
  along the fold. Unit: dimensionless. Status: AUTHOR.
\item
  \(f\): The label of a kind of Standard-Model multiplet:
  \(f \in \{Q, u, d, \ell, e, \nu\}\) (quark doublets, up-quark
  singlets, down-quark singlets, lepton doublets, charged-lepton
  singlets, right-handed neutrinos); a sum over \(f\) runs over the six.
  Status: STANDARD.
\item
  \(c^{f}_X\): The current-portal coupling of the half \(X\) of
  dirac16complex with the fermion-number current of the multiplets of
  kind \(f\) (flavour blind). Unit: GeV\(^{-7}\). Status: INVENTED.
\item
  \(J_f^\mu\): The fermion-number current of the Standard-Model
  multiplets of kind \(f\), \(J_f^\mu = \bar\psi_f\hat\gamma^\mu\psi_f\)
  summed over the three generations and the gauge components; with other
  labels the gauge currents (FE-16 to FE-18). Unit: GeV\(^3\). Status:
  STANDARD.
\item
  \(D_M,\ D_\mu\): The covariant derivative: on dirac16complex
  \(D_M = \partial_M + \Omega_M\) (with
  \(D_M\bar\Psi = \partial_M\bar\Psi - \bar\Psi\Omega_M\)); on the
  fields of the fold the gauge and spin covariant derivative \(D_\mu\)
  of IN-13; \(D^\mu = \hat h^{\mu\nu}D_\nu\); on an adjoint field
  \((D_\mu X)^A = \partial_\mu X^A + g_3f^{ABC}G^B_\mu X^C\) (likewise
  for \(SU(2)\) with \(g_2\epsilon^{IJK}\)). Unit: 1/length. Status:
  STANDARD.
\item
  \(\hat\nabla\): The covariant derivative of the fold: with the
  Levi-Civita connection of the metric of the fold (\(\hat h_{\mu\nu}\)
  for four-dimensional tensors, the induced metric of the
  seven-dimensional fold for \(\hat\nabla_a\mathcal{S}^a{}_b\)). Status:
  STANDARD.
\end{itemize}

The equation of dirac16complex at the fold: the variation of the portals
IN-16 with respect to \(\bar\Psi\) adds
\(-\delta_{\mathrm{D}}(\sigma)\,\mathbb{M}\Psi/V_*\) to FE-12.
\(\mathbb{M}\) is the fold operator, a \(16 \times 16\) matrix (with a
derivative along the fold from the current portals) built from the
Standard-Model fields at the fold. Integrating FE-12 with this term
across a fold of vanishing thickness gives the matching condition in the
first line: the values of dirac16complex on the two sides of the fold
differ by the exponential of \(\gamma^{(0)}\mathbb{M}/V_*\) (the exact
solution of
\(\partial_\sigma\Psi = \gamma^{(0)}\mathbb{M}\,\delta_{\mathrm{D}}(\sigma)\Psi/V_*\)
in a layer in which the fold term dominates;
\((\gamma^{(0)})^2 = \mathbf{1}\)). The Higgs term acts on the
left-handed half only. Status: DERIVED (from INVENTED terms).

\[
\gamma^M\Omega_M = -3H\,\mathrm{sgn}(\sigma)\,\gamma^{(0)} + \frac32\,\frac{b_4'}{b_4}\,\gamma^{(4)} \ne 0 \tag{FE-14}
\] \textbf{Symbols in (FE-14).}

\begin{itemize}
\tightlist
\item
  \(\gamma^M,\ \gamma^{(a)}\): The Dirac matrices of dirac16complex
  along the coordinates, \(\gamma^M = \mathbf{e}_a{}^MT^a\),
  \(\gamma_M = g_{MN}\gamma^N\); \(\gamma^{(a)} = T^a\) along the frame
  direction \((a)\) (\(\gamma^{(0)}\) along \(\sigma\), \(\gamma^{(4)}\)
  along the time). Unit: dimensionless. Status: AUTHOR.
\item
  \(M\): A coordinate index of eight dimensions (also with a numbered
  copy \(M_1, M_2, \dots\)): it runs over \(x_0, x_1, \dots, x_7\); a
  repeated index is summed. Status: STANDARD.
\item
  \(\Omega_M\): The spin connection acting on dirac16complex,
  \(\Omega_M = \frac12\omega_{Mab}\Sigma^{ab}\). Unit: 1/length. Status:
  RECORD.
\item
  \(H\): The author's constant (an inverse length) that fixes the period
  \(\pi/(6H)\) of \(x_0\), the Lovelock couplings and the cosmological
  constant \(\Lambda = -10H^2\); in UM-3 a free parameter. Unit: GeV
  (1/length). Status: AUTHOR.
\item
  \(\mathrm{sgn}\): The sign function: \(\mathrm{sgn}(y) = +1\) for
  \(y > 0\), \(-1\) for \(y < 0\). Status: STANDARD.
\item
  \(\sigma\): The proper coordinate of the hidden direction:
  \(\mathrm{d}\sigma = \lvert\cot z\rvert\,\mathrm{d}x_0\),
  \(\sigma = \mathrm{sgn}(\pi/2 - z)\ln(\sin z)/(6H)\) on one period
  \(0 < z < \pi\); \(\sigma = 0\) is the fold \(z = \pi/2\),
  \(\sigma \to \mp\infty\) the tips \(z \to 0, \pi\); one period of
  \(x_0\) is the whole line \(-\infty < \sigma < \infty\). Unit: length.
  Status: DERIVED.
\item
  \(b_4\): The author's metric function of the time \(x_4\) in the
  extra-time components, \(b_4 > 0\), never identically 1 (the author's
  order of 2026-10-24); the scale factor of the extra times is
  \(\mathrm{e}^{-a_4}b_4\) times the warp. Unit: dimensionless. Status:
  AUTHOR.
\item
  \({}'\): The derivative of a function of the time \(x_4\) alone with
  respect to \(x_4\) (\(a_4' = \mathrm{d}a_4/\mathrm{d}x_4\); two primes
  the second derivative); on a function of \(\sigma\) alone (a profile)
  the derivative with respect to \(\sigma\). Status: STANDARD.
\end{itemize}

The coupling of dirac16complex to gravity through the canonical spin
connection, for the author's metric: the contraction
\(\gamma^M\Omega_M\) that enters the Dirac operator. The first term
comes from the warp of the six directions
\(x_1, x_2, x_3, x_5, x_6, x_7\) (each contributes
\(\frac12\partial_\sigma\) of the logarithm of its scale factor), the
second from the time dependence of their scale factors (\(a_4\) cancels
between ordinary space and the extra times). The first term never
vanishes (\(H > 0\), and it only changes sign at the fold); the second
is non-zero whenever \(b_4\) depends on \(x_4\), as the author requires.
Hence the coupling of dirac16complex to gravity through the spin
connection is never identically zero (DERIVED,
\texttt{um3\_geometry.py}, G5, with the author's matrices; the vielbein
postulate \(D_M\gamma^N = 0\) holds, G5). For the general warped metric
FE-34,
\(\gamma^M\Omega_M = \frac12\partial_\sigma(3\mathfrak{a} + 3\mathfrak{b} + \mathfrak{n})\gamma^{(0)} + \frac12\mathrm{e}^{-\mathfrak{n}}\partial_{x_4}(3\mathfrak{a} + 3\mathfrak{b})\gamma^{(4)}\).
Status: DERIVED.

\hypertarget{the-higgs-doublet}{%
\subsection{The Higgs doublet}\label{the-higgs-doublet}}

\[
\begin{aligned}
&\frac{1}{\sqrt{-\hat h}\,\mathcal{Z}}\,D_\mu\big(\sqrt{-\hat h}\,\mathcal{Z}\,\hat h^{\mu\nu}D_\nu\Phi\big) - \mu^2\Phi
 - 2\lambda\big(\Phi^\dagger\Phi\big)\Phi - c_\Phi\,\bar\Psi P_{\mathrm{L}}\Psi\,\Phi = \Theta,\\
&\Theta = \begin{pmatrix}\bar d_{\mathrm{R}}(\mathcal{Y}^{d})^\dagger u_{\mathrm{L}} + \bar e_{\mathrm{R}}(\mathcal{Y}^{e})^\dagger\nu_{\mathrm{L}}
 - \bar d_{\mathrm{L}}\mathcal{Y}^{u}u_{\mathrm{R}} - \bar e_{\mathrm{L}}\mathcal{Y}^{\nu}\nu_{\mathrm{R}}\\
 \bar d_{\mathrm{R}}(\mathcal{Y}^{d})^\dagger d_{\mathrm{L}} + \bar e_{\mathrm{R}}(\mathcal{Y}^{e})^\dagger e_{\mathrm{L}}
 + \bar u_{\mathrm{L}}\mathcal{Y}^{u}u_{\mathrm{R}} + \bar\nu_{\mathrm{L}}\mathcal{Y}^{\nu}\nu_{\mathrm{R}}\end{pmatrix}
\end{aligned} \tag{FE-15}
\] \textbf{Symbols in (FE-15).}

\begin{itemize}
\tightlist
\item
  \(\hat h_{\mu\nu}\): The four-dimensional metric of the fold,
  \(\hat h_{\mu\nu} = \mathrm{diag}(\mathrm{e}^{2a_4}, \mathrm{e}^{2a_4}, \mathrm{e}^{2a_4}, -1)\)
  in the coordinates \(x_1, \dots, x_4\) (the induced metric at
  \(\sigma = 0\) without the extra times); \(\hat h^{\mu\nu}\) its
  inverse, \(\hat h\) its determinant,
  \(\sqrt{-\hat h} = \mathrm{e}^{3a_4}\); \(\hat h^{(1)}_{\mu\nu}\) its
  first-order cosmological perturbation. Unit: dimensionless. Status:
  DERIVED.
\item
  \(\mathcal{Z}\): The fold factor,
  \(\mathcal{Z}(x_4) = V_{\mathrm{t}}(\mathrm{e}^{-a_4}b_4)^3/V_*\): the
  factor by which the size of the extra times multiplies every term of
  the fold; \(\mathcal{Z} = 1\) when the extra times are frozen at their
  late value. Unit: dimensionless. Value: 1 in the late universe.
  Status: DERIVED.
\item
  \(D_M,\ D_\mu\): The covariant derivative: on dirac16complex
  \(D_M = \partial_M + \Omega_M\) (with
  \(D_M\bar\Psi = \partial_M\bar\Psi - \bar\Psi\Omega_M\)); on the
  fields of the fold the gauge and spin covariant derivative \(D_\mu\)
  of IN-13; \(D^\mu = \hat h^{\mu\nu}D_\nu\); on an adjoint field
  \((D_\mu X)^A = \partial_\mu X^A + g_3f^{ABC}G^B_\mu X^C\) (likewise
  for \(SU(2)\) with \(g_2\epsilon^{IJK}\)). Unit: 1/length. Status:
  STANDARD.
\item
  \(\mu\): A four-dimensional coordinate index running over the
  observable coordinates \(x_1, x_2, x_3, x_4\) of the fold (signature
  \(+,+,+,-\)); a repeated index is summed. Status: STANDARD.
\item
  \(\nu\): A four-dimensional coordinate index (as \(\mu\)). Status:
  STANDARD.
\item
  \(\Phi\): The Higgs doublet (two complex components, hypercharge
  \(1/2\)), a field of the fold. Unit: GeV. Status: AUTHOR.
\item
  \(\mu^2\): The mass parameter of the Higgs potential, \(\mu^2 < 0\)
  (it enters only as \(\mu^2\)). Unit: GeV\(^2\). Value:
  \(\mu^2 = -\lambda v^2 = -m_h^2/2 = -(88.6\ \mathrm{GeV})^2\). Status:
  AUTHOR.
\item
  \(\lambda\): The quartic coupling of the Higgs potential,
  \(\lambda > 0\). Unit: dimensionless. Value:
  \(\lambda = m_h^2/(2v^2) = 0.129\). Status: AUTHOR.
\item
  \({}^\dagger\): Hermitian conjugation: the complex conjugate of the
  transpose of a matrix or spinor (for field operators the adjoint
  operator). Status: STANDARD.
\item
  \(c_\Phi\): The Higgs-portal coupling of the left-handed half of
  dirac16complex (the right-handed half has none). Unit: GeV\(^{-5}\).
  Status: INVENTED.
\item
  \(\Psi,\ \bar\Psi\): dirac16complex: the author's 16-component complex
  spinor field of eight dimensions;
  \(\bar\Psi = \Psi^\dagger\sigma_{16}\) is its Dirac adjoint; in a term
  of the fold, \(\Psi\) is its value at \(\sigma = 0\). Unit:
  GeV\(^{7/2}\). Status: AUTHOR.
\item
  \(P_{\mathrm{L}}\): The author's left-handed chiral projector,
  \(P_{\mathrm{L}} = (\mathbf{1} - \Gamma_8)/2\). Unit: dimensionless.
  Status: AUTHOR.
\item
  \(\Theta\): The Yukawa source of the Higgs equation,
  \(\Theta = -\partial\mathcal{L}_{\mathrm{Y}}/\partial\Phi^\dagger\) (a
  doublet of fermion bilinears). Unit: GeV\(^3\). Status: STANDARD.
\item
  \(d_{\mathrm{R}}\): The right-handed down-type quarks \(d, s, b\)
  (hypercharge \(-1/3\)). Unit: GeV\(^{3/2}\). Status: STANDARD.
\item
  \(\mathcal{Y}^{u},\ \mathcal{Y}^{d},\ \mathcal{Y}^{e},\ \mathcal{Y}^{\nu}\):
  The Yukawa matrices (\(3 \times 3\) complex matrices in generation
  space) of the up quarks, the down quarks, the charged leptons and the
  neutrinos. Unit: dimensionless. Status: STANDARD.
\item
  \(u_{\mathrm{L}}\): The upper (up-type, left-handed) components of the
  quark doublets. Unit: GeV\(^{3/2}\). Status: STANDARD.
\item
  \(e_{\mathrm{R}}\): The right-handed electron, muon and tau (singlets,
  hypercharge \(-1\)). Unit: GeV\(^{3/2}\). Status: AUTHOR.
\item
  \(\nu_{\mathrm{L}}\): The upper (neutral, left-handed) components of
  the lepton doublets: the left-handed neutrinos. Unit: GeV\(^{3/2}\).
  Status: STANDARD.
\item
  \(d_{\mathrm{L}}\): The lower (down-type, left-handed) components of
  the quark doublets. Unit: GeV\(^{3/2}\). Status: STANDARD.
\item
  \(u_{\mathrm{R}}\): The right-handed up-type quarks \(u, c, t\)
  (singlets of \(SU(2)_L\), hypercharge \(2/3\)). Unit: GeV\(^{3/2}\).
  Status: STANDARD.
\item
  \(\nu_{\mathrm{R}}\): The three right-handed neutrinos (gauge
  singlets). Unit: GeV\(^{3/2}\). Status: ASSUMED.
\item
  \(e_{\mathrm{L}}\): The lower (charged, left-handed) components of the
  lepton doublets. Unit: GeV\(^{3/2}\). Status: STANDARD.
\end{itemize}

The field equation of the Higgs doublet on the fold: the variation of
IN-11, IN-15 and IN-16 with respect to \(\Phi^\dagger\). The factor
\(\mathcal{Z}\) (IN-8) is 1 in the late universe. The Higgs portal adds
the term \(-c_\Phi\bar\Psi P_{\mathrm{L}}\Psi\,\Phi\): the scalar
density of the left-handed half of dirac16complex at the fold shifts the
mass parameter of the Higgs field. \(\Theta\) is the Yukawa source,
\(-\partial\mathcal{L}_{\mathrm{Y}}/\partial\Phi^\dagger\), written with
the components \(Q = (u_{\mathrm{L}}, d_{\mathrm{L}})\) and
\(\ell = (\nu_{\mathrm{L}}, e_{\mathrm{L}})\) of the doublets
(generation indices contracted with the Yukawa matrices). Status:
STANDARD (AUTHOR potential; INVENTED portal term).

\hypertarget{the-gauge-fields-1}{%
\subsection{The gauge fields}\label{the-gauge-fields-1}}

\[
\begin{aligned}
&\frac{1}{\sqrt{-\hat h}\,\mathcal{Z}}\Big(\partial_\mu\big[\sqrt{-\hat h}\,\mathcal{Z}\,\mathcal{K}_{G}\,G^{A\mu\nu}\big]
 + g_3\,f^{ABC}G^B_\mu\,\sqrt{-\hat h}\,\mathcal{Z}\,\mathcal{K}_{G}\,G^{C\mu\nu}\Big) = -g_3\,J_{\mathrm{c}}^{A\nu},\qquad
 \mathcal{K}_{G} = 1 + \sum_X c^{G}_X\,\bar\Psi P_X\Psi,\\
&J_{\mathrm{c}}^{A\nu} = \bar Q\,\hat\gamma^\nu t^AQ + \bar u_{\mathrm{R}}\,\hat\gamma^\nu t^Au_{\mathrm{R}}
 + \bar d_{\mathrm{R}}\,\hat\gamma^\nu t^Ad_{\mathrm{R}}
\end{aligned} \tag{FE-16}
\] \textbf{Symbols in (FE-16).}

\begin{itemize}
\tightlist
\item
  \(\hat h_{\mu\nu}\): The four-dimensional metric of the fold,
  \(\hat h_{\mu\nu} = \mathrm{diag}(\mathrm{e}^{2a_4}, \mathrm{e}^{2a_4}, \mathrm{e}^{2a_4}, -1)\)
  in the coordinates \(x_1, \dots, x_4\) (the induced metric at
  \(\sigma = 0\) without the extra times); \(\hat h^{\mu\nu}\) its
  inverse, \(\hat h\) its determinant,
  \(\sqrt{-\hat h} = \mathrm{e}^{3a_4}\); \(\hat h^{(1)}_{\mu\nu}\) its
  first-order cosmological perturbation. Unit: dimensionless. Status:
  DERIVED.
\item
  \(\mathcal{Z}\): The fold factor,
  \(\mathcal{Z}(x_4) = V_{\mathrm{t}}(\mathrm{e}^{-a_4}b_4)^3/V_*\): the
  factor by which the size of the extra times multiplies every term of
  the fold; \(\mathcal{Z} = 1\) when the extra times are frozen at their
  late value. Unit: dimensionless. Value: 1 in the late universe.
  Status: DERIVED.
\item
  \(\partial\): A partial derivative: \(\partial_M\) with respect to the
  coordinate \(X^M\), \(\partial_\mu\) with respect to \(x^\mu\),
  \(\partial_\sigma\) with respect to \(\sigma\),
  \(\partial X/\partial Y\) the derivative of \(X\) with respect to
  \(Y\) at fixed other variables. Status: STANDARD.
\item
  \(\mu\): A four-dimensional coordinate index running over the
  observable coordinates \(x_1, x_2, x_3, x_4\) of the fold (signature
  \(+,+,+,-\)); a repeated index is summed. Status: STANDARD.
\item
  \(\mathcal{K}_{G},\ \mathcal{K}_{W},\ \mathcal{K}_{B}\): The portal
  factors of the gauge kinetic terms,
  \(\mathcal{K}_{G} = 1 + \sum_Xc^{G}_X\bar\Psi P_X\Psi\) (likewise for
  \(W\) and \(B\)): the gauge couplings at the fold are divided by
  \(\sqrt{\mathcal{K}}\). Unit: dimensionless. Value: 1 where
  dirac16complex vanishes. Status: DERIVED.
\item
  \(G^A_\mu,\ G^A_{\mu\nu}\): The gluon field of \(SU(3)_C\) on the fold
  (\(A = 1, \dots, 8\)) and its field strength. Unit: GeV (field),
  GeV\(^2\) (field strength). Status: AUTHOR.
\item
  \(A\): A colour (adjoint SU(3)) index running over \(1, \dots, 8\); a
  repeated index is summed. Status: STANDARD.
\item
  \(\nu\): A four-dimensional coordinate index (as \(\mu\)). Status:
  STANDARD.
\item
  \(g_3\): The \(SU(3)_C\) gauge coupling (the strong coupling \(g_s\)).
  Unit: dimensionless. Value: \(g_3(m_Z) = 1.2177\)
  (\(\alpha_s = 0.1180\)). Status: STANDARD.
\item
  \(f^{ABC}\): The structure constants of \(SU(3)\),
  \([t^A, t^B] = \mathrm{i}f^{ABC}t^C\). Unit: dimensionless. Status:
  STANDARD.
\item
  \(B\): A colour (adjoint SU(3)) index (as \(A\)). Status: STANDARD.
\item
  \(C\): A colour (adjoint SU(3)) index (as \(A\)). Status: STANDARD.
\item
  \(J_{\mathrm{c}}^{A\nu}\): The colour current of the quarks. Unit:
  GeV\(^3\). Status: STANDARD.
\item
  \(X\): The label of a half of dirac16complex:
  \(X \in \{\mathrm{L}, \mathrm{R}\}\) (left-handed or right-handed
  half); a sum over \(X\) runs over both. Status: STANDARD.
\item
  \(c^{G}_X\): The gluon-kinetic portal coupling of the half \(X\) of
  dirac16complex. Unit: GeV\(^{-7}\). Status: INVENTED.
\item
  \(\Psi,\ \bar\Psi\): dirac16complex: the author's 16-component complex
  spinor field of eight dimensions;
  \(\bar\Psi = \Psi^\dagger\sigma_{16}\) is its Dirac adjoint; in a term
  of the fold, \(\Psi\) is its value at \(\sigma = 0\). Unit:
  GeV\(^{7/2}\). Status: AUTHOR.
\item
  \(P_X\): The chiral projector of the half \(X\):
  \(P_X = P_{\mathrm{L}}\) for \(X = \mathrm{L}\), \(P_{\mathrm{R}}\)
  for \(X = \mathrm{R}\). Unit: dimensionless. Status: AUTHOR.
\item
  \(Q,\ \bar Q\): The left-handed quark doublets of the three
  generations, \(Q = (u_{\mathrm{L}}, d_{\mathrm{L}})\) (colour
  triplets, hypercharge \(1/6\)). Unit: GeV\(^{3/2}\). Status: AUTHOR.
\item
  \(\hat\gamma^\mu\): The four-dimensional complex \(4 \times 4\) Dirac
  matrices of the fold,
  \(\hat\gamma^\mu = \hat{\mathbf{e}}_a{}^\mu\hat\gamma^{(a)}\),
  \(\{\hat\gamma^\mu, \hat\gamma^\nu\} = -2\hat h^{\mu\nu}\mathbf{1}\);
  \(\hat\gamma_\mu = \hat h_{\mu\nu}\hat\gamma^\nu\). Unit:
  dimensionless. Status: STANDARD.
\item
  \(t^A\): The generators of \(SU(3)\) on colour triplets,
  \(t^A = \lambda_{\mathrm{GM}}^A/2\) (the Gell-Mann matrices halved),
  \(\mathrm{tr}(t^At^B) = \delta^{AB}/2\). Unit: dimensionless. Status:
  STANDARD.
\item
  \(u_{\mathrm{R}}\): The right-handed up-type quarks \(u, c, t\)
  (singlets of \(SU(2)_L\), hypercharge \(2/3\)). Unit: GeV\(^{3/2}\).
  Status: STANDARD.
\item
  \(d_{\mathrm{R}}\): The right-handed down-type quarks \(d, s, b\)
  (hypercharge \(-1/3\)). Unit: GeV\(^{3/2}\). Status: STANDARD.
\end{itemize}

The Yang-Mills equation of the gluons (with the convention of the
signature \((+,+,+,-)\): \(\nabla_\mu F^{\mu\nu} = -J^\nu\) for a
Lagrangian \(-\frac14F^2 + A_\mu J^\mu\)). The gluon-kinetic portals
multiply the gluon kinetic term by \(\mathcal{K}_{G}\): the strong
coupling at the fold is \(g_3/\sqrt{\mathcal{K}_{G}}\) where
dirac16complex is dense. The colour current \(J_{\mathrm{c}}^{A\nu}\) of
the quarks is the source. Status: STANDARD (INVENTED portal factor).

\[
\begin{aligned}
&\frac{1}{\sqrt{-\hat h}\,\mathcal{Z}}\Big(\partial_\mu\big[\sqrt{-\hat h}\,\mathcal{Z}\,\mathcal{K}_{W}\,W^{I\mu\nu}\big]
 + g_2\,\epsilon^{IJK}W^J_\mu\,\sqrt{-\hat h}\,\mathcal{Z}\,\mathcal{K}_{W}\,W^{K\mu\nu}\Big)
 = -g_2\big(J_{\mathrm{w}}^{I\nu} + J_\Phi^{I\nu}\big),\\
&\mathcal{K}_{W} = 1 + \sum_X c^{W}_X\,\bar\Psi P_X\Psi,\qquad J_{\mathrm{w}}^{I\nu} = \bar Q\,\hat\gamma^\nu\frac{\tau^I}{2}Q + \bar\ell\,\hat\gamma^\nu\frac{\tau^I}{2}\ell,\qquad
 J_\Phi^{I\nu} = \mathrm{i}\Big[(D^\nu\Phi)^\dagger\frac{\tau^I}{2}\Phi - \Phi^\dagger\frac{\tau^I}{2}D^\nu\Phi\Big]
\end{aligned} \tag{FE-17}
\] \textbf{Symbols in (FE-17).}

\begin{itemize}
\tightlist
\item
  \(\hat h_{\mu\nu}\): The four-dimensional metric of the fold,
  \(\hat h_{\mu\nu} = \mathrm{diag}(\mathrm{e}^{2a_4}, \mathrm{e}^{2a_4}, \mathrm{e}^{2a_4}, -1)\)
  in the coordinates \(x_1, \dots, x_4\) (the induced metric at
  \(\sigma = 0\) without the extra times); \(\hat h^{\mu\nu}\) its
  inverse, \(\hat h\) its determinant,
  \(\sqrt{-\hat h} = \mathrm{e}^{3a_4}\); \(\hat h^{(1)}_{\mu\nu}\) its
  first-order cosmological perturbation. Unit: dimensionless. Status:
  DERIVED.
\item
  \(\mathcal{Z}\): The fold factor,
  \(\mathcal{Z}(x_4) = V_{\mathrm{t}}(\mathrm{e}^{-a_4}b_4)^3/V_*\): the
  factor by which the size of the extra times multiplies every term of
  the fold; \(\mathcal{Z} = 1\) when the extra times are frozen at their
  late value. Unit: dimensionless. Value: 1 in the late universe.
  Status: DERIVED.
\item
  \(\partial\): A partial derivative: \(\partial_M\) with respect to the
  coordinate \(X^M\), \(\partial_\mu\) with respect to \(x^\mu\),
  \(\partial_\sigma\) with respect to \(\sigma\),
  \(\partial X/\partial Y\) the derivative of \(X\) with respect to
  \(Y\) at fixed other variables. Status: STANDARD.
\item
  \(\mu\): A four-dimensional coordinate index running over the
  observable coordinates \(x_1, x_2, x_3, x_4\) of the fold (signature
  \(+,+,+,-\)); a repeated index is summed. Status: STANDARD.
\item
  \(\mathcal{K}_{G},\ \mathcal{K}_{W},\ \mathcal{K}_{B}\): The portal
  factors of the gauge kinetic terms,
  \(\mathcal{K}_{G} = 1 + \sum_Xc^{G}_X\bar\Psi P_X\Psi\) (likewise for
  \(W\) and \(B\)): the gauge couplings at the fold are divided by
  \(\sqrt{\mathcal{K}}\). Unit: dimensionless. Value: 1 where
  dirac16complex vanishes. Status: DERIVED.
\item
  \(W^I_\mu,\ W^I_{\mu\nu}\): The gauge field of \(SU(2)_L\) on the
  fold, \(W^1_\mu, W^2_\mu, W^3_\mu\), and its field strength. Unit: GeV
  (field), GeV\(^2\) (field strength). Status: AUTHOR.
\item
  \(I\): A weak-isospin (adjoint SU(2)) index running over \(1, 2, 3\);
  a repeated index is summed. Status: STANDARD.
\item
  \(\nu\): A four-dimensional coordinate index (as \(\mu\)). Status:
  STANDARD.
\item
  \(g_2\): The \(SU(2)_L\) gauge coupling (usually \(g\)). Unit:
  dimensionless. Value: \(g_2(m_Z) = 0.6517\). Status: STANDARD.
\item
  \(\epsilon^{IJK}\): The Levi-Civita symbol of three indices (the
  structure constants of \(SU(2)\)), \(\epsilon^{123} = 1\). Unit:
  dimensionless. Status: STANDARD.
\item
  \(J\): A weak-isospin (adjoint SU(2)) index (as \(I\)). Status:
  STANDARD.
\item
  \(K\): A weak-isospin (adjoint SU(2)) index (as \(I\)). Status:
  STANDARD.
\item
  \(J_{\mathrm{w}}^{I\nu}\): The weak-isospin current of the left-handed
  fermion doublets. Unit: GeV\(^3\). Status: STANDARD.
\item
  \(J_\Phi^{I\nu}\): The weak-isospin current of the Higgs doublet (it
  contains the mass terms of the \(W\) fields); \(J_\Phi^{2\,x_4}\) its
  isospin charge density with \(I = 2\). Unit: GeV\(^3\). Status:
  STANDARD.
\item
  \(X\): The label of a half of dirac16complex:
  \(X \in \{\mathrm{L}, \mathrm{R}\}\) (left-handed or right-handed
  half); a sum over \(X\) runs over both. Status: STANDARD.
\item
  \(c^{W}_X\): The \(W\)-kinetic portal coupling of the half \(X\).
  Unit: GeV\(^{-7}\). Status: INVENTED.
\item
  \(\Psi,\ \bar\Psi\): dirac16complex: the author's 16-component complex
  spinor field of eight dimensions;
  \(\bar\Psi = \Psi^\dagger\sigma_{16}\) is its Dirac adjoint; in a term
  of the fold, \(\Psi\) is its value at \(\sigma = 0\). Unit:
  GeV\(^{7/2}\). Status: AUTHOR.
\item
  \(P_X\): The chiral projector of the half \(X\):
  \(P_X = P_{\mathrm{L}}\) for \(X = \mathrm{L}\), \(P_{\mathrm{R}}\)
  for \(X = \mathrm{R}\). Unit: dimensionless. Status: AUTHOR.
\item
  \(Q,\ \bar Q\): The left-handed quark doublets of the three
  generations, \(Q = (u_{\mathrm{L}}, d_{\mathrm{L}})\) (colour
  triplets, hypercharge \(1/6\)). Unit: GeV\(^{3/2}\). Status: AUTHOR.
\item
  \(\hat\gamma^\mu\): The four-dimensional complex \(4 \times 4\) Dirac
  matrices of the fold,
  \(\hat\gamma^\mu = \hat{\mathbf{e}}_a{}^\mu\hat\gamma^{(a)}\),
  \(\{\hat\gamma^\mu, \hat\gamma^\nu\} = -2\hat h^{\mu\nu}\mathbf{1}\);
  \(\hat\gamma_\mu = \hat h_{\mu\nu}\hat\gamma^\nu\). Unit:
  dimensionless. Status: STANDARD.
\item
  \(\tau^I\): The Pauli matrices acting on weak doublets
  (\(\tau^1, \tau^2, \tau^3\); \(\tau^2\) is the second one, not a
  square). Unit: dimensionless. Status: STANDARD.
\item
  \(\ell,\ \bar\ell\): The left-handed lepton doublets of the three
  generations, \(\ell = (\nu_{\mathrm{L}}, e_{\mathrm{L}})\)
  (hypercharge \(-1/2\)). Unit: GeV\(^{3/2}\). Status: AUTHOR.
\item
  \(\mathrm{i}\): The imaginary unit, \(\mathrm{i}^2 = -1\). Unit:
  dimensionless. Status: STANDARD.
\item
  \(D_M,\ D_\mu\): The covariant derivative: on dirac16complex
  \(D_M = \partial_M + \Omega_M\) (with
  \(D_M\bar\Psi = \partial_M\bar\Psi - \bar\Psi\Omega_M\)); on the
  fields of the fold the gauge and spin covariant derivative \(D_\mu\)
  of IN-13; \(D^\mu = \hat h^{\mu\nu}D_\nu\); on an adjoint field
  \((D_\mu X)^A = \partial_\mu X^A + g_3f^{ABC}G^B_\mu X^C\) (likewise
  for \(SU(2)\) with \(g_2\epsilon^{IJK}\)). Unit: 1/length. Status:
  STANDARD.
\item
  \(\Phi\): The Higgs doublet (two complex components, hypercharge
  \(1/2\)), a field of the fold. Unit: GeV. Status: AUTHOR.
\item
  \({}^\dagger\): Hermitian conjugation: the complex conjugate of the
  transpose of a matrix or spinor (for field operators the adjoint
  operator). Status: STANDARD.
\end{itemize}

The Yang-Mills equation of the \(SU(2)_L\) fields. The sources are the
weak-isospin current of the left-handed doublets and the isospin current
of the Higgs doublet (which contains the mass terms of the \(W\) fields:
in the vacuum \(J_\Phi^{I\nu} \ni -g_2(v^2 + \mathcal{A}^2)W^{I\nu}/4\),
so that FE-17 becomes
\(\nabla_\mu W^{I\mu\nu} = m_{\mathrm{W}}^2W^{I\nu} + \dots\)). Status:
STANDARD (INVENTED portal factor).

\[
\begin{aligned}
&\frac{1}{\sqrt{-\hat h}\,\mathcal{Z}}\,\partial_\mu\Big[\sqrt{-\hat h}\,\mathcal{Z}\Big(\mathcal{K}_{B}\,B^{\mu\nu}
 + \sum_X d_X\,\mathcal{D}_X^{\mu\nu}\Big)\Big] = -g_1\,J_Y^\nu,\qquad \mathcal{K}_{B} = 1 + \sum_X c^{B}_X\,\bar\Psi P_X\Psi,\\
&J_Y^\nu = \sum_\psi Y_\psi\,\bar\psi\,\hat\gamma^\nu\psi + \frac{\mathrm{i}}{2}\big[(D^\nu\Phi)^\dagger\Phi - \Phi^\dagger D^\nu\Phi\big]
\end{aligned} \tag{FE-18}
\] \textbf{Symbols in (FE-18).}

\begin{itemize}
\tightlist
\item
  \(\hat h_{\mu\nu}\): The four-dimensional metric of the fold,
  \(\hat h_{\mu\nu} = \mathrm{diag}(\mathrm{e}^{2a_4}, \mathrm{e}^{2a_4}, \mathrm{e}^{2a_4}, -1)\)
  in the coordinates \(x_1, \dots, x_4\) (the induced metric at
  \(\sigma = 0\) without the extra times); \(\hat h^{\mu\nu}\) its
  inverse, \(\hat h\) its determinant,
  \(\sqrt{-\hat h} = \mathrm{e}^{3a_4}\); \(\hat h^{(1)}_{\mu\nu}\) its
  first-order cosmological perturbation. Unit: dimensionless. Status:
  DERIVED.
\item
  \(\mathcal{Z}\): The fold factor,
  \(\mathcal{Z}(x_4) = V_{\mathrm{t}}(\mathrm{e}^{-a_4}b_4)^3/V_*\): the
  factor by which the size of the extra times multiplies every term of
  the fold; \(\mathcal{Z} = 1\) when the extra times are frozen at their
  late value. Unit: dimensionless. Value: 1 in the late universe.
  Status: DERIVED.
\item
  \(\partial\): A partial derivative: \(\partial_M\) with respect to the
  coordinate \(X^M\), \(\partial_\mu\) with respect to \(x^\mu\),
  \(\partial_\sigma\) with respect to \(\sigma\),
  \(\partial X/\partial Y\) the derivative of \(X\) with respect to
  \(Y\) at fixed other variables. Status: STANDARD.
\item
  \(\mu\): A four-dimensional coordinate index running over the
  observable coordinates \(x_1, x_2, x_3, x_4\) of the fold (signature
  \(+,+,+,-\)); a repeated index is summed. Status: STANDARD.
\item
  \(\mathcal{K}_{G},\ \mathcal{K}_{W},\ \mathcal{K}_{B}\): The portal
  factors of the gauge kinetic terms,
  \(\mathcal{K}_{G} = 1 + \sum_Xc^{G}_X\bar\Psi P_X\Psi\) (likewise for
  \(W\) and \(B\)): the gauge couplings at the fold are divided by
  \(\sqrt{\mathcal{K}}\). Unit: dimensionless. Value: 1 where
  dirac16complex vanishes. Status: DERIVED.
\item
  \(B_\mu,\ B_{\mu\nu}\): The hypercharge gauge field of \(U(1)_Y\) on
  the fold and its field strength. Unit: GeV (field), GeV\(^2\) (field
  strength). Status: AUTHOR.
\item
  \(\nu\): A four-dimensional coordinate index (as \(\mu\)). Status:
  STANDARD.
\item
  \(X\): The label of a half of dirac16complex:
  \(X \in \{\mathrm{L}, \mathrm{R}\}\) (left-handed or right-handed
  half); a sum over \(X\) runs over both. Status: STANDARD.
\item
  \(d_X\): The dipole-portal coupling of the half \(X\) of
  dirac16complex with the hypercharge field strength. Unit:
  GeV\(^{-5}\). Status: INVENTED.
\item
  \(\mathcal{D}_X^{\mu\nu}\): The dipole (antisymmetric tensor) density
  of the half \(X\) of dirac16complex at the fold,
  \(\mathcal{D}_X^{\mu\nu} = \mathrm{i}\bar\Psi\Sigma^{\mu\nu}P_X\Psi\)
  (real, DERIVED). Unit: GeV\(^7\). Status: DERIVED.
\item
  \(g_1\): The \(U(1)_Y\) gauge coupling (usually \(g'\); not normalised
  for grand unification). Unit: dimensionless. Value:
  \(g_1(m_Z) = 0.3574\). Status: STANDARD.
\item
  \(J_Y^\nu\): The hypercharge current of the fermions and of the Higgs
  doublet. Unit: GeV\(^3\). Status: STANDARD.
\item
  \(c^{B}_X\): The \(B\)-kinetic portal coupling of the half \(X\).
  Unit: GeV\(^{-7}\). Status: INVENTED.
\item
  \(\Psi,\ \bar\Psi\): dirac16complex: the author's 16-component complex
  spinor field of eight dimensions;
  \(\bar\Psi = \Psi^\dagger\sigma_{16}\) is its Dirac adjoint; in a term
  of the fold, \(\Psi\) is its value at \(\sigma = 0\). Unit:
  GeV\(^{7/2}\). Status: AUTHOR.
\item
  \(P_X\): The chiral projector of the half \(X\):
  \(P_X = P_{\mathrm{L}}\) for \(X = \mathrm{L}\), \(P_{\mathrm{R}}\)
  for \(X = \mathrm{R}\). Unit: dimensionless. Status: AUTHOR.
\item
  \(\psi\): The label of a Standard-Model fermion field in a sum:
  \(\psi\) runs over
  \(Q, u_{\mathrm{R}}, d_{\mathrm{R}}, \ell, e_{\mathrm{R}}, \nu_{\mathrm{R}}\).
  Status: STANDARD.
\item
  \(Y_f\): The hypercharge of a field (IN-13: \(1/6\), \(2/3\),
  \(-1/3\), \(-1/2\), \(-1\), \(0\) for
  \(Q, u_{\mathrm{R}}, d_{\mathrm{R}}, \ell, e_{\mathrm{R}}, \nu_{\mathrm{R}}\),
  \(1/2\) for \(\Phi\)). Unit: dimensionless. Status: STANDARD.
\item
  \(\psi,\ \bar\psi\): A Standard-Model fermion field (a
  four-dimensional Weyl field of the fold with three generations,
  written as a four-component spinor of fixed chirality);
  \(\bar\psi = \psi^\dagger\beta\); \(\psi_f\) the fields of the kind
  \(f\). Unit: GeV\(^{3/2}\). Status: STANDARD.
\item
  \(\hat\gamma^\mu\): The four-dimensional complex \(4 \times 4\) Dirac
  matrices of the fold,
  \(\hat\gamma^\mu = \hat{\mathbf{e}}_a{}^\mu\hat\gamma^{(a)}\),
  \(\{\hat\gamma^\mu, \hat\gamma^\nu\} = -2\hat h^{\mu\nu}\mathbf{1}\);
  \(\hat\gamma_\mu = \hat h_{\mu\nu}\hat\gamma^\nu\). Unit:
  dimensionless. Status: STANDARD.
\item
  \(\mathrm{i}\): The imaginary unit, \(\mathrm{i}^2 = -1\). Unit:
  dimensionless. Status: STANDARD.
\item
  \(D_M,\ D_\mu\): The covariant derivative: on dirac16complex
  \(D_M = \partial_M + \Omega_M\) (with
  \(D_M\bar\Psi = \partial_M\bar\Psi - \bar\Psi\Omega_M\)); on the
  fields of the fold the gauge and spin covariant derivative \(D_\mu\)
  of IN-13; \(D^\mu = \hat h^{\mu\nu}D_\nu\); on an adjoint field
  \((D_\mu X)^A = \partial_\mu X^A + g_3f^{ABC}G^B_\mu X^C\) (likewise
  for \(SU(2)\) with \(g_2\epsilon^{IJK}\)). Unit: 1/length. Status:
  STANDARD.
\item
  \(\Phi\): The Higgs doublet (two complex components, hypercharge
  \(1/2\)), a field of the fold. Unit: GeV. Status: AUTHOR.
\item
  \({}^\dagger\): Hermitian conjugation: the complex conjugate of the
  transpose of a matrix or spinor (for field operators the adjoint
  operator). Status: STANDARD.
\end{itemize}

The Maxwell equation of the hypercharge field. The dipole portal adds
the dipole density of dirac16complex to the field strength, as a
polarisation adds to the electric field in a medium: dirac16complex at
the fold carries electric and magnetic dipole moments. The source is the
hypercharge current of the fermions and of the Higgs doublet. Status:
STANDARD (INVENTED portal terms).

\hypertarget{the-quarks-and-leptons-1}{%
\subsection{The quarks and leptons}\label{the-quarks-and-leptons-1}}

\[
\mathrm{i}\hat\gamma^\mu\big(D_\mu + \mathrm{i}\,\mathcal{V}^{Q}_\mu\big)Q = \mathcal{Y}^{u}\tilde\Phi\,u_{\mathrm{R}}
 + \mathcal{Y}^{d}\Phi\,d_{\mathrm{R}},\qquad \mathcal{V}^{f}_\mu = \sum_X c^{f}_X\,j_{X\mu} \tag{FE-19}
\] \textbf{Symbols in (FE-19).}

\begin{itemize}
\tightlist
\item
  \(\mathrm{i}\): The imaginary unit, \(\mathrm{i}^2 = -1\). Unit:
  dimensionless. Status: STANDARD.
\item
  \(\hat\gamma^\mu\): The four-dimensional complex \(4 \times 4\) Dirac
  matrices of the fold,
  \(\hat\gamma^\mu = \hat{\mathbf{e}}_a{}^\mu\hat\gamma^{(a)}\),
  \(\{\hat\gamma^\mu, \hat\gamma^\nu\} = -2\hat h^{\mu\nu}\mathbf{1}\);
  \(\hat\gamma_\mu = \hat h_{\mu\nu}\hat\gamma^\nu\). Unit:
  dimensionless. Status: STANDARD.
\item
  \(\mu\): A four-dimensional coordinate index running over the
  observable coordinates \(x_1, x_2, x_3, x_4\) of the fold (signature
  \(+,+,+,-\)); a repeated index is summed. Status: STANDARD.
\item
  \(D_M,\ D_\mu\): The covariant derivative: on dirac16complex
  \(D_M = \partial_M + \Omega_M\) (with
  \(D_M\bar\Psi = \partial_M\bar\Psi - \bar\Psi\Omega_M\)); on the
  fields of the fold the gauge and spin covariant derivative \(D_\mu\)
  of IN-13; \(D^\mu = \hat h^{\mu\nu}D_\nu\); on an adjoint field
  \((D_\mu X)^A = \partial_\mu X^A + g_3f^{ABC}G^B_\mu X^C\) (likewise
  for \(SU(2)\) with \(g_2\epsilon^{IJK}\)). Unit: 1/length. Status:
  STANDARD.
\item
  \(\mathcal{V}^{f}_\mu\): The vector potential that dirac16complex
  exerts on the Standard-Model fermions of kind \(f\) through the
  current portals, \(\mathcal{V}^{f}_\mu = \sum_Xc^{f}_Xj_{X\mu}\)
  (\(\mathcal{V}^{Q}, \mathcal{V}^{u}, \dots, \mathcal{V}^{\nu}\) for
  the six kinds); flavour blind. Unit: GeV. Status: DERIVED.
\item
  \(Q,\ \bar Q\): The left-handed quark doublets of the three
  generations, \(Q = (u_{\mathrm{L}}, d_{\mathrm{L}})\) (colour
  triplets, hypercharge \(1/6\)). Unit: GeV\(^{3/2}\). Status: AUTHOR.
\item
  \(\mathcal{Y}^{u},\ \mathcal{Y}^{d},\ \mathcal{Y}^{e},\ \mathcal{Y}^{\nu}\):
  The Yukawa matrices (\(3 \times 3\) complex matrices in generation
  space) of the up quarks, the down quarks, the charged leptons and the
  neutrinos. Unit: dimensionless. Status: STANDARD.
\item
  \(\tilde\Phi\): The conjugate Higgs doublet,
  \(\tilde\Phi = \mathrm{i}\tau^2\Phi^*\) (hypercharge \(-1/2\)). Unit:
  GeV. Status: STANDARD.
\item
  \(u_{\mathrm{R}}\): The right-handed up-type quarks \(u, c, t\)
  (singlets of \(SU(2)_L\), hypercharge \(2/3\)). Unit: GeV\(^{3/2}\).
  Status: STANDARD.
\item
  \(\Phi\): The Higgs doublet (two complex components, hypercharge
  \(1/2\)), a field of the fold. Unit: GeV. Status: AUTHOR.
\item
  \(d_{\mathrm{R}}\): The right-handed down-type quarks \(d, s, b\)
  (hypercharge \(-1/3\)). Unit: GeV\(^{3/2}\). Status: STANDARD.
\item
  \(X\): The label of a half of dirac16complex:
  \(X \in \{\mathrm{L}, \mathrm{R}\}\) (left-handed or right-handed
  half); a sum over \(X\) runs over both. Status: STANDARD.
\item
  \(c^{f}_X\): The current-portal coupling of the half \(X\) of
  dirac16complex with the fermion-number current of the multiplets of
  kind \(f\) (flavour blind). Unit: GeV\(^{-7}\). Status: INVENTED.
\item
  \(j_X^\mu\): The convective current of the half \(X\) of
  dirac16complex,
  \(j_X^\mu = \mathrm{i}(\bar\Psi P_XD^\mu\Psi - (D^\mu\bar\Psi)P_X\Psi)\)
  (real and joining the half with itself, DERIVED). Unit: GeV\(^8\).
  Status: DERIVED.
\item
  \(f\): The label of a kind of Standard-Model multiplet:
  \(f \in \{Q, u, d, \ell, e, \nu\}\) (quark doublets, up-quark
  singlets, down-quark singlets, lepton doublets, charged-lepton
  singlets, right-handed neutrinos); a sum over \(f\) runs over the six.
  Status: STANDARD.
\end{itemize}

The Dirac equation of the quark doublets: the variation of IN-14, IN-15
and IN-16 with respect to \(\bar Q\). The current portals act on every
Standard-Model fermion of kind \(f\) as a vector potential
\(\mathcal{V}^{f}_\mu\) built from the convective currents of the two
halves of dirac16complex at the fold; it is the same for the three
generations (flavour blind). Status: STANDARD (INVENTED potential).

\[
\mathrm{i}\hat\gamma^\mu\big(D_\mu + \mathrm{i}\,\mathcal{V}^{u}_\mu\big)u_{\mathrm{R}} = (\mathcal{Y}^{u})^\dagger\tilde\Phi^\dagger Q \tag{FE-20}
\] \textbf{Symbols in (FE-20).}

\begin{itemize}
\tightlist
\item
  \(\mathrm{i}\): The imaginary unit, \(\mathrm{i}^2 = -1\). Unit:
  dimensionless. Status: STANDARD.
\item
  \(\hat\gamma^\mu\): The four-dimensional complex \(4 \times 4\) Dirac
  matrices of the fold,
  \(\hat\gamma^\mu = \hat{\mathbf{e}}_a{}^\mu\hat\gamma^{(a)}\),
  \(\{\hat\gamma^\mu, \hat\gamma^\nu\} = -2\hat h^{\mu\nu}\mathbf{1}\);
  \(\hat\gamma_\mu = \hat h_{\mu\nu}\hat\gamma^\nu\). Unit:
  dimensionless. Status: STANDARD.
\item
  \(\mu\): A four-dimensional coordinate index running over the
  observable coordinates \(x_1, x_2, x_3, x_4\) of the fold (signature
  \(+,+,+,-\)); a repeated index is summed. Status: STANDARD.
\item
  \(D_M,\ D_\mu\): The covariant derivative: on dirac16complex
  \(D_M = \partial_M + \Omega_M\) (with
  \(D_M\bar\Psi = \partial_M\bar\Psi - \bar\Psi\Omega_M\)); on the
  fields of the fold the gauge and spin covariant derivative \(D_\mu\)
  of IN-13; \(D^\mu = \hat h^{\mu\nu}D_\nu\); on an adjoint field
  \((D_\mu X)^A = \partial_\mu X^A + g_3f^{ABC}G^B_\mu X^C\) (likewise
  for \(SU(2)\) with \(g_2\epsilon^{IJK}\)). Unit: 1/length. Status:
  STANDARD.
\item
  \(\mathcal{V}^{f}_\mu\): The vector potential that dirac16complex
  exerts on the Standard-Model fermions of kind \(f\) through the
  current portals, \(\mathcal{V}^{f}_\mu = \sum_Xc^{f}_Xj_{X\mu}\)
  (\(\mathcal{V}^{Q}, \mathcal{V}^{u}, \dots, \mathcal{V}^{\nu}\) for
  the six kinds); flavour blind. Unit: GeV. Status: DERIVED.
\item
  \(u_{\mathrm{R}}\): The right-handed up-type quarks \(u, c, t\)
  (singlets of \(SU(2)_L\), hypercharge \(2/3\)). Unit: GeV\(^{3/2}\).
  Status: STANDARD.
\item
  \(\mathcal{Y}^{u},\ \mathcal{Y}^{d},\ \mathcal{Y}^{e},\ \mathcal{Y}^{\nu}\):
  The Yukawa matrices (\(3 \times 3\) complex matrices in generation
  space) of the up quarks, the down quarks, the charged leptons and the
  neutrinos. Unit: dimensionless. Status: STANDARD.
\item
  \({}^\dagger\): Hermitian conjugation: the complex conjugate of the
  transpose of a matrix or spinor (for field operators the adjoint
  operator). Status: STANDARD.
\item
  \(\tilde\Phi\): The conjugate Higgs doublet,
  \(\tilde\Phi = \mathrm{i}\tau^2\Phi^*\) (hypercharge \(-1/2\)). Unit:
  GeV. Status: STANDARD.
\item
  \(Q,\ \bar Q\): The left-handed quark doublets of the three
  generations, \(Q = (u_{\mathrm{L}}, d_{\mathrm{L}})\) (colour
  triplets, hypercharge \(1/6\)). Unit: GeV\(^{3/2}\). Status: AUTHOR.
\end{itemize}

\[
\mathrm{i}\hat\gamma^\mu\big(D_\mu + \mathrm{i}\,\mathcal{V}^{d}_\mu\big)d_{\mathrm{R}} = (\mathcal{Y}^{d})^\dagger\Phi^\dagger Q \tag{FE-21}
\] \textbf{Symbols in (FE-21).}

\begin{itemize}
\tightlist
\item
  \(\mathrm{i}\): The imaginary unit, \(\mathrm{i}^2 = -1\). Unit:
  dimensionless. Status: STANDARD.
\item
  \(\hat\gamma^\mu\): The four-dimensional complex \(4 \times 4\) Dirac
  matrices of the fold,
  \(\hat\gamma^\mu = \hat{\mathbf{e}}_a{}^\mu\hat\gamma^{(a)}\),
  \(\{\hat\gamma^\mu, \hat\gamma^\nu\} = -2\hat h^{\mu\nu}\mathbf{1}\);
  \(\hat\gamma_\mu = \hat h_{\mu\nu}\hat\gamma^\nu\). Unit:
  dimensionless. Status: STANDARD.
\item
  \(\mu\): A four-dimensional coordinate index running over the
  observable coordinates \(x_1, x_2, x_3, x_4\) of the fold (signature
  \(+,+,+,-\)); a repeated index is summed. Status: STANDARD.
\item
  \(D_M,\ D_\mu\): The covariant derivative: on dirac16complex
  \(D_M = \partial_M + \Omega_M\) (with
  \(D_M\bar\Psi = \partial_M\bar\Psi - \bar\Psi\Omega_M\)); on the
  fields of the fold the gauge and spin covariant derivative \(D_\mu\)
  of IN-13; \(D^\mu = \hat h^{\mu\nu}D_\nu\); on an adjoint field
  \((D_\mu X)^A = \partial_\mu X^A + g_3f^{ABC}G^B_\mu X^C\) (likewise
  for \(SU(2)\) with \(g_2\epsilon^{IJK}\)). Unit: 1/length. Status:
  STANDARD.
\item
  \(\mathcal{V}^{f}_\mu\): The vector potential that dirac16complex
  exerts on the Standard-Model fermions of kind \(f\) through the
  current portals, \(\mathcal{V}^{f}_\mu = \sum_Xc^{f}_Xj_{X\mu}\)
  (\(\mathcal{V}^{Q}, \mathcal{V}^{u}, \dots, \mathcal{V}^{\nu}\) for
  the six kinds); flavour blind. Unit: GeV. Status: DERIVED.
\item
  \(d_{\mathrm{R}}\): The right-handed down-type quarks \(d, s, b\)
  (hypercharge \(-1/3\)). Unit: GeV\(^{3/2}\). Status: STANDARD.
\item
  \(\mathcal{Y}^{u},\ \mathcal{Y}^{d},\ \mathcal{Y}^{e},\ \mathcal{Y}^{\nu}\):
  The Yukawa matrices (\(3 \times 3\) complex matrices in generation
  space) of the up quarks, the down quarks, the charged leptons and the
  neutrinos. Unit: dimensionless. Status: STANDARD.
\item
  \({}^\dagger\): Hermitian conjugation: the complex conjugate of the
  transpose of a matrix or spinor (for field operators the adjoint
  operator). Status: STANDARD.
\item
  \(\Phi\): The Higgs doublet (two complex components, hypercharge
  \(1/2\)), a field of the fold. Unit: GeV. Status: AUTHOR.
\item
  \(Q,\ \bar Q\): The left-handed quark doublets of the three
  generations, \(Q = (u_{\mathrm{L}}, d_{\mathrm{L}})\) (colour
  triplets, hypercharge \(1/6\)). Unit: GeV\(^{3/2}\). Status: AUTHOR.
\end{itemize}

\[
\mathrm{i}\hat\gamma^\mu\big(D_\mu + \mathrm{i}\,\mathcal{V}^{\ell}_\mu\big)\ell = \mathcal{Y}^{e}\Phi\,e_{\mathrm{R}}
 + \mathcal{Y}^{\nu}\tilde\Phi\,\nu_{\mathrm{R}} \tag{FE-22}
\] \textbf{Symbols in (FE-22).}

\begin{itemize}
\tightlist
\item
  \(\mathrm{i}\): The imaginary unit, \(\mathrm{i}^2 = -1\). Unit:
  dimensionless. Status: STANDARD.
\item
  \(\hat\gamma^\mu\): The four-dimensional complex \(4 \times 4\) Dirac
  matrices of the fold,
  \(\hat\gamma^\mu = \hat{\mathbf{e}}_a{}^\mu\hat\gamma^{(a)}\),
  \(\{\hat\gamma^\mu, \hat\gamma^\nu\} = -2\hat h^{\mu\nu}\mathbf{1}\);
  \(\hat\gamma_\mu = \hat h_{\mu\nu}\hat\gamma^\nu\). Unit:
  dimensionless. Status: STANDARD.
\item
  \(\mu\): A four-dimensional coordinate index running over the
  observable coordinates \(x_1, x_2, x_3, x_4\) of the fold (signature
  \(+,+,+,-\)); a repeated index is summed. Status: STANDARD.
\item
  \(D_M,\ D_\mu\): The covariant derivative: on dirac16complex
  \(D_M = \partial_M + \Omega_M\) (with
  \(D_M\bar\Psi = \partial_M\bar\Psi - \bar\Psi\Omega_M\)); on the
  fields of the fold the gauge and spin covariant derivative \(D_\mu\)
  of IN-13; \(D^\mu = \hat h^{\mu\nu}D_\nu\); on an adjoint field
  \((D_\mu X)^A = \partial_\mu X^A + g_3f^{ABC}G^B_\mu X^C\) (likewise
  for \(SU(2)\) with \(g_2\epsilon^{IJK}\)). Unit: 1/length. Status:
  STANDARD.
\item
  \(\mathcal{V}^{f}_\mu\): The vector potential that dirac16complex
  exerts on the Standard-Model fermions of kind \(f\) through the
  current portals, \(\mathcal{V}^{f}_\mu = \sum_Xc^{f}_Xj_{X\mu}\)
  (\(\mathcal{V}^{Q}, \mathcal{V}^{u}, \dots, \mathcal{V}^{\nu}\) for
  the six kinds); flavour blind. Unit: GeV. Status: DERIVED.
\item
  \(\ell,\ \bar\ell\): The left-handed lepton doublets of the three
  generations, \(\ell = (\nu_{\mathrm{L}}, e_{\mathrm{L}})\)
  (hypercharge \(-1/2\)). Unit: GeV\(^{3/2}\). Status: AUTHOR.
\item
  \(\mathcal{Y}^{u},\ \mathcal{Y}^{d},\ \mathcal{Y}^{e},\ \mathcal{Y}^{\nu}\):
  The Yukawa matrices (\(3 \times 3\) complex matrices in generation
  space) of the up quarks, the down quarks, the charged leptons and the
  neutrinos. Unit: dimensionless. Status: STANDARD.
\item
  \(\Phi\): The Higgs doublet (two complex components, hypercharge
  \(1/2\)), a field of the fold. Unit: GeV. Status: AUTHOR.
\item
  \(e_{\mathrm{R}}\): The right-handed electron, muon and tau (singlets,
  hypercharge \(-1\)). Unit: GeV\(^{3/2}\). Status: AUTHOR.
\item
  \(\tilde\Phi\): The conjugate Higgs doublet,
  \(\tilde\Phi = \mathrm{i}\tau^2\Phi^*\) (hypercharge \(-1/2\)). Unit:
  GeV. Status: STANDARD.
\item
  \(\nu_{\mathrm{R}}\): The three right-handed neutrinos (gauge
  singlets). Unit: GeV\(^{3/2}\). Status: ASSUMED.
\end{itemize}

\[
\mathrm{i}\hat\gamma^\mu\big(D_\mu + \mathrm{i}\,\mathcal{V}^{e}_\mu\big)e_{\mathrm{R}} = (\mathcal{Y}^{e})^\dagger\Phi^\dagger\ell \tag{FE-23}
\] \textbf{Symbols in (FE-23).}

\begin{itemize}
\tightlist
\item
  \(\mathrm{i}\): The imaginary unit, \(\mathrm{i}^2 = -1\). Unit:
  dimensionless. Status: STANDARD.
\item
  \(\hat\gamma^\mu\): The four-dimensional complex \(4 \times 4\) Dirac
  matrices of the fold,
  \(\hat\gamma^\mu = \hat{\mathbf{e}}_a{}^\mu\hat\gamma^{(a)}\),
  \(\{\hat\gamma^\mu, \hat\gamma^\nu\} = -2\hat h^{\mu\nu}\mathbf{1}\);
  \(\hat\gamma_\mu = \hat h_{\mu\nu}\hat\gamma^\nu\). Unit:
  dimensionless. Status: STANDARD.
\item
  \(\mu\): A four-dimensional coordinate index running over the
  observable coordinates \(x_1, x_2, x_3, x_4\) of the fold (signature
  \(+,+,+,-\)); a repeated index is summed. Status: STANDARD.
\item
  \(D_M,\ D_\mu\): The covariant derivative: on dirac16complex
  \(D_M = \partial_M + \Omega_M\) (with
  \(D_M\bar\Psi = \partial_M\bar\Psi - \bar\Psi\Omega_M\)); on the
  fields of the fold the gauge and spin covariant derivative \(D_\mu\)
  of IN-13; \(D^\mu = \hat h^{\mu\nu}D_\nu\); on an adjoint field
  \((D_\mu X)^A = \partial_\mu X^A + g_3f^{ABC}G^B_\mu X^C\) (likewise
  for \(SU(2)\) with \(g_2\epsilon^{IJK}\)). Unit: 1/length. Status:
  STANDARD.
\item
  \(\mathcal{V}^{f}_\mu\): The vector potential that dirac16complex
  exerts on the Standard-Model fermions of kind \(f\) through the
  current portals, \(\mathcal{V}^{f}_\mu = \sum_Xc^{f}_Xj_{X\mu}\)
  (\(\mathcal{V}^{Q}, \mathcal{V}^{u}, \dots, \mathcal{V}^{\nu}\) for
  the six kinds); flavour blind. Unit: GeV. Status: DERIVED.
\item
  \(e_{\mathrm{R}}\): The right-handed electron, muon and tau (singlets,
  hypercharge \(-1\)). Unit: GeV\(^{3/2}\). Status: AUTHOR.
\item
  \(\mathcal{Y}^{u},\ \mathcal{Y}^{d},\ \mathcal{Y}^{e},\ \mathcal{Y}^{\nu}\):
  The Yukawa matrices (\(3 \times 3\) complex matrices in generation
  space) of the up quarks, the down quarks, the charged leptons and the
  neutrinos. Unit: dimensionless. Status: STANDARD.
\item
  \({}^\dagger\): Hermitian conjugation: the complex conjugate of the
  transpose of a matrix or spinor (for field operators the adjoint
  operator). Status: STANDARD.
\item
  \(\Phi\): The Higgs doublet (two complex components, hypercharge
  \(1/2\)), a field of the fold. Unit: GeV. Status: AUTHOR.
\item
  \(\ell,\ \bar\ell\): The left-handed lepton doublets of the three
  generations, \(\ell = (\nu_{\mathrm{L}}, e_{\mathrm{L}})\)
  (hypercharge \(-1/2\)). Unit: GeV\(^{3/2}\). Status: AUTHOR.
\end{itemize}

\[
\mathrm{i}\hat\gamma^\mu\big(D_\mu + \mathrm{i}\,\mathcal{V}^{\nu}_\mu\big)\nu_{\mathrm{R}} = (\mathcal{Y}^{\nu})^\dagger\tilde\Phi^\dagger\ell \tag{FE-24}
\] \textbf{Symbols in (FE-24).}

\begin{itemize}
\tightlist
\item
  \(\mathrm{i}\): The imaginary unit, \(\mathrm{i}^2 = -1\). Unit:
  dimensionless. Status: STANDARD.
\item
  \(\hat\gamma^\mu\): The four-dimensional complex \(4 \times 4\) Dirac
  matrices of the fold,
  \(\hat\gamma^\mu = \hat{\mathbf{e}}_a{}^\mu\hat\gamma^{(a)}\),
  \(\{\hat\gamma^\mu, \hat\gamma^\nu\} = -2\hat h^{\mu\nu}\mathbf{1}\);
  \(\hat\gamma_\mu = \hat h_{\mu\nu}\hat\gamma^\nu\). Unit:
  dimensionless. Status: STANDARD.
\item
  \(\mu\): A four-dimensional coordinate index running over the
  observable coordinates \(x_1, x_2, x_3, x_4\) of the fold (signature
  \(+,+,+,-\)); a repeated index is summed. Status: STANDARD.
\item
  \(D_M,\ D_\mu\): The covariant derivative: on dirac16complex
  \(D_M = \partial_M + \Omega_M\) (with
  \(D_M\bar\Psi = \partial_M\bar\Psi - \bar\Psi\Omega_M\)); on the
  fields of the fold the gauge and spin covariant derivative \(D_\mu\)
  of IN-13; \(D^\mu = \hat h^{\mu\nu}D_\nu\); on an adjoint field
  \((D_\mu X)^A = \partial_\mu X^A + g_3f^{ABC}G^B_\mu X^C\) (likewise
  for \(SU(2)\) with \(g_2\epsilon^{IJK}\)). Unit: 1/length. Status:
  STANDARD.
\item
  \(\mathcal{V}^{f}_\mu\): The vector potential that dirac16complex
  exerts on the Standard-Model fermions of kind \(f\) through the
  current portals, \(\mathcal{V}^{f}_\mu = \sum_Xc^{f}_Xj_{X\mu}\)
  (\(\mathcal{V}^{Q}, \mathcal{V}^{u}, \dots, \mathcal{V}^{\nu}\) for
  the six kinds); flavour blind. Unit: GeV. Status: DERIVED.
\item
  \(\nu_{\mathrm{R}}\): The three right-handed neutrinos (gauge
  singlets). Unit: GeV\(^{3/2}\). Status: ASSUMED.
\item
  \(\mathcal{Y}^{u},\ \mathcal{Y}^{d},\ \mathcal{Y}^{e},\ \mathcal{Y}^{\nu}\):
  The Yukawa matrices (\(3 \times 3\) complex matrices in generation
  space) of the up quarks, the down quarks, the charged leptons and the
  neutrinos. Unit: dimensionless. Status: STANDARD.
\item
  \({}^\dagger\): Hermitian conjugation: the complex conjugate of the
  transpose of a matrix or spinor (for field operators the adjoint
  operator). Status: STANDARD.
\item
  \(\tilde\Phi\): The conjugate Higgs doublet,
  \(\tilde\Phi = \mathrm{i}\tau^2\Phi^*\) (hypercharge \(-1/2\)). Unit:
  GeV. Status: STANDARD.
\item
  \(\ell,\ \bar\ell\): The left-handed lepton doublets of the three
  generations, \(\ell = (\nu_{\mathrm{L}}, e_{\mathrm{L}})\)
  (hypercharge \(-1/2\)). Unit: GeV\(^{3/2}\). Status: AUTHOR.
\end{itemize}

The Dirac equations of the up-quark singlets, the down-quark singlets,
the lepton doublets, the charged-lepton singlets and the right-handed
neutrinos (FE-20 to FE-24), each the variation of IN-14, IN-15 and IN-16
with respect to the conjugate field. Each fermion sees the gauge fields
(in \(D_\mu\)), the Higgs doublet (through the Yukawa matrices) and
dirac16complex (through \(\mathcal{V}^{f}_\mu\)). The right-handed
neutrinos see no gauge field; they couple to the Standard Model only
through the neutrino Yukawa matrix, which gives the neutrinos their
Dirac masses. Status: STANDARD (INVENTED potentials).

\[
\begin{aligned}
\hat\nabla_\mu J_{\ell}^\mu &= -\hat\nabla_\mu\big(J_{e}^\mu + J_{\nu}^\mu\big)
 = -\mathrm{i}\big(\bar\ell\,\mathcal{Y}^{e}\Phi\,e_{\mathrm{R}} + \bar\ell\,\mathcal{Y}^{\nu}\tilde\Phi\,\nu_{\mathrm{R}}\big) + \mathrm{h.c.},\\
\hat\nabla_\mu J_{Q}^\mu &= -\hat\nabla_\mu\big(J_{u}^\mu + J_{d}^\mu\big)
 = -\mathrm{i}\big(\bar Q\,\mathcal{Y}^{u}\tilde\Phi\,u_{\mathrm{R}} + \bar Q\,\mathcal{Y}^{d}\Phi\,d_{\mathrm{R}}\big) + \mathrm{h.c.}
\end{aligned} \tag{FE-25}
\] \textbf{Symbols in (FE-25).}

\begin{itemize}
\tightlist
\item
  \(\hat\nabla\): The covariant derivative of the fold: with the
  Levi-Civita connection of the metric of the fold (\(\hat h_{\mu\nu}\)
  for four-dimensional tensors, the induced metric of the
  seven-dimensional fold for \(\hat\nabla_a\mathcal{S}^a{}_b\)). Status:
  STANDARD.
\item
  \(\mu\): A four-dimensional coordinate index running over the
  observable coordinates \(x_1, x_2, x_3, x_4\) of the fold (signature
  \(+,+,+,-\)); a repeated index is summed. Status: STANDARD.
\item
  \(J_f^\mu\): The fermion-number current of the Standard-Model
  multiplets of kind \(f\), \(J_f^\mu = \bar\psi_f\hat\gamma^\mu\psi_f\)
  summed over the three generations and the gauge components; with other
  labels the gauge currents (FE-16 to FE-18). Unit: GeV\(^3\). Status:
  STANDARD.
\item
  \(\mathrm{i}\): The imaginary unit, \(\mathrm{i}^2 = -1\). Unit:
  dimensionless. Status: STANDARD.
\item
  \(\ell,\ \bar\ell\): The left-handed lepton doublets of the three
  generations, \(\ell = (\nu_{\mathrm{L}}, e_{\mathrm{L}})\)
  (hypercharge \(-1/2\)). Unit: GeV\(^{3/2}\). Status: AUTHOR.
\item
  \(\mathcal{Y}^{u},\ \mathcal{Y}^{d},\ \mathcal{Y}^{e},\ \mathcal{Y}^{\nu}\):
  The Yukawa matrices (\(3 \times 3\) complex matrices in generation
  space) of the up quarks, the down quarks, the charged leptons and the
  neutrinos. Unit: dimensionless. Status: STANDARD.
\item
  \(\Phi\): The Higgs doublet (two complex components, hypercharge
  \(1/2\)), a field of the fold. Unit: GeV. Status: AUTHOR.
\item
  \(e_{\mathrm{R}}\): The right-handed electron, muon and tau (singlets,
  hypercharge \(-1\)). Unit: GeV\(^{3/2}\). Status: AUTHOR.
\item
  \(\tilde\Phi\): The conjugate Higgs doublet,
  \(\tilde\Phi = \mathrm{i}\tau^2\Phi^*\) (hypercharge \(-1/2\)). Unit:
  GeV. Status: STANDARD.
\item
  \(\nu_{\mathrm{R}}\): The three right-handed neutrinos (gauge
  singlets). Unit: GeV\(^{3/2}\). Status: ASSUMED.
\item
  \(\mathrm{h.c.}\): The Hermitian conjugate of the preceding terms.
  Status: STANDARD.
\item
  \(Q,\ \bar Q\): The left-handed quark doublets of the three
  generations, \(Q = (u_{\mathrm{L}}, d_{\mathrm{L}})\) (colour
  triplets, hypercharge \(1/6\)). Unit: GeV\(^{3/2}\). Status: AUTHOR.
\item
  \(u_{\mathrm{R}}\): The right-handed up-type quarks \(u, c, t\)
  (singlets of \(SU(2)_L\), hypercharge \(2/3\)). Unit: GeV\(^{3/2}\).
  Status: STANDARD.
\item
  \(d_{\mathrm{R}}\): The right-handed down-type quarks \(d, s, b\)
  (hypercharge \(-1/3\)). Unit: GeV\(^{3/2}\). Status: STANDARD.
\end{itemize}

The divergences of the fermion-number currents of the six kinds of
multiplets (DERIVED from FE-19 to FE-24): the Yukawa couplings turn
doublets into singlets, so the current of each kind is not conserved by
itself, while the lepton number \(J_\ell + J_e + J_\nu\) and the baryon
number \((J_Q + J_u + J_d)/3\) are conserved. The divergences enter the
fold operator \(\mathbb{M}\) (FE-13). Status: DERIVED.

\hypertarget{the-ad-hoc-higgs-vacuum}{%
\subsection{The ad hoc Higgs vacuum}\label{the-ad-hoc-higgs-vacuum}}

\[
\begin{aligned}
&\mathcal{R}_{\mathcal{A}} = \mathcal{A}'' + 3\frac{b_4'}{b_4}\mathcal{A}' + \lambda\mathcal{A}^3 + c_\Phi\,\bar\Psi P_{\mathrm{L}}\Psi\,\mathcal{A}
 = \Big(4a_{40}^2 - 6a_{40}\frac{b_4'}{b_4}\Big)\mathcal{A} + \lambda\mathcal{A}^3 + c_\Phi\,\bar\Psi P_{\mathrm{L}}\Psi\,\mathcal{A},\\
&\mathcal{R}_v = \lambda\mathcal{A}^2v + c_\Phi\,\bar\Psi P_{\mathrm{L}}\Psi\,v,\qquad
 J_\Phi^{2\,x_4} = a_{40}\,v\,\mathcal{A},\qquad J_\Phi^{1\,x_4} = J_\Phi^{3\,x_4} = 0
\end{aligned} \tag{FE-26}
\] \textbf{Symbols in (FE-26).}

\begin{itemize}
\tightlist
\item
  \(\mathcal{R}_{\mathcal{A}},\ \mathcal{R}_v\): The residuals of the ad
  hoc Higgs vacuum in the upper (\(\mathcal{R}_{\mathcal{A}}\)) and
  lower (\(\mathcal{R}_v\)) components of the Higgs equation: what the
  ad hoc vacuum leaves unbalanced. Unit: GeV\(^3\). Status: DERIVED.
\item
  \(\mathcal{A}\): The author's ad hoc upper component of the Higgs
  vacuum, almightyAdHoc,
  \(\mathcal{A}(x_4) = v_{\mathcal{A}}\mathrm{e}^{-2a_{40}x_4}\). Unit:
  GeV. Status: AUTHOR.
\item
  \({}'\): The derivative of a function of the time \(x_4\) alone with
  respect to \(x_4\) (\(a_4' = \mathrm{d}a_4/\mathrm{d}x_4\); two primes
  the second derivative); on a function of \(\sigma\) alone (a profile)
  the derivative with respect to \(\sigma\). Status: STANDARD.
\item
  \(b_4\): The author's metric function of the time \(x_4\) in the
  extra-time components, \(b_4 > 0\), never identically 1 (the author's
  order of 2026-10-24); the scale factor of the extra times is
  \(\mathrm{e}^{-a_4}b_4\) times the warp. Unit: dimensionless. Status:
  AUTHOR.
\item
  \(\lambda\): The quartic coupling of the Higgs potential,
  \(\lambda > 0\). Unit: dimensionless. Value:
  \(\lambda = m_h^2/(2v^2) = 0.129\). Status: AUTHOR.
\item
  \(c_\Phi\): The Higgs-portal coupling of the left-handed half of
  dirac16complex (the right-handed half has none). Unit: GeV\(^{-5}\).
  Status: INVENTED.
\item
  \(\Psi,\ \bar\Psi\): dirac16complex: the author's 16-component complex
  spinor field of eight dimensions;
  \(\bar\Psi = \Psi^\dagger\sigma_{16}\) is its Dirac adjoint; in a term
  of the fold, \(\Psi\) is its value at \(\sigma = 0\). Unit:
  GeV\(^{7/2}\). Status: AUTHOR.
\item
  \(P_{\mathrm{L}}\): The author's left-handed chiral projector,
  \(P_{\mathrm{L}} = (\mathbf{1} - \Gamma_8)/2\). Unit: dimensionless.
  Status: AUTHOR.
\item
  \(a_{40}\): The decay constant of almightyAdHoc, \(a_{40} > 0\):
  \(\mathcal{A}\) falls by \(\mathrm{e}\) in the time \(1/(2a_{40})\).
  Unit: GeV (1/length). Status: AUTHOR.
\item
  \(v\): The vacuum expectation value of the lower component of the
  Higgs doublet, \(v^2 = -\mu^2/\lambda\). Unit: GeV. Value: 246.22 GeV.
  Status: STANDARD.
\item
  \(J_\Phi^{I\nu}\): The weak-isospin current of the Higgs doublet (it
  contains the mass terms of the \(W\) fields); \(J_\Phi^{2\,x_4}\) its
  isospin charge density with \(I = 2\). Unit: GeV\(^3\). Status:
  STANDARD.
\item
  \(x_4\): The time (\(x_4 \ge 0\) in the ad hoc vacuum); a prime on a
  function of \(x_4\) alone is \(\mathrm{d}/\mathrm{d}x_4\). Unit:
  length (= time, \(c = 1\)). Status: AUTHOR.
\end{itemize}

The author's ad hoc vacuum IN-12 inserted into the Higgs equation FE-15
(homogeneous fields, no gauge fields, the fold factor
\(\mathcal{Z} \propto b_4^3\mathrm{e}^{-3a_4}\) so that the time
derivative of \(\sqrt{-\hat h}\mathcal{Z}\) gives \(3b_4'/b_4\)). The
upper and lower components of FE-15 are not satisfied:
\(\mathcal{R}_{\mathcal{A}}\) and \(\mathcal{R}_v\) are what remains
(the residuals; FE-15 holds exactly if the source \(\Theta\) is replaced
by \(\Theta - (\mathcal{R}_{\mathcal{A}}, \mathcal{R}_v)/\sqrt2\)). The
ad hoc vacuum is an exact solution only in the limit
\(\mathcal{A} \to 0\). It also carries a weak-isospin charge density
\(J_\Phi^{2\,x_4} = a_{40}v\mathcal{A}\) (the component along \(x_4\) of
the isospin current of FE-17 with \(I = 2\)), which the \(W^2\) equation
FE-17 must balance with a background field \(W^2_{x_4}\) (DERIVED,
\texttt{um3\_fields.py}, F5: the masses of the \(W\) and \(Z\), the
direction of the photon, the charge density, the residual of the upper
component). Status: DERIVED.

\hypertarget{the-background}{%
\section{The background}\label{the-background}}

The background is the author's metric with the two functions
\(a_4(x_4)\) and \(b_4(x_4)\) of the time, on which every other field
moves; the matter fields correct it by the next terms of a power series
(FE-35). This section writes out its geometry and the exact equations
that its two functions obey, in the bulk and at the fold.

\hypertarget{the-authors-metric-and-the-proper-chart}{%
\subsection{The author's metric and the proper
chart}\label{the-authors-metric-and-the-proper-chart}}

\[
\begin{aligned}
&g_{MN} = \mathrm{diag}\Big(\cot^2z,\ \mathrm{e}^{2a_4}(\sin^2z)^{1/6},\ \mathrm{e}^{2a_4}(\sin^2z)^{1/6},
 \ \mathrm{e}^{2a_4}(\sin^2z)^{1/6},\ -1,\\
&\qquad\qquad -\mathrm{e}^{-2a_4}b_4^2(\sin^2z)^{1/6},\ -\mathrm{e}^{-2a_4}b_4^2(\sin^2z)^{1/6},
 \ -\mathrm{e}^{-2a_4}b_4^2(\sin^2z)^{1/6}\Big),\qquad z = 6Hx_0,\\
&\sigma = \mathrm{sgn}\Big(\frac{\pi}{2} - z\Big)\frac{\ln\sin z}{6H},\qquad \mathrm{d}\sigma = \lvert\cot z\rvert\,\mathrm{d}x_0,\\
&g_{MN}\,\mathrm{d}X^M\mathrm{d}X^N = \mathrm{d}\sigma^2 + \mathrm{e}^{-2H\lvert\sigma\rvert}\Big[\mathrm{e}^{2a_4}\big(\mathrm{d}x_1^2
 + \mathrm{d}x_2^2 + \mathrm{d}x_3^2\big) - \mathrm{e}^{-2a_4}b_4^2\big(\mathrm{d}x_5^2 + \mathrm{d}x_6^2 + \mathrm{d}x_7^2\big)\Big]
 - \mathrm{d}x_4^2
\end{aligned} \tag{FE-27}
\] \textbf{Symbols in (FE-27).}

\begin{itemize}
\tightlist
\item
  \(g_{MN}\): The metric tensor of the eight-dimensional spacetime
  (signature \(+,+,+,+,-,-,-,-\) in the order \(x_0, \dots, x_7\)):
  \(g_{MN}\,\mathrm{d}X^M\mathrm{d}X^N\) is the squared interval,
  \(g^{MN}\) its inverse, \(\lvert g\rvert\) the absolute value of its
  determinant (\(\sqrt{\lvert g\rvert}\,\mathrm{d}^8X\) is the proper
  volume element). Unit: dimensionless. Status: AUTHOR.
\item
  \(M\): A coordinate index of eight dimensions (also with a numbered
  copy \(M_1, M_2, \dots\)): it runs over \(x_0, x_1, \dots, x_7\); a
  repeated index is summed. Status: STANDARD.
\item
  \(N\): A coordinate index of eight dimensions (also
  \(N_1, N_2, \dots\)), running over \(x_0, \dots, x_7\); a repeated
  index is summed. Status: STANDARD.
\item
  \(\mathrm{diag}\): A diagonal matrix with the listed entries on its
  diagonal (in the order of the coordinates \(x_0, x_1, \dots, x_7\) for
  a tensor of eight dimensions, \(x_1, \dots, x_7\) for a tensor of the
  fold, \(x_1, \dots, x_4\) for a four-dimensional tensor). Status:
  STANDARD.
\item
  \(z\): The angle variable of the hidden direction, \(z = 6Hx_0\); one
  period of \(x_0\) is \(0 < z < \pi\), the fold is \(z = \pi/2\). Unit:
  dimensionless. Status: AUTHOR.
\item
  \(\mathrm{e}\): Euler's number, the base of the natural logarithm;
  \(\mathrm{e}^{X}\) is the exponential of \(X\). Unit: dimensionless.
  Value: 2.718281828. Status: STANDARD.
\item
  \(a_4\): The author's metric function of the time \(x_4\) that scales
  ordinary 3-space: on the fold \(\mathrm{e}^{a_4}\) is the scale factor
  of cosmology and \(a_4'\) the Hubble rate. Unit: dimensionless.
  Status: AUTHOR.
\item
  \(b_4\): The author's metric function of the time \(x_4\) in the
  extra-time components, \(b_4 > 0\), never identically 1 (the author's
  order of 2026-10-24); the scale factor of the extra times is
  \(\mathrm{e}^{-a_4}b_4\) times the warp. Unit: dimensionless. Status:
  AUTHOR.
\item
  \(H\): The author's constant (an inverse length) that fixes the period
  \(\pi/(6H)\) of \(x_0\), the Lovelock couplings and the cosmological
  constant \(\Lambda = -10H^2\); in UM-3 a free parameter. Unit: GeV
  (1/length). Status: AUTHOR.
\item
  \(x_0\): The hidden space-like coordinate (listed first); the metric
  depends on it through \(z = 6Hx_0\), periodically with the period
  \(\pi/(6H)\). Unit: length. Status: AUTHOR.
\item
  \(\sigma\): The proper coordinate of the hidden direction:
  \(\mathrm{d}\sigma = \lvert\cot z\rvert\,\mathrm{d}x_0\),
  \(\sigma = \mathrm{sgn}(\pi/2 - z)\ln(\sin z)/(6H)\) on one period
  \(0 < z < \pi\); \(\sigma = 0\) is the fold \(z = \pi/2\),
  \(\sigma \to \mp\infty\) the tips \(z \to 0, \pi\); one period of
  \(x_0\) is the whole line \(-\infty < \sigma < \infty\). Unit: length.
  Status: DERIVED.
\item
  \(\mathrm{sgn}\): The sign function: \(\mathrm{sgn}(y) = +1\) for
  \(y > 0\), \(-1\) for \(y < 0\). Status: STANDARD.
\item
  \(\pi\): The number pi, the ratio of the circumference of a circle to
  its diameter. Unit: dimensionless. Value: 3.141592654. Status:
  STANDARD.
\item
  \(\mathrm{d}\): The differential: \(\mathrm{d}X\) is an infinitesimal
  change of \(X\), \(\mathrm{d}X/\mathrm{d}Y\) the derivative of \(X\)
  with respect to \(Y\), \(\mathrm{d}^nX\) the volume element of \(n\)
  variables. Status: STANDARD.
\item
  \(X = \{x_0, \dots, x_7\}\): The author's coordinates of the
  eight-dimensional spacetime, \(X = \{x_0, x_1, \dots, x_7\}\) with the
  hidden direction first (Dirac\_Einstein-Lovelock\_4+4\_Nash.nb,
  section DEFCHART); \(X^M\) is the coordinate with the index \(M\) and
  \(\mathrm{d}^8X = \mathrm{d}x_0\,\mathrm{d}x_1\cdots\mathrm{d}x_7\)
  the coordinate volume element. Unit: length (each coordinate). Status:
  AUTHOR.
\item
  \(x_1\): The first coordinate of ordinary 3-space. Unit: length.
  Status: AUTHOR.
\item
  \(x_2\): The second coordinate of ordinary 3-space. Unit: length.
  Status: AUTHOR.
\item
  \(x_3\): The third coordinate of ordinary 3-space. Unit: length.
  Status: AUTHOR.
\item
  \(x_4\): The time (\(x_4 \ge 0\) in the ad hoc vacuum); a prime on a
  function of \(x_4\) alone is \(\mathrm{d}/\mathrm{d}x_4\). Unit:
  length (= time, \(c = 1\)). Status: AUTHOR.
\item
  \(x_5\): The first of the three extra time-like coordinates. Unit:
  length. Status: AUTHOR.
\item
  \(x_6\): The second of the three extra time-like coordinates. Unit:
  length. Status: AUTHOR.
\item
  \(x_7\): The third of the three extra time-like coordinates. Unit:
  length. Status: AUTHOR.
\end{itemize}

The author's metric of the primordial gravitational field (his order of
2026-10-24: the metric of
\texttt{Dirac\_Einstein-Lovelock\_4+4\_Nash.nb} with \(b_4 \to b_4^2\);
\((\sin^2z)^{1/6}\) is the positive branch of
\texttt{Sin{[}6\ H\ x0{]}\^{}(1/3)}, the author's rule of 2026-10-16),
and the same metric in the proper chart \(\sigma\) of the hidden
direction. In the author's chart the metric is smooth but degenerate at
the fold \(z = \pi/2\) (\(\cot z = 0\)); the proper distance from the
fold is \(\lvert\sigma\rvert\), and \(\sigma\) runs from \(-\infty\)
(the tip \(z \to 0\)) through \(0\) (the fold) to \(+\infty\) (the tip
\(z \to \pi\)): one period of \(x_0\) is the whole line, and the two
halves of a period are mirror images under \(z \to \pi - z\)
(\(\sigma \to -\sigma\)). Ordinary space and the extra times are warped
by \(\mathrm{e}^{-H\lvert\sigma\rvert}\) (with a kink at the fold), the
time \(x_4\) is not warped. \(a_4\) and \(b_4\) are functions of
\(x_4\); \(b_4\) is never identically 1. Status: AUTHOR (the metric),
DERIVED (the chart, \texttt{um3\_fold.py}, J1).

\[
\begin{aligned}
&\mathbf{e}^{(0)} = \mathrm{d}\sigma,\quad \mathbf{e}^{(i)} = \mathrm{e}^{a_4 - H\lvert\sigma\rvert}\mathrm{d}x_i,\quad
 \mathbf{e}^{(4)} = \mathrm{d}x_4,\quad \mathbf{e}^{(r)} = b_4\,\mathrm{e}^{-a_4 - H\lvert\sigma\rvert}\mathrm{d}x_r,\\
&\sqrt{\lvert g\rvert} = b_4^3\,\mathrm{e}^{-6H\lvert\sigma\rvert},\qquad
 \mathsf{A} = \frac{a_4'}{H},\qquad \mathsf{U} = \frac{b_4'}{Hb_4},\qquad \mathsf{A}_1 = \frac{\mathsf{A}'}{H},\qquad
 \mathsf{U}_1 = \frac{\mathsf{U}'}{H},\\
&R^{(0)(i)}{}_{(0)(i)} = R^{(0)(r)}{}_{(0)(r)} = -H^2,\qquad R^{(0)(4)}{}_{(0)(4)} = 0,\qquad
 R^{(i)(j)}{}_{(i)(j)} = (\mathsf{A}^2 - 1)H^2,\\
&R^{(i)(4)}{}_{(i)(4)} = (\mathsf{A}^2 + \mathsf{A}_1)H^2,\qquad
 R^{(i)(r)}{}_{(i)(r)} = -(1 + \mathsf{A}^2 - \mathsf{A}\mathsf{U})H^2,\\
&R^{(4)(r)}{}_{(4)(r)} = \big((\mathsf{U} - \mathsf{A})^2 + \mathsf{U}_1 - \mathsf{A}_1\big)H^2,\qquad
 R^{(r)(s)}{}_{(r)(s)} = \big((\mathsf{U} - \mathsf{A})^2 - 1\big)H^2,\\
&R^{(0)(i)}{}_{(4)(i)} = \mathrm{sgn}(\sigma)\,\mathsf{A}H^2,\qquad
 R^{(0)(r)}{}_{(4)(r)} = \mathrm{sgn}(\sigma)\,(\mathsf{U} - \mathsf{A})H^2
\end{aligned} \tag{FE-28}
\] \textbf{Symbols in (FE-28).}

\begin{itemize}
\tightlist
\item
  \(\mathbf{e}_a{}^M,\ \mathbf{e}^a{}_M\): The vielbein (frame field) of
  the metric: \(g_{MN} = \eta_{ab}\,\mathbf{e}^a{}_M\mathbf{e}^b{}_N\);
  \(\mathbf{e}_a{}^M\) is its inverse;
  \(\mathbf{e}_{bN} = \eta_{bc}\mathbf{e}^c{}_N\); \(\mathbf{e}^{(a)}\)
  is the frame one-form \(\mathbf{e}^a{}_M\,\mathrm{d}X^M\). Unit:
  dimensionless. Status: STANDARD.
\item
  \(\mathrm{d}\): The differential: \(\mathrm{d}X\) is an infinitesimal
  change of \(X\), \(\mathrm{d}X/\mathrm{d}Y\) the derivative of \(X\)
  with respect to \(Y\), \(\mathrm{d}^nX\) the volume element of \(n\)
  variables. Status: STANDARD.
\item
  \(\sigma\): The proper coordinate of the hidden direction:
  \(\mathrm{d}\sigma = \lvert\cot z\rvert\,\mathrm{d}x_0\),
  \(\sigma = \mathrm{sgn}(\pi/2 - z)\ln(\sin z)/(6H)\) on one period
  \(0 < z < \pi\); \(\sigma = 0\) is the fold \(z = \pi/2\),
  \(\sigma \to \mp\infty\) the tips \(z \to 0, \pi\); one period of
  \(x_0\) is the whole line \(-\infty < \sigma < \infty\). Unit: length.
  Status: DERIVED.
\item
  \(\mathrm{e}\): Euler's number, the base of the natural logarithm;
  \(\mathrm{e}^{X}\) is the exponential of \(X\). Unit: dimensionless.
  Value: 2.718281828. Status: STANDARD.
\item
  \(a_4\): The author's metric function of the time \(x_4\) that scales
  ordinary 3-space: on the fold \(\mathrm{e}^{a_4}\) is the scale factor
  of cosmology and \(a_4'\) the Hubble rate. Unit: dimensionless.
  Status: AUTHOR.
\item
  \(H\): The author's constant (an inverse length) that fixes the period
  \(\pi/(6H)\) of \(x_0\), the Lovelock couplings and the cosmological
  constant \(\Lambda = -10H^2\); in UM-3 a free parameter. Unit: GeV
  (1/length). Status: AUTHOR.
\item
  \(x\): The coordinates of the fold: \(x_i\) (\(i = 1, 2, 3\)) of
  ordinary space, \(x_4\) the time, \(x_r\) (\(r = 5, 6, 7\)) the extra
  times; a field of the fold is a function of the four observable
  coordinates \(x = (x_1, x_2, x_3, x_4)\), and
  \(\mathrm{d}^4x = \mathrm{d}x_1\,\mathrm{d}x_2\,\mathrm{d}x_3\,\mathrm{d}x_4\).
  Unit: length (each coordinate). Status: AUTHOR.
\item
  \(i\): A generation index (\(i = 1, 2, 3\)), which also labels the
  neutrino mass states \(\nu_1, \nu_2, \nu_3\); a repeated index is not
  summed unless a sum sign is written. Status: STANDARD.
\item
  \(x_4\): The time (\(x_4 \ge 0\) in the ad hoc vacuum); a prime on a
  function of \(x_4\) alone is \(\mathrm{d}/\mathrm{d}x_4\). Unit:
  length (= time, \(c = 1\)). Status: AUTHOR.
\item
  \(b_4\): The author's metric function of the time \(x_4\) in the
  extra-time components, \(b_4 > 0\), never identically 1 (the author's
  order of 2026-10-24); the scale factor of the extra times is
  \(\mathrm{e}^{-a_4}b_4\) times the warp. Unit: dimensionless. Status:
  AUTHOR.
\item
  \(r\): An extra-time direction of the fold: \(r \in \{5, 6, 7\}\)
  (\(x_5, x_6, x_7\)); not summed (the three components are equal).
  Status: STANDARD.
\item
  \(g_{MN}\): The metric tensor of the eight-dimensional spacetime
  (signature \(+,+,+,+,-,-,-,-\) in the order \(x_0, \dots, x_7\)):
  \(g_{MN}\,\mathrm{d}X^M\mathrm{d}X^N\) is the squared interval,
  \(g^{MN}\) its inverse, \(\lvert g\rvert\) the absolute value of its
  determinant (\(\sqrt{\lvert g\rvert}\,\mathrm{d}^8X\) is the proper
  volume element). Unit: dimensionless. Status: AUTHOR.
\item
  \(\mathsf{A}\): The expansion rate of ordinary space in units of
  \(H\), \(\mathsf{A} = a_4'/H\). Unit: dimensionless. Value: 1 at the
  author's vacuum; \(H_0/H\) today (\(H_0\) the Hubble constant).
  Status: DERIVED.
\item
  \(\mathsf{U}\): The rate of change of \(b_4\) in units of \(H\),
  \(\mathsf{U} = b_4'/(Hb_4)\). Unit: dimensionless. Value:
  \(\mathsf{U} = \mathsf{A}\) with frozen extra times. Status: DERIVED.
\item
  \(\mathsf{A}_1\): The time derivative of \(\mathsf{A}\) in units of
  \(H\), \(\mathsf{A}_1 = \mathsf{A}'/H = a_4''/H^2\). Unit:
  dimensionless. Status: DERIVED.
\item
  \(\mathsf{U}_1\): The time derivative of \(\mathsf{U}\) in units of
  \(H\), \(\mathsf{U}_1 = \mathsf{U}'/H\). Unit: dimensionless. Status:
  DERIVED.
\item
  \({}'\): The derivative of a function of the time \(x_4\) alone with
  respect to \(x_4\) (\(a_4' = \mathrm{d}a_4/\mathrm{d}x_4\); two primes
  the second derivative); on a function of \(\sigma\) alone (a profile)
  the derivative with respect to \(\sigma\). Status: STANDARD.
\item
  \(R^{MN}{}_{PQ},\ R_{MN},\ R\): The curvature of the metric
  \(g_{MN}\): with four indices the Riemann tensor (convention
  \([\nabla_P, \nabla_Q]V^M = R^M{}_{NPQ}V^N\)), with two indices the
  Ricci tensor \(R_{MN} = R^P{}_{MPN}\), without indices the scalar
  curvature \(R = g^{MN}R_{MN}\); with frame indices
  \(R^{(a)(b)}{}_{(c)(d)}\) its frame components,
  \(R^{(a)(b)}{}_{(a)(b)}\) (no sum) the sectional curvature of the
  frame plane \((a, b)\). Unit: 1/length\(^2\). Status: STANDARD.
\item
  \(j\): A second generation or mass-state index (as \(i\)). Status:
  STANDARD.
\item
  \(s\): A second extra-time direction of the fold (as \(r\), different
  from it). Status: STANDARD.
\item
  \(\mathrm{sgn}\): The sign function: \(\mathrm{sgn}(y) = +1\) for
  \(y > 0\), \(-1\) for \(y < 0\). Status: STANDARD.
\end{itemize}

The vielbein and the volume element (\(i, j\) any two different
directions of ordinary space, \(r, s\) any two different extra times;
\(x_i\) and \(x_r\) the coordinates of these directions), the four
dimensionless rates of the background (\(\mathsf{A}\) the expansion rate
of ordinary space in units of \(H\), \(\mathsf{U}\) the rate of change
of \(b_4\)), and every non-zero frame component of the Riemann tensor,
up to the symmetries
\(R^{ab}{}_{cd} = -R^{ba}{}_{cd} = -R^{ab}{}_{dc} = R_{cd}{}^{ab}\) (the
last with the frame metric). No component depends on \(\sigma\) except
through the sign of the two mixed ones, which change sign at the fold:
the bulk on each side of the fold is homogeneous along \(\sigma\)
(DERIVED, \texttt{um3\_geometry.py}, G2, exact; \texttt{um3\_fold.py},
J2). The frame curvature of the hidden direction is \(-H^2\): each half
of a period is a space of constant negative curvature along \(\sigma\),
like a slice of anti-de Sitter space, with the time not warped. Status:
DERIVED.

\hypertarget{the-bulk-equations-of-the-background}{%
\subsection{The bulk equations of the
background}\label{the-bulk-equations-of-the-background}}

\[
\begin{aligned}
\frac{\kappa\,\mathcal{T}_{\mathrm{bulk}}{}^{(0)}{}_{(0)}}{H^2} = -\frac{1}{16}\Big(&\mathsf{A}^6 - 3\mathsf{A}^5\mathsf{U} + 12\mathsf{A}^4\mathsf{U}^2
 + 3\mathsf{A}^4\mathsf{U}_1 + 3\mathsf{A}^4 + 12\mathsf{A}^3\mathsf{A}_1\mathsf{U} - 19\mathsf{A}^3\mathsf{U}^3 - 12\mathsf{A}^3\mathsf{U}\mathsf{U}_1
 - 6\mathsf{A}^3\mathsf{U} - 18\mathsf{A}^2\mathsf{A}_1\mathsf{U}^2 + 9\mathsf{A}^2\mathsf{U}^4\\
&+ 9\mathsf{A}^2\mathsf{U}^2\mathsf{U}_1 - 18\mathsf{A}^2\mathsf{U}^2 - 6\mathsf{A}^2\mathsf{U}_1 - 9\mathsf{A}^2 + 6\mathsf{A}\mathsf{A}_1\mathsf{U}^3
 - 12\mathsf{A}\mathsf{A}_1\mathsf{U} + 21\mathsf{A}\mathsf{U}^3 + 12\mathsf{A}\mathsf{U}\mathsf{U}_1 + 9\mathsf{A}\mathsf{U}\\
&+ 6\mathsf{A}_1\mathsf{U}^2 + 3\mathsf{U}^4 + 3\mathsf{U}^2\mathsf{U}_1 + 18\mathsf{U}^2 + 3\mathsf{U}_1 + 5\Big),\\
\frac{\kappa\,\mathcal{T}_{\mathrm{bulk}}{}^{(i)}{}_{(i)}}{H^2} = -\frac{1}{16}\Big(&\mathsf{A}^6 - 6\mathsf{A}^5\mathsf{U} - 5\mathsf{A}^4\mathsf{A}_1
 + 12\mathsf{A}^4\mathsf{U}^2 + 3\mathsf{A}^4\mathsf{U}_1 + 3\mathsf{A}^4 + 12\mathsf{A}^3\mathsf{A}_1\mathsf{U} - 10\mathsf{A}^3\mathsf{U}^3
 - 6\mathsf{A}^3\mathsf{U}\mathsf{U}_1 + 12\mathsf{A}^3\mathsf{U} - 9\mathsf{A}^2\mathsf{A}_1\mathsf{U}^2\\
&+ 6\mathsf{A}^2\mathsf{A}_1 + 3\mathsf{A}^2\mathsf{U}^4 + 3\mathsf{A}^2\mathsf{U}^2\mathsf{U}_1 - 36\mathsf{A}^2\mathsf{U}^2 - 6\mathsf{A}^2\mathsf{U}_1
 - 9\mathsf{A}^2 + 2\mathsf{A}\mathsf{A}_1\mathsf{U}^3 - 12\mathsf{A}\mathsf{A}_1\mathsf{U} + 18\mathsf{A}\mathsf{U}^3 + 6\mathsf{A}\mathsf{U}\mathsf{U}_1\\
&- 6\mathsf{A}\mathsf{U} + 3\mathsf{A}_1\mathsf{U}^2 - \mathsf{A}_1 + 3\mathsf{U}^4 + 3\mathsf{U}^2\mathsf{U}_1 + 24\mathsf{U}^2 + 3\mathsf{U}_1 + 5\Big),\\
\frac{\kappa\,\mathcal{T}_{\mathrm{bulk}}{}^{(4)}{}_{(4)}}{H^2} = \frac{1}{16}\Big(&5\mathsf{A}^6 - 15\mathsf{A}^5\mathsf{U} + 15\mathsf{A}^4\mathsf{U}^2
 - 9\mathsf{A}^4 - 5\mathsf{A}^3\mathsf{U}^3 + 18\mathsf{A}^3\mathsf{U} + 3\mathsf{A}^2 - 9\mathsf{A}\mathsf{U}^3 - 3\mathsf{A}\mathsf{U} - 3\mathsf{U}^2 + 1\Big),\\
\frac{\kappa\,\mathcal{T}_{\mathrm{bulk}}{}^{(r)}{}_{(r)}}{H^2} = -\frac{1}{16}\Big(&\mathsf{A}^6 + 5\mathsf{A}^4\mathsf{A}_1 - 3\mathsf{A}^4\mathsf{U}^2
 - 2\mathsf{A}^4\mathsf{U}_1 + 3\mathsf{A}^4 - 8\mathsf{A}^3\mathsf{A}_1\mathsf{U} + 2\mathsf{A}^3\mathsf{U}^3 + 2\mathsf{A}^3\mathsf{U}\mathsf{U}_1
 - 24\mathsf{A}^3\mathsf{U} + 3\mathsf{A}^2\mathsf{A}_1\mathsf{U}^2 - 6\mathsf{A}^2\mathsf{A}_1\\
&+ 18\mathsf{A}^2\mathsf{U}^2 - 9\mathsf{A}^2 + 6\mathsf{A}\mathsf{U}^3 + 6\mathsf{A}\mathsf{U}\mathsf{U}_1 + 24\mathsf{A}\mathsf{U} + 3\mathsf{A}_1\mathsf{U}^2
 + \mathsf{A}_1 + 9\mathsf{U}^2 + 2\mathsf{U}_1 + 5\Big),\\
\frac{\kappa\,\mathcal{T}_{\mathrm{bulk}}{}^{(0)}{}_{(4)}}{H^2} = \mathrm{sgn}(\sigma)\,\frac{3}{16}\,\mathsf{U}\,&\mathcal{G},\qquad
 \mathcal{G} = 5\mathsf{A}^2(\mathsf{A} - \mathsf{U})^2 - 6\mathsf{A}^2 + 6\mathsf{A}\mathsf{U} + \mathsf{U}^2 + 1
\end{aligned} \tag{FE-29}
\] \textbf{Symbols in (FE-29).}

\begin{itemize}
\tightlist
\item
  \(\kappa\): The eight-dimensional gravitational coupling,
  \(\kappa = 8\pi G_8\). Unit: GeV\(^{-6}\). Status: AUTHOR.
\item
  \(\mathcal{T}_{\mathrm{bulk}}\): The energy-momentum tensor of all
  sources of the bulk,
  \(\mathcal{T}_\Psi + \mathcal{T}_{\mathrm{stab}}\);
  \(\mathcal{T}_{\mathrm{bulk}}{}^{(0)}{}_b\) is its flux along
  \(\sigma\) (frame components). Unit: GeV\(^8\). Status: STANDARD.
\item
  \(i\): A generation index (\(i = 1, 2, 3\)), which also labels the
  neutrino mass states \(\nu_1, \nu_2, \nu_3\); a repeated index is not
  summed unless a sum sign is written. Status: STANDARD.
\item
  \(r\): An extra-time direction of the fold: \(r \in \{5, 6, 7\}\)
  (\(x_5, x_6, x_7\)); not summed (the three components are equal).
  Status: STANDARD.
\item
  \(H\): The author's constant (an inverse length) that fixes the period
  \(\pi/(6H)\) of \(x_0\), the Lovelock couplings and the cosmological
  constant \(\Lambda = -10H^2\); in UM-3 a free parameter. Unit: GeV
  (1/length). Status: AUTHOR.
\item
  \(\mathsf{A}\): The expansion rate of ordinary space in units of
  \(H\), \(\mathsf{A} = a_4'/H\). Unit: dimensionless. Value: 1 at the
  author's vacuum; \(H_0/H\) today (\(H_0\) the Hubble constant).
  Status: DERIVED.
\item
  \(\mathsf{U}\): The rate of change of \(b_4\) in units of \(H\),
  \(\mathsf{U} = b_4'/(Hb_4)\). Unit: dimensionless. Value:
  \(\mathsf{U} = \mathsf{A}\) with frozen extra times. Status: DERIVED.
\item
  \(\mathsf{A}_1\): The time derivative of \(\mathsf{A}\) in units of
  \(H\), \(\mathsf{A}_1 = \mathsf{A}'/H = a_4''/H^2\). Unit:
  dimensionless. Status: DERIVED.
\item
  \(\mathsf{U}_1\): The time derivative of \(\mathsf{U}\) in units of
  \(H\), \(\mathsf{U}_1 = \mathsf{U}'/H\). Unit: dimensionless. Status:
  DERIVED.
\item
  \(\mathrm{sgn}\): The sign function: \(\mathrm{sgn}(y) = +1\) for
  \(y > 0\), \(-1\) for \(y < 0\). Status: STANDARD.
\item
  \(\sigma\): The proper coordinate of the hidden direction:
  \(\mathrm{d}\sigma = \lvert\cot z\rvert\,\mathrm{d}x_0\),
  \(\sigma = \mathrm{sgn}(\pi/2 - z)\ln(\sin z)/(6H)\) on one period
  \(0 < z < \pi\); \(\sigma = 0\) is the fold \(z = \pi/2\),
  \(\sigma \to \mp\infty\) the tips \(z \to 0, \pi\); one period of
  \(x_0\) is the whole line \(-\infty < \sigma < \infty\). Unit: length.
  Status: DERIVED.
\item
  \(\mathcal{G}\): The factor of the \((\sigma, x_4)\) bulk equation of
  the author's metric,
  \(\mathcal{G} = 5\mathsf{A}^2(\mathsf{A} - \mathsf{U})^2 - 6\mathsf{A}^2 + 6\mathsf{A}\mathsf{U} + \mathsf{U}^2 + 1\).
  Unit: dimensionless. Status: DERIVED.
\end{itemize}

The bulk equations FE-1 for the author's metric, divided by \(H^2\): the
left-hand side of FE-1 with the author's weights and
\(\Lambda = -10H^2\), evaluated exactly (DERIVED,
\texttt{um3\_geometry.py}, G3 and G4, with the identities FE-3 checked;
\texttt{um3\_fold.py}, J2, by an independent general computation). They
state what the bulk sources must be for given \(a_4(x_4)\),
\(b_4(x_4)\): \(\mathcal{T}^{(4)}{}_{(4)} = -\rho\) gives the energy
density, \(\mathcal{T}^{(0)}{}_{(0)}\), \(\mathcal{T}^{(i)}{}_{(i)}\),
\(\mathcal{T}^{(r)}{}_{(r)}\) the pressures along \(\sigma\), ordinary
space and the extra times, and \(\mathcal{T}^{(0)}{}_{(4)}\) the energy
flux along \(\sigma\), which has opposite signs on the two sides of the
fold (it flows toward or away from the fold on both sides). Because the
frame components do not depend on \(\sigma\), the author's metric can be
exact only with sources whose frame components do not depend on
\(\sigma\); sources that depend on \(\sigma\) require the next terms of
the power series of the metric (FE-35). The time-time equation contains
no \(\mathsf{A}_1\) and no \(\mathsf{U}_1\): it is the energy constraint
of the background. The flux equation factorises: with a source without
flux along \(\sigma\), either \(\mathsf{U} = 0\) (\(b_4\) constant) or
\(\mathcal{G}(\mathsf{A}, \mathsf{U}) = 0\), which has real roots
\(\mathsf{U}\) only for \(0.3249 \le \lvert\mathsf{A}\rvert \le 1.3764\)
(\(5\mathsf{A}^4 - 10\mathsf{A}^2 + 1 \le 0\)). Status: DERIVED.

\hypertarget{the-fold-equations-of-the-background}{%
\subsection{The fold equations of the
background}\label{the-fold-equations-of-the-background}}

\[
\begin{aligned}
\frac{\kappa\,\mathcal{S}^{(i)}{}_{(i)}}{H} = \frac{1}{24}\Big(&9\mathsf{A}^4 - 24\mathsf{A}^3\mathsf{U} - 12\mathsf{A}^2\mathsf{A}_1
 + 27\mathsf{A}^2\mathsf{U}^2 + 12\mathsf{A}^2\mathsf{U}_1 - 66\mathsf{A}^2 + 24\mathsf{A}\mathsf{A}_1\mathsf{U} - 6\mathsf{A}\mathsf{U}^3
 - 12\mathsf{A}\mathsf{U}\mathsf{U}_1 + 120\mathsf{A}\mathsf{U}\\
&- 6\mathsf{A}_1\mathsf{U}^2 + 20\mathsf{A}_1 - 6\mathsf{U}^4 - 6\mathsf{U}^2\mathsf{U}_1 - 111\mathsf{U}^2 - 60\mathsf{U}_1 - 183\Big),\\
\frac{\kappa\,\mathcal{S}^{(4)}{}_{(4)}}{H} = -\frac14\Big(&3\mathsf{A}^4 - 6\mathsf{A}^3\mathsf{U} - 10\mathsf{A}^2 + 3\mathsf{A}\mathsf{U}^3
 + 10\mathsf{A}\mathsf{U} + 10\mathsf{U}^2 + 31\Big),\\
\frac{\kappa\,\mathcal{S}^{(r)}{}_{(r)}}{H} = \frac{1}{24}\Big(&9\mathsf{A}^4 - 12\mathsf{A}^3\mathsf{U} + 12\mathsf{A}^2\mathsf{A}_1
 + 9\mathsf{A}^2\mathsf{U}^2 - 66\mathsf{A}^2 - 12\mathsf{A}\mathsf{U}^3 - 12\mathsf{A}\mathsf{U}\mathsf{U}_1 + 12\mathsf{A}\mathsf{U}
 - 6\mathsf{A}_1\mathsf{U}^2 - 20\mathsf{A}_1\\
&- 57\mathsf{U}^2 - 40\mathsf{U}_1 - 183\Big)
\end{aligned} \tag{FE-30}
\] \textbf{Symbols in (FE-30).}

\begin{itemize}
\tightlist
\item
  \(\kappa\): The eight-dimensional gravitational coupling,
  \(\kappa = 8\pi G_8\). Unit: GeV\(^{-6}\). Status: AUTHOR.
\item
  \(\mathcal{S}^a{}_b\): The thin source of the fold: its
  energy-momentum tensor per unit proper length across the fold
  (eight-dimensional normalisation, frame components along the fold);
  \(\kappa\mathcal{S}\) is the jump of the potential \(\mathfrak{F}\).
  Unit: GeV\(^7\). Status: DERIVED.
\item
  \(i\): A generation index (\(i = 1, 2, 3\)), which also labels the
  neutrino mass states \(\nu_1, \nu_2, \nu_3\); a repeated index is not
  summed unless a sum sign is written. Status: STANDARD.
\item
  \(r\): An extra-time direction of the fold: \(r \in \{5, 6, 7\}\)
  (\(x_5, x_6, x_7\)); not summed (the three components are equal).
  Status: STANDARD.
\item
  \(H\): The author's constant (an inverse length) that fixes the period
  \(\pi/(6H)\) of \(x_0\), the Lovelock couplings and the cosmological
  constant \(\Lambda = -10H^2\); in UM-3 a free parameter. Unit: GeV
  (1/length). Status: AUTHOR.
\item
  \(\mathsf{A}\): The expansion rate of ordinary space in units of
  \(H\), \(\mathsf{A} = a_4'/H\). Unit: dimensionless. Value: 1 at the
  author's vacuum; \(H_0/H\) today (\(H_0\) the Hubble constant).
  Status: DERIVED.
\item
  \(\mathsf{U}\): The rate of change of \(b_4\) in units of \(H\),
  \(\mathsf{U} = b_4'/(Hb_4)\). Unit: dimensionless. Value:
  \(\mathsf{U} = \mathsf{A}\) with frozen extra times. Status: DERIVED.
\item
  \(\mathsf{A}_1\): The time derivative of \(\mathsf{A}\) in units of
  \(H\), \(\mathsf{A}_1 = \mathsf{A}'/H = a_4''/H^2\). Unit:
  dimensionless. Status: DERIVED.
\item
  \(\mathsf{U}_1\): The time derivative of \(\mathsf{U}\) in units of
  \(H\), \(\mathsf{U}_1 = \mathsf{U}'/H\). Unit: dimensionless. Status:
  DERIVED.
\end{itemize}

The junction conditions FE-4 for the author's metric: the thin source
that the fold must carry for given \(a_4(x_4)\), \(b_4(x_4)\) (the
off-diagonal and \((\sigma, \sigma)\) components have no thin part).
They follow from the potential \(\mathfrak{F}\) evaluated with
\(\partial_\sigma\mathfrak{a} = \partial_\sigma\mathfrak{b} = -H\,\mathrm{sgn}(\sigma)\)
on the two sides (DERIVED, \texttt{um3\_fold.py}, J4 and J5, and
independently by smoothing the kink,
\(\partial_\sigma\mathfrak{a} = -H\tanh(\sigma/\varepsilon)\), and
integrating the full left-hand side of FE-1 across the layer: the result
converges linearly in \(\varepsilon\) to FE-30). With the bulk flux of
FE-29 they satisfy the conservation law FE-11 identically (DERIVED, J6).
Status: DERIVED.

\hypertarget{the-authors-vacuum}{%
\subsection{The author's vacuum}\label{the-authors-vacuum}}

\[
\begin{aligned}
&a_4 = Hx_4,\quad b_4 = \text{constant}\ (\mathsf{A} = 1,\ \mathsf{U} = \mathsf{A}_1 = \mathsf{U}_1 = 0):\qquad
 \sum_{p=1}^{3}\alpha_{(p)}\,\mathcal{E}_{(p)}{}^M{}_N + \Lambda\,\delta^M_N = 0\quad(\sigma \ne 0),\\
&\kappa\,\mathcal{S}^{(i)}{}_{(i)} = \kappa\,\mathcal{S}^{(r)}{}_{(r)} = -10H,\qquad \kappa\,\mathcal{S}^{(4)}{}_{(4)} = -6H,\\
&\kappa\rho_{\mathrm{bulk}} = -H^2\Big[\frac34\big((\mathsf{A} - 1 - \mathsf{U})^2 + (\mathsf{A} - 1)^2\big) + \dots\Big],\qquad
 \mathcal{G}(1, \mathsf{U}) = 2\mathsf{U}(3\mathsf{U} - 2)
\end{aligned} \tag{FE-31}
\] \textbf{Symbols in (FE-31).}

\begin{itemize}
\tightlist
\item
  \(a_4\): The author's metric function of the time \(x_4\) that scales
  ordinary 3-space: on the fold \(\mathrm{e}^{a_4}\) is the scale factor
  of cosmology and \(a_4'\) the Hubble rate. Unit: dimensionless.
  Status: AUTHOR.
\item
  \(H\): The author's constant (an inverse length) that fixes the period
  \(\pi/(6H)\) of \(x_0\), the Lovelock couplings and the cosmological
  constant \(\Lambda = -10H^2\); in UM-3 a free parameter. Unit: GeV
  (1/length). Status: AUTHOR.
\item
  \(x_4\): The time (\(x_4 \ge 0\) in the ad hoc vacuum); a prime on a
  function of \(x_4\) alone is \(\mathrm{d}/\mathrm{d}x_4\). Unit:
  length (= time, \(c = 1\)). Status: AUTHOR.
\item
  \(b_4\): The author's metric function of the time \(x_4\) in the
  extra-time components, \(b_4 > 0\), never identically 1 (the author's
  order of 2026-10-24); the scale factor of the extra times is
  \(\mathrm{e}^{-a_4}b_4\) times the warp. Unit: dimensionless. Status:
  AUTHOR.
\item
  \(\mathsf{A}\): The expansion rate of ordinary space in units of
  \(H\), \(\mathsf{A} = a_4'/H\). Unit: dimensionless. Value: 1 at the
  author's vacuum; \(H_0/H\) today (\(H_0\) the Hubble constant).
  Status: DERIVED.
\item
  \(\mathsf{U}\): The rate of change of \(b_4\) in units of \(H\),
  \(\mathsf{U} = b_4'/(Hb_4)\). Unit: dimensionless. Value:
  \(\mathsf{U} = \mathsf{A}\) with frozen extra times. Status: DERIVED.
\item
  \(\mathsf{A}_1\): The time derivative of \(\mathsf{A}\) in units of
  \(H\), \(\mathsf{A}_1 = \mathsf{A}'/H = a_4''/H^2\). Unit:
  dimensionless. Status: DERIVED.
\item
  \(\mathsf{U}_1\): The time derivative of \(\mathsf{U}\) in units of
  \(H\), \(\mathsf{U}_1 = \mathsf{U}'/H\). Unit: dimensionless. Status:
  DERIVED.
\item
  \(p\): The Lovelock order \(p = 1, 2, 3\) (Einstein, Gauss-Bonnet,
  cubic), as an index and as a number. Status: STANDARD.
\item
  \(\alpha_{(p)}\): The Lovelock coupling of order \(p\),
  \(\alpha_{(p)} = w_{(p)}H^{2-2p}\) (\(\alpha_{(1)} = 1\)). Unit:
  GeV\(^{2-2p}\). Status: AUTHOR.
\item
  \(\mathcal{E}_{(p)}{}^M{}_N\): The Lovelock tensor of order \(p\): the
  variation of \(\sqrt{\lvert g\rvert}\,\mathcal{L}_{(p)}\) with respect
  to the metric (\(\mathcal{E}_{(1)}\) the Einstein tensor,
  \(\mathcal{E}_{(2)}\) the Gauss-Bonnet tensor); with frame indices its
  frame components. Unit: 1/length\(^{2p}\). Status: STANDARD.
\item
  \(M\): A coordinate index of eight dimensions (also with a numbered
  copy \(M_1, M_2, \dots\)): it runs over \(x_0, x_1, \dots, x_7\); a
  repeated index is summed. Status: STANDARD.
\item
  \(N\): A coordinate index of eight dimensions (also
  \(N_1, N_2, \dots\)), running over \(x_0, \dots, x_7\); a repeated
  index is summed. Status: STANDARD.
\item
  \(\Lambda\): The eight-dimensional cosmological constant,
  \(\Lambda = -10H^2\) (the author writes the pure number
  \(\Lambda_{\mathrm{N}} = 10\) in
  \(0 = -\Lambda_{\mathrm{N}} + \sum_p H^{-2p}w_{(p)}\mathcal{E}_{(p)}\)).
  Unit: GeV\(^2\). Value: \(-10H^2\). Status: AUTHOR.
\item
  \(\delta\): The Kronecker delta: with two indices \(\delta^M_N = 1\)
  if \(M = N\) and 0 otherwise (likewise \(\delta_{X\mathrm{L}} = 1\)
  for \(X = \mathrm{L}\)); with upper and lower index strings the
  generalised Kronecker delta
  \(\delta^{M_1\cdots M_n}_{N_1\cdots N_n}\), the determinant of the
  \(n \times n\) matrix of the Kronecker deltas \(\delta^{M_a}_{N_b}\)
  (\(+1\) or \(-1\) if the upper indices are an even or odd permutation
  of the lower ones, 0 otherwise). Status: STANDARD.
\item
  \(\sigma\): The proper coordinate of the hidden direction:
  \(\mathrm{d}\sigma = \lvert\cot z\rvert\,\mathrm{d}x_0\),
  \(\sigma = \mathrm{sgn}(\pi/2 - z)\ln(\sin z)/(6H)\) on one period
  \(0 < z < \pi\); \(\sigma = 0\) is the fold \(z = \pi/2\),
  \(\sigma \to \mp\infty\) the tips \(z \to 0, \pi\); one period of
  \(x_0\) is the whole line \(-\infty < \sigma < \infty\). Unit: length.
  Status: DERIVED.
\item
  \(\kappa\): The eight-dimensional gravitational coupling,
  \(\kappa = 8\pi G_8\). Unit: GeV\(^{-6}\). Status: AUTHOR.
\item
  \(\mathcal{S}^a{}_b\): The thin source of the fold: its
  energy-momentum tensor per unit proper length across the fold
  (eight-dimensional normalisation, frame components along the fold);
  \(\kappa\mathcal{S}\) is the jump of the potential \(\mathfrak{F}\).
  Unit: GeV\(^7\). Status: DERIVED.
\item
  \(i\): A generation index (\(i = 1, 2, 3\)), which also labels the
  neutrino mass states \(\nu_1, \nu_2, \nu_3\); a repeated index is not
  summed unless a sum sign is written. Status: STANDARD.
\item
  \(r\): An extra-time direction of the fold: \(r \in \{5, 6, 7\}\)
  (\(x_5, x_6, x_7\)); not summed (the three components are equal).
  Status: STANDARD.
\item
  \(\rho_{\mathrm{bulk}}\): The energy density of the bulk sources,
  \(\rho_{\mathrm{bulk}} = -\mathcal{T}_{\mathrm{bulk}}{}^{(4)}{}_{(4)}\).
  Unit: GeV\(^8\). Status: DERIVED.
\item
  \(\mathcal{G}\): The factor of the \((\sigma, x_4)\) bulk equation of
  the author's metric,
  \(\mathcal{G} = 5\mathsf{A}^2(\mathsf{A} - \mathsf{U})^2 - 6\mathsf{A}^2 + 6\mathsf{A}\mathsf{U} + \mathsf{U}^2 + 1\).
  Unit: dimensionless. Status: DERIVED.
\end{itemize}

The author's solution (\texttt{Dirac\_Einstein-Lovelock\_4+4\_Nash.nb},
In{[}134{]}: \(a_4' = H\), the weights \((1, 1/32, 1/1152)\) and
\(\Lambda_{\mathrm{N}} = 10\)): ordinary space inflates with the rate
\(H\) forever and the extra times deflate
(\(\mathrm{e}^{-a_4}b_4 \propto \mathrm{e}^{-Hx_4}\)), ``eternal
inflation/deflation''. With the new metric (\(b_4^2\)) it is still an
exact vacuum solution of the bulk equations for every constant \(b_4\)
(DERIVED, \texttt{um3\_geometry.py}, G4). Three facts about it
(DERIVED):

\begin{itemize}
\tightlist
\item
  The fold must carry a thin source:
  \(\kappa\mathcal{S} = \mathrm{diag}(-10H, -6H, -10H)\) for (ordinary
  space, time, extra times), an energy density \(6H/\kappa\) with the
  pressure \(-10H/\kappa\) (\texttt{um3\_fold.py}, J5).
\item
  The vacuum is degenerate: near it the energy density of any bulk
  source is
  \(\kappa\rho_{\mathrm{bulk}} = -H^2[\frac34((\mathsf{A} - 1 - \mathsf{U})^2 + (\mathsf{A} - 1)^2) + \dots]\)
  (the dots are of third order in \(\mathsf{A} - 1\) and
  \(\mathsf{U}\)): the linear terms vanish, so a small positive energy
  density of the bulk cannot be carried by a small change of \(a_4\) and
  \(b_4\) (\texttt{um3\_geometry.py}, G4). The transverse traceless
  gravitons of the bulk have no time-kinetic term at this point
  (\texttt{um3\_bulk.py}, B5).
\item
  \(b_4\) is a function of \(x_4\), never identically 1 (the author's
  order): a constant \(b_4 \ne 1\) is equivalent to \(b_4 = 1\) by a
  rescaling of \(x_5, x_6, x_7\) unless their coordinate period is
  fixed, and \(\mathsf{U} \ne 0\) is not a vacuum solution near
  \(\mathsf{A} = 1\): it needs a bulk source with a flux along
  \(\sigma\),
  \(\kappa\mathcal{T}^{(0)}{}_{(4)} = \mathrm{sgn}(\sigma)\frac{3}{16}\mathsf{U}\,\mathcal{G}(1, \mathsf{U})H^2 = \mathrm{sgn}(\sigma)\frac38\mathsf{U}^2(3\mathsf{U} - 2)H^2\),
  and a non-positive energy density. At \(\mathsf{A} = 1\) the flux
  vanishes for \(\mathsf{U} = 2/3\)
  (\(b_4 \propto \mathrm{e}^{2Hx_4/3}\)), which needs the constant
  source \(\kappa\mathcal{T} = (-14/27, -2/9, 2/27, -22/27)H^2\) for
  (\(\sigma\), ordinary space, time, extra times), with
  \(\kappa\rho_{\mathrm{bulk}} = -2H^2/27\) (\texttt{um3\_geometry.py},
  G4).
\end{itemize}

Status: AUTHOR (the solution), DERIVED (the three facts).

\hypertarget{the-static-block-and-the-frozen-extra-times}{%
\subsection{The static block and the frozen extra
times}\label{the-static-block-and-the-frozen-extra-times}}

\[
\begin{aligned}
&\mathsf{A} = \mathsf{U} = \mathsf{A}_1 = \mathsf{U}_1 = 0:\qquad
 \kappa\,\mathcal{T}_{\mathrm{bulk}}{}^a{}_b = H^2\,\mathrm{diag}\Big(-\frac{5}{16}, -\frac{5}{16}, -\frac{5}{16}, -\frac{5}{16},
 \frac{1}{16}, -\frac{5}{16}, -\frac{5}{16}, -\frac{5}{16}\Big),\\
&\kappa\,\mathcal{S}^{(i)}{}_{(i)} = \kappa\,\mathcal{S}^{(r)}{}_{(r)} = -\frac{61}{8}H,\qquad \kappa\,\mathcal{S}^{(4)}{}_{(4)} = -\frac{31}{4}H
\end{aligned} \tag{FE-32}
\] \textbf{Symbols in (FE-32).}

\begin{itemize}
\tightlist
\item
  \(\mathsf{A}\): The expansion rate of ordinary space in units of
  \(H\), \(\mathsf{A} = a_4'/H\). Unit: dimensionless. Value: 1 at the
  author's vacuum; \(H_0/H\) today (\(H_0\) the Hubble constant).
  Status: DERIVED.
\item
  \(\mathsf{U}\): The rate of change of \(b_4\) in units of \(H\),
  \(\mathsf{U} = b_4'/(Hb_4)\). Unit: dimensionless. Value:
  \(\mathsf{U} = \mathsf{A}\) with frozen extra times. Status: DERIVED.
\item
  \(\mathsf{A}_1\): The time derivative of \(\mathsf{A}\) in units of
  \(H\), \(\mathsf{A}_1 = \mathsf{A}'/H = a_4''/H^2\). Unit:
  dimensionless. Status: DERIVED.
\item
  \(\mathsf{U}_1\): The time derivative of \(\mathsf{U}\) in units of
  \(H\), \(\mathsf{U}_1 = \mathsf{U}'/H\). Unit: dimensionless. Status:
  DERIVED.
\item
  \(\kappa\): The eight-dimensional gravitational coupling,
  \(\kappa = 8\pi G_8\). Unit: GeV\(^{-6}\). Status: AUTHOR.
\item
  \(\mathcal{T}_{\mathrm{bulk}}\): The energy-momentum tensor of all
  sources of the bulk,
  \(\mathcal{T}_\Psi + \mathcal{T}_{\mathrm{stab}}\);
  \(\mathcal{T}_{\mathrm{bulk}}{}^{(0)}{}_b\) is its flux along
  \(\sigma\) (frame components). Unit: GeV\(^8\). Status: STANDARD.
\item
  \(a\): A frame (orthonormal) index running over \(0, 1, \dots, 7\)
  (frame direction \((0)\) along \(\sigma\), \((1), (2), (3)\) along
  ordinary space, \((4)\) along the time, \((5), (6), (7)\) along the
  extra times); raised and lowered with \(\eta_{ab}\); on the fold
  (\(a, b \in \{1, \dots, 7\}\)) the frame of the fold; a repeated index
  is summed. Status: STANDARD.
\item
  \(b\): A frame index (as \(a\)). Status: STANDARD.
\item
  \(H\): The author's constant (an inverse length) that fixes the period
  \(\pi/(6H)\) of \(x_0\), the Lovelock couplings and the cosmological
  constant \(\Lambda = -10H^2\); in UM-3 a free parameter. Unit: GeV
  (1/length). Status: AUTHOR.
\item
  \(\mathrm{diag}\): A diagonal matrix with the listed entries on its
  diagonal (in the order of the coordinates \(x_0, x_1, \dots, x_7\) for
  a tensor of eight dimensions, \(x_1, \dots, x_7\) for a tensor of the
  fold, \(x_1, \dots, x_4\) for a four-dimensional tensor). Status:
  STANDARD.
\item
  \(\mathcal{S}^a{}_b\): The thin source of the fold: its
  energy-momentum tensor per unit proper length across the fold
  (eight-dimensional normalisation, frame components along the fold);
  \(\kappa\mathcal{S}\) is the jump of the potential \(\mathfrak{F}\).
  Unit: GeV\(^7\). Status: DERIVED.
\item
  \(i\): A generation index (\(i = 1, 2, 3\)), which also labels the
  neutrino mass states \(\nu_1, \nu_2, \nu_3\); a repeated index is not
  summed unless a sum sign is written. Status: STANDARD.
\item
  \(r\): An extra-time direction of the fold: \(r \in \{5, 6, 7\}\)
  (\(x_5, x_6, x_7\)); not summed (the three components are equal).
  Status: STANDARD.
\end{itemize}

The static block: \(a_4\) and \(b_4\) constant. Its bulk is (each half
of a period) a seven-dimensional space of constant negative curvature
\(-H^2\) of the directions \(\sigma, x_1, x_2, x_3, x_5, x_6, x_7\)
times the unwarped time (FE-28 with \(\mathsf{A} = \mathsf{U} = 0\)). It
is not a vacuum solution: it requires the bulk source of the first line,
which is the stabiliser IN-7 (\(\kappa\rho = -H^2/16\),
\(\kappa p = -5H^2/16\)), and the thin source of the second line at the
fold, which is the fold fluid at rest
(\(\kappa\rho_{\mathrm{ff}}/V_* = 31H/4\),
\(\kappa p_{\mathrm{ff}}/V_* = \kappa q_{\mathrm{ff}}/V_* = -61H/8\),
\(\rho + p = H/(8\kappa)\) per \(V_*\), positive). The late universe,
whose expansion rate is tiny in units of \(H\), is the static block plus
small corrections (FE-33, FE-59). Status: DERIVED.

\[
\begin{aligned}
&\mathsf{U} = \mathsf{A},\quad \mathsf{U}_1 = \mathsf{A}_1\ \ \big(\mathrm{e}^{-a_4}b_4 = \text{constant},\ \mathcal{Z} = 1\big):\\
&\frac{\kappa\rho_{\mathrm{bulk}}}{H^2} = \frac{3\mathsf{A}^2 - 1}{16},\qquad
 \frac{\kappa\,\mathcal{T}_{\mathrm{bulk}}{}^{(0)}{}_{(0)}}{H^2} = -\frac{3\mathsf{A}^4 + 3\mathsf{A}^2\mathsf{A}_1 + 18\mathsf{A}^2 + 3\mathsf{A}_1 + 5}{16},\\
&\frac{\kappa\,\mathcal{T}_{\mathrm{bulk}}{}^{(i)}{}_{(i)}}{H^2} = -\frac{9\mathsf{A}^2 + 2\mathsf{A}_1 + 5}{16},\qquad
 \frac{\kappa\,\mathcal{T}_{\mathrm{bulk}}{}^{(r)}{}_{(r)}}{H^2} = -\frac{3\mathsf{A}^4 + 3\mathsf{A}^2\mathsf{A}_1 + 24\mathsf{A}^2 + 3\mathsf{A}_1 + 5}{16},\\
&\frac{\kappa\,\mathcal{T}_{\mathrm{bulk}}{}^{(0)}{}_{(4)}}{H^2} = \mathrm{sgn}(\sigma)\,\frac{3\mathsf{A}(1 + \mathsf{A}^2)}{16},\qquad
 \frac{\kappa\,\mathcal{S}^{(i)}{}_{(i)}}{H} = -\frac{57\mathsf{A}^2 + 40\mathsf{A}_1 + 183}{24},\\
&\frac{\kappa\,\mathcal{S}^{(4)}{}_{(4)}}{H} = -\frac{10\mathsf{A}^2 + 31}{4},\qquad
 \frac{\kappa\,\mathcal{S}^{(r)}{}_{(r)}}{H} = -\frac{2\mathsf{A}^4 + 2\mathsf{A}^2\mathsf{A}_1 + 37\mathsf{A}^2 + 20\mathsf{A}_1 + 61}{8}
\end{aligned} \tag{FE-33}
\] \textbf{Symbols in (FE-33).}

\begin{itemize}
\tightlist
\item
  \(\mathsf{U}\): The rate of change of \(b_4\) in units of \(H\),
  \(\mathsf{U} = b_4'/(Hb_4)\). Unit: dimensionless. Value:
  \(\mathsf{U} = \mathsf{A}\) with frozen extra times. Status: DERIVED.
\item
  \(\mathsf{A}\): The expansion rate of ordinary space in units of
  \(H\), \(\mathsf{A} = a_4'/H\). Unit: dimensionless. Value: 1 at the
  author's vacuum; \(H_0/H\) today (\(H_0\) the Hubble constant).
  Status: DERIVED.
\item
  \(\mathsf{U}_1\): The time derivative of \(\mathsf{U}\) in units of
  \(H\), \(\mathsf{U}_1 = \mathsf{U}'/H\). Unit: dimensionless. Status:
  DERIVED.
\item
  \(\mathsf{A}_1\): The time derivative of \(\mathsf{A}\) in units of
  \(H\), \(\mathsf{A}_1 = \mathsf{A}'/H = a_4''/H^2\). Unit:
  dimensionless. Status: DERIVED.
\item
  \(\mathrm{e}\): Euler's number, the base of the natural logarithm;
  \(\mathrm{e}^{X}\) is the exponential of \(X\). Unit: dimensionless.
  Value: 2.718281828. Status: STANDARD.
\item
  \(a_4\): The author's metric function of the time \(x_4\) that scales
  ordinary 3-space: on the fold \(\mathrm{e}^{a_4}\) is the scale factor
  of cosmology and \(a_4'\) the Hubble rate. Unit: dimensionless.
  Status: AUTHOR.
\item
  \(b_4\): The author's metric function of the time \(x_4\) in the
  extra-time components, \(b_4 > 0\), never identically 1 (the author's
  order of 2026-10-24); the scale factor of the extra times is
  \(\mathrm{e}^{-a_4}b_4\) times the warp. Unit: dimensionless. Status:
  AUTHOR.
\item
  \(\mathcal{Z}\): The fold factor,
  \(\mathcal{Z}(x_4) = V_{\mathrm{t}}(\mathrm{e}^{-a_4}b_4)^3/V_*\): the
  factor by which the size of the extra times multiplies every term of
  the fold; \(\mathcal{Z} = 1\) when the extra times are frozen at their
  late value. Unit: dimensionless. Value: 1 in the late universe.
  Status: DERIVED.
\item
  \(\kappa\): The eight-dimensional gravitational coupling,
  \(\kappa = 8\pi G_8\). Unit: GeV\(^{-6}\). Status: AUTHOR.
\item
  \(\rho_{\mathrm{bulk}}\): The energy density of the bulk sources,
  \(\rho_{\mathrm{bulk}} = -\mathcal{T}_{\mathrm{bulk}}{}^{(4)}{}_{(4)}\).
  Unit: GeV\(^8\). Status: DERIVED.
\item
  \(H\): The author's constant (an inverse length) that fixes the period
  \(\pi/(6H)\) of \(x_0\), the Lovelock couplings and the cosmological
  constant \(\Lambda = -10H^2\); in UM-3 a free parameter. Unit: GeV
  (1/length). Status: AUTHOR.
\item
  \(\mathcal{T}_{\mathrm{bulk}}\): The energy-momentum tensor of all
  sources of the bulk,
  \(\mathcal{T}_\Psi + \mathcal{T}_{\mathrm{stab}}\);
  \(\mathcal{T}_{\mathrm{bulk}}{}^{(0)}{}_b\) is its flux along
  \(\sigma\) (frame components). Unit: GeV\(^8\). Status: STANDARD.
\item
  \(i\): A generation index (\(i = 1, 2, 3\)), which also labels the
  neutrino mass states \(\nu_1, \nu_2, \nu_3\); a repeated index is not
  summed unless a sum sign is written. Status: STANDARD.
\item
  \(r\): An extra-time direction of the fold: \(r \in \{5, 6, 7\}\)
  (\(x_5, x_6, x_7\)); not summed (the three components are equal).
  Status: STANDARD.
\item
  \(\mathrm{sgn}\): The sign function: \(\mathrm{sgn}(y) = +1\) for
  \(y > 0\), \(-1\) for \(y < 0\). Status: STANDARD.
\item
  \(\sigma\): The proper coordinate of the hidden direction:
  \(\mathrm{d}\sigma = \lvert\cot z\rvert\,\mathrm{d}x_0\),
  \(\sigma = \mathrm{sgn}(\pi/2 - z)\ln(\sin z)/(6H)\) on one period
  \(0 < z < \pi\); \(\sigma = 0\) is the fold \(z = \pi/2\),
  \(\sigma \to \mp\infty\) the tips \(z \to 0, \pi\); one period of
  \(x_0\) is the whole line \(-\infty < \sigma < \infty\). Unit: length.
  Status: DERIVED.
\item
  \(\mathcal{S}^a{}_b\): The thin source of the fold: its
  energy-momentum tensor per unit proper length across the fold
  (eight-dimensional normalisation, frame components along the fold);
  \(\kappa\mathcal{S}\) is the jump of the potential \(\mathfrak{F}\).
  Unit: GeV\(^7\). Status: DERIVED.
\end{itemize}

The frozen extra times: the physical size \(\mathrm{e}^{-a_4}b_4\) of
the extra times is constant while ordinary space expands. Then the fold
factor \(\mathcal{Z}\) is constant (IN-8) and every coupling of the
Standard Model is constant in time, as observation requires (the
fine-structure constant has not changed by more than parts in \(10^{6}\)
since the CMB); this is the condition of the late universe. It gives
\(b_4 = c_b\,\mathrm{e}^{a_4}\) with a constant \(c_b\): \(b_4\) grows
with the scale factor of ordinary space, so it is a non-trivial function
of \(x_4\), never identically 1, as the author requires. The equations
of FE-29 and FE-30 on this branch are written out (DERIVED,
\texttt{um3\_geometry.py}, G4, \texttt{um3\_fold.py}, J7). For
\(\mathsf{A} \to 0\) they reduce to the static block FE-32. The extra
energy of the fold above the static value is \((5/2)a_4'^2/(\kappa H)\)
per \(V_*\): on the separable metric the fold energy fixes the expansion
rate; a separately conserved fold fluid fits the separable metric only
for the equation of state \(p = -\rho/40\) (DERIVED, J7), so ordinary
matter on the fold requires the next terms of the power series (FE-35,
FE-59). Status: DERIVED.

\hypertarget{the-general-warped-metric-and-the-power-series-of-the-metric}{%
\subsection{The general warped metric and the power series of the
metric}\label{the-general-warped-metric-and-the-power-series-of-the-metric}}

\[
\begin{aligned}
&g_{MN}\,\mathrm{d}X^M\mathrm{d}X^N = \mathrm{d}\sigma^2 + \mathrm{e}^{2\mathfrak{a}}\big(\mathrm{d}x_1^2 + \mathrm{d}x_2^2 + \mathrm{d}x_3^2\big)
 - \mathrm{e}^{2\mathfrak{n}}\,\mathrm{d}x_4^2 - \mathrm{e}^{2\mathfrak{b}}\big(\mathrm{d}x_5^2 + \mathrm{d}x_6^2 + \mathrm{d}x_7^2\big),\\
&R^{(0)(i)}{}_{(0)(i)} = -\partial_\sigma^2\mathfrak{a} - (\partial_\sigma\mathfrak{a})^2,\qquad
 R^{(0)(4)}{}_{(0)(4)} = -\partial_\sigma^2\mathfrak{n} - (\partial_\sigma\mathfrak{n})^2,\qquad
 R^{(i)(j)}{}_{(i)(j)} = -(\partial_\sigma\mathfrak{a})^2 + \mathrm{e}^{-2\mathfrak{n}}(\partial_{x_4}\mathfrak{a})^2,\\
&R^{(i)(r)}{}_{(i)(r)} = -\partial_\sigma\mathfrak{a}\,\partial_\sigma\mathfrak{b} + \mathrm{e}^{-2\mathfrak{n}}\,\partial_{x_4}\mathfrak{a}\,\partial_{x_4}\mathfrak{b},\qquad
 R^{(i)(4)}{}_{(i)(4)} = -\partial_\sigma\mathfrak{a}\,\partial_\sigma\mathfrak{n} + \mathrm{e}^{-2\mathfrak{n}}\big(\partial_{x_4}^2\mathfrak{a}
 + (\partial_{x_4}\mathfrak{a})^2 - \partial_{x_4}\mathfrak{a}\,\partial_{x_4}\mathfrak{n}\big),\\
&R^{(0)(i)}{}_{(4)(i)} = -\mathrm{e}^{-\mathfrak{n}}\big(\partial_\sigma\partial_{x_4}\mathfrak{a} + \partial_\sigma\mathfrak{a}\,\partial_{x_4}\mathfrak{a}
 - \partial_{x_4}\mathfrak{a}\,\partial_\sigma\mathfrak{n}\big)
\end{aligned} \tag{FE-34}
\] \textbf{Symbols in (FE-34).}

\begin{itemize}
\tightlist
\item
  \(g_{MN}\): The metric tensor of the eight-dimensional spacetime
  (signature \(+,+,+,+,-,-,-,-\) in the order \(x_0, \dots, x_7\)):
  \(g_{MN}\,\mathrm{d}X^M\mathrm{d}X^N\) is the squared interval,
  \(g^{MN}\) its inverse, \(\lvert g\rvert\) the absolute value of its
  determinant (\(\sqrt{\lvert g\rvert}\,\mathrm{d}^8X\) is the proper
  volume element). Unit: dimensionless. Status: AUTHOR.
\item
  \(M\): A coordinate index of eight dimensions (also with a numbered
  copy \(M_1, M_2, \dots\)): it runs over \(x_0, x_1, \dots, x_7\); a
  repeated index is summed. Status: STANDARD.
\item
  \(N\): A coordinate index of eight dimensions (also
  \(N_1, N_2, \dots\)), running over \(x_0, \dots, x_7\); a repeated
  index is summed. Status: STANDARD.
\item
  \(\mathrm{d}\): The differential: \(\mathrm{d}X\) is an infinitesimal
  change of \(X\), \(\mathrm{d}X/\mathrm{d}Y\) the derivative of \(X\)
  with respect to \(Y\), \(\mathrm{d}^nX\) the volume element of \(n\)
  variables. Status: STANDARD.
\item
  \(X = \{x_0, \dots, x_7\}\): The author's coordinates of the
  eight-dimensional spacetime, \(X = \{x_0, x_1, \dots, x_7\}\) with the
  hidden direction first (Dirac\_Einstein-Lovelock\_4+4\_Nash.nb,
  section DEFCHART); \(X^M\) is the coordinate with the index \(M\) and
  \(\mathrm{d}^8X = \mathrm{d}x_0\,\mathrm{d}x_1\cdots\mathrm{d}x_7\)
  the coordinate volume element. Unit: length (each coordinate). Status:
  AUTHOR.
\item
  \(\sigma\): The proper coordinate of the hidden direction:
  \(\mathrm{d}\sigma = \lvert\cot z\rvert\,\mathrm{d}x_0\),
  \(\sigma = \mathrm{sgn}(\pi/2 - z)\ln(\sin z)/(6H)\) on one period
  \(0 < z < \pi\); \(\sigma = 0\) is the fold \(z = \pi/2\),
  \(\sigma \to \mp\infty\) the tips \(z \to 0, \pi\); one period of
  \(x_0\) is the whole line \(-\infty < \sigma < \infty\). Unit: length.
  Status: DERIVED.
\item
  \(\mathrm{e}\): Euler's number, the base of the natural logarithm;
  \(\mathrm{e}^{X}\) is the exponential of \(X\). Unit: dimensionless.
  Value: 2.718281828. Status: STANDARD.
\item
  \(\mathfrak{a}\): The logarithm of the scale factor of ordinary space
  in the general warped metric FE-34, a function of \(\sigma\) and
  \(x_4\); for the author's metric
  \(\mathfrak{a} = a_4 - H\lvert\sigma\rvert\); \(\mathfrak{a}_n\) its
  term of order \(n\) in the power series FE-35. Unit: dimensionless.
  Status: DERIVED.
\item
  \(x_1\): The first coordinate of ordinary 3-space. Unit: length.
  Status: AUTHOR.
\item
  \(x_2\): The second coordinate of ordinary 3-space. Unit: length.
  Status: AUTHOR.
\item
  \(x_3\): The third coordinate of ordinary 3-space. Unit: length.
  Status: AUTHOR.
\item
  \(\mathfrak{n}\): The logarithm of the lapse in the general warped
  metric FE-34 (\(g_{x_4x_4} = -\mathrm{e}^{2\mathfrak{n}}\)); for the
  author's metric \(\mathfrak{n} = 0\); \(\mathfrak{n}_n\) its term of
  order \(n\). Unit: dimensionless. Status: DERIVED.
\item
  \(x_4\): The time (\(x_4 \ge 0\) in the ad hoc vacuum); a prime on a
  function of \(x_4\) alone is \(\mathrm{d}/\mathrm{d}x_4\). Unit:
  length (= time, \(c = 1\)). Status: AUTHOR.
\item
  \(\mathfrak{b}\): The logarithm of the scale factor of the extra times
  in the general warped metric FE-34; for the author's metric
  \(\mathfrak{b} = \ln b_4 - a_4 - H\lvert\sigma\rvert\);
  \(\mathfrak{b}_n\) its term of order \(n\). Unit: dimensionless.
  Status: DERIVED.
\item
  \(x_5\): The first of the three extra time-like coordinates. Unit:
  length. Status: AUTHOR.
\item
  \(x_6\): The second of the three extra time-like coordinates. Unit:
  length. Status: AUTHOR.
\item
  \(x_7\): The third of the three extra time-like coordinates. Unit:
  length. Status: AUTHOR.
\item
  \(R^{MN}{}_{PQ},\ R_{MN},\ R\): The curvature of the metric
  \(g_{MN}\): with four indices the Riemann tensor (convention
  \([\nabla_P, \nabla_Q]V^M = R^M{}_{NPQ}V^N\)), with two indices the
  Ricci tensor \(R_{MN} = R^P{}_{MPN}\), without indices the scalar
  curvature \(R = g^{MN}R_{MN}\); with frame indices
  \(R^{(a)(b)}{}_{(c)(d)}\) its frame components,
  \(R^{(a)(b)}{}_{(a)(b)}\) (no sum) the sectional curvature of the
  frame plane \((a, b)\). Unit: 1/length\(^2\). Status: STANDARD.
\item
  \(i\): A generation index (\(i = 1, 2, 3\)), which also labels the
  neutrino mass states \(\nu_1, \nu_2, \nu_3\); a repeated index is not
  summed unless a sum sign is written. Status: STANDARD.
\item
  \(j\): A second generation or mass-state index (as \(i\)). Status:
  STANDARD.
\item
  \(r\): An extra-time direction of the fold: \(r \in \{5, 6, 7\}\)
  (\(x_5, x_6, x_7\)); not summed (the three components are equal).
  Status: STANDARD.
\item
  \(\partial\): A partial derivative: \(\partial_M\) with respect to the
  coordinate \(X^M\), \(\partial_\mu\) with respect to \(x^\mu\),
  \(\partial_\sigma\) with respect to \(\sigma\),
  \(\partial X/\partial Y\) the derivative of \(X\) with respect to
  \(Y\) at fixed other variables. Status: STANDARD.
\item
  \(\sigma\): The label \(\sigma\) in \(\partial_\sigma\): the
  derivative with respect to the proper hidden coordinate \(\sigma\).
  Status: STANDARD.
\end{itemize}

The general metric with the symmetries of the author's metric
(homogeneous and isotropic in ordinary space and in the extra times,
independent of the extra times) but with three arbitrary functions
\(\mathfrak{a}\), \(\mathfrak{b}\), \(\mathfrak{n}\) of \(\sigma\) and
\(x_4\), and its frame curvature: the components with the extra times
follow from those with ordinary space by replacing \(\mathfrak{a}\) by
\(\mathfrak{b}\) and \((i)\) by \((r)\) (\(R^{(0)(r)}{}_{(0)(r)}\),
\(R^{(r)(s)}{}_{(r)(s)}\), \(R^{(r)(4)}{}_{(r)(4)}\),
\(R^{(0)(r)}{}_{(4)(r)}\)); all others vanish (DERIVED,
\texttt{um3\_fold.py}, J2). Its Lovelock tensor is FE-2 with these
components; it has only diagonal components and the \((\sigma, x_4)\)
pair. The author's metric is
\(\mathfrak{a} = a_4 - H\lvert\sigma\rvert\),
\(\mathfrak{b} = \ln b_4 - a_4 - H\lvert\sigma\rvert\),
\(\mathfrak{n} = 0\), and the general equations reduce to FE-29 on both
sides of the fold (DERIVED, J2). Status: DERIVED.

\[
\begin{aligned}
&\mathfrak{a} = a_4 - H\lvert\sigma\rvert + \sum_{n\ge1}\epsilon_{\mathrm{r}}^{\,n}\,\mathfrak{a}_n(\sigma, x_4),\qquad
 \mathfrak{b} = \ln b_4 - a_4 - H\lvert\sigma\rvert + \sum_{n\ge1}\epsilon_{\mathrm{r}}^{\,n}\,\mathfrak{b}_n(\sigma, x_4),\\
&\mathfrak{n} = \sum_{n\ge1}\epsilon_{\mathrm{r}}^{\,n}\,\mathfrak{n}_n(\sigma, x_4),\qquad
 \hat h_{\mu\nu} \to \hat h_{\mu\nu} + \epsilon_{\mathrm{c}}\,\hat h^{(1)}_{\mu\nu} + \dots
\end{aligned} \tag{FE-35}
\] \textbf{Symbols in (FE-35).}

\begin{itemize}
\tightlist
\item
  \(\mathfrak{a}\): The logarithm of the scale factor of ordinary space
  in the general warped metric FE-34, a function of \(\sigma\) and
  \(x_4\); for the author's metric
  \(\mathfrak{a} = a_4 - H\lvert\sigma\rvert\); \(\mathfrak{a}_n\) its
  term of order \(n\) in the power series FE-35. Unit: dimensionless.
  Status: DERIVED.
\item
  \(a_4\): The author's metric function of the time \(x_4\) that scales
  ordinary 3-space: on the fold \(\mathrm{e}^{a_4}\) is the scale factor
  of cosmology and \(a_4'\) the Hubble rate. Unit: dimensionless.
  Status: AUTHOR.
\item
  \(H\): The author's constant (an inverse length) that fixes the period
  \(\pi/(6H)\) of \(x_0\), the Lovelock couplings and the cosmological
  constant \(\Lambda = -10H^2\); in UM-3 a free parameter. Unit: GeV
  (1/length). Status: AUTHOR.
\item
  \(\sigma\): The proper coordinate of the hidden direction:
  \(\mathrm{d}\sigma = \lvert\cot z\rvert\,\mathrm{d}x_0\),
  \(\sigma = \mathrm{sgn}(\pi/2 - z)\ln(\sin z)/(6H)\) on one period
  \(0 < z < \pi\); \(\sigma = 0\) is the fold \(z = \pi/2\),
  \(\sigma \to \mp\infty\) the tips \(z \to 0, \pi\); one period of
  \(x_0\) is the whole line \(-\infty < \sigma < \infty\). Unit: length.
  Status: DERIVED.
\item
  \(n\): The order \(n = 1, 2, \dots\) of a term of the power series of
  the metric (also its exponent of \(\epsilon_{\mathrm{r}}\));
  \(\mathfrak{a}_n\), \(\mathfrak{b}_n\), \(\mathfrak{n}_n\) are the
  terms of order \(n\). Status: STANDARD.
\item
  \(\epsilon_{\mathrm{r}}\): The bookkeeping parameter of the response
  of the bulk to the matter of the fold: the ratio of the bulk to the
  fold gravity,
  \(\epsilon_{\mathrm{r}} \sim r = \kappa_4V_*/(\kappa H)\). Unit:
  dimensionless. Value: \(\lesssim 10^{-15}\). Status: DERIVED.
\item
  \(x_4\): The time (\(x_4 \ge 0\) in the ad hoc vacuum); a prime on a
  function of \(x_4\) alone is \(\mathrm{d}/\mathrm{d}x_4\). Unit:
  length (= time, \(c = 1\)). Status: AUTHOR.
\item
  \(\mathfrak{b}\): The logarithm of the scale factor of the extra times
  in the general warped metric FE-34; for the author's metric
  \(\mathfrak{b} = \ln b_4 - a_4 - H\lvert\sigma\rvert\);
  \(\mathfrak{b}_n\) its term of order \(n\). Unit: dimensionless.
  Status: DERIVED.
\item
  \(b_4\): The author's metric function of the time \(x_4\) in the
  extra-time components, \(b_4 > 0\), never identically 1 (the author's
  order of 2026-10-24); the scale factor of the extra times is
  \(\mathrm{e}^{-a_4}b_4\) times the warp. Unit: dimensionless. Status:
  AUTHOR.
\item
  \(\mathfrak{n}\): The logarithm of the lapse in the general warped
  metric FE-34 (\(g_{x_4x_4} = -\mathrm{e}^{2\mathfrak{n}}\)); for the
  author's metric \(\mathfrak{n} = 0\); \(\mathfrak{n}_n\) its term of
  order \(n\). Unit: dimensionless. Status: DERIVED.
\item
  \(\hat h_{\mu\nu}\): The four-dimensional metric of the fold,
  \(\hat h_{\mu\nu} = \mathrm{diag}(\mathrm{e}^{2a_4}, \mathrm{e}^{2a_4}, \mathrm{e}^{2a_4}, -1)\)
  in the coordinates \(x_1, \dots, x_4\) (the induced metric at
  \(\sigma = 0\) without the extra times); \(\hat h^{\mu\nu}\) its
  inverse, \(\hat h\) its determinant,
  \(\sqrt{-\hat h} = \mathrm{e}^{3a_4}\); \(\hat h^{(1)}_{\mu\nu}\) its
  first-order cosmological perturbation. Unit: dimensionless. Status:
  DERIVED.
\item
  \(\mu\): A four-dimensional coordinate index running over the
  observable coordinates \(x_1, x_2, x_3, x_4\) of the fold (signature
  \(+,+,+,-\)); a repeated index is summed. Status: STANDARD.
\item
  \(\nu\): A four-dimensional coordinate index (as \(\mu\)). Status:
  STANDARD.
\item
  \(\epsilon_{\mathrm{c}}\): The bookkeeping parameter of the
  cosmological perturbations of the metric of the fold (the size of the
  CMB anisotropies). Unit: dimensionless. Value: \(\sim 10^{-5}\).
  Status: STANDARD.
\end{itemize}

The metric as a power series of terms of decreasing magnitude (the
author's request): the order-0 term is the author's metric; the terms of
order \(n\) in \(\epsilon_{\mathrm{r}}\) are the response of the bulk to
the matter of the fold and of dirac16complex, which depends on
\(\sigma\); \(\epsilon_{\mathrm{c}} \sim 10^{-5}\) (the size of the CMB
anisotropies) counts the cosmological perturbations of the metric of the
fold. The bookkeeping parameter \(\epsilon_{\mathrm{r}}\) is the ratio
of the bulk to the fold gravity, \(r = \kappa_4V_*/(\kappa H)\) (FE-59).
Each order obeys the linearisation of FE-1 and FE-4 around the previous
orders, with the sources of that order. Status: DERIVED (the scheme).

\end{document}
